% Options for packages loaded elsewhere
% Options for packages loaded elsewhere
\PassOptionsToPackage{unicode}{hyperref}
\PassOptionsToPackage{hyphens}{url}
%
\documentclass[
  11pt,
  letterpaper,
]{scrbook}
\usepackage{xcolor}
\usepackage[paperwidth=6in,paperheight=9in,inner=0.85in,outer=0.95in,top=0.9in,bottom=1.1in,includehead,heightrounded]{geometry}
\usepackage{amsmath,amssymb}
\setcounter{secnumdepth}{2}
\usepackage{iftex}
\ifPDFTeX
  \usepackage[T1]{fontenc}
  \usepackage[utf8]{inputenc}
  \usepackage{textcomp} % provide euro and other symbols
\else % if luatex or xetex
  \usepackage{unicode-math} % this also loads fontspec
  \defaultfontfeatures{Scale=MatchLowercase}
  \defaultfontfeatures[\rmfamily]{Ligatures=TeX,Scale=1}
\fi
\usepackage{lmodern}
\ifPDFTeX\else
  % xetex/luatex font selection
  \setmainfont[BoldFont=EB Garamond SemiBold,ItalicFont=EB Garamond
Italic,BoldItalicFont=EB Garamond SemiBold Italic]{EB Garamond Math}
  \setmathfont[]{Garamond-Math}
\fi
% Use upquote if available, for straight quotes in verbatim environments
\IfFileExists{upquote.sty}{\usepackage{upquote}}{}
\IfFileExists{microtype.sty}{% use microtype if available
  \usepackage[]{microtype}
  \UseMicrotypeSet[protrusion]{basicmath} % disable protrusion for tt fonts
}{}
\makeatletter
\@ifundefined{KOMAClassName}{% if non-KOMA class
  \IfFileExists{parskip.sty}{%
    \usepackage{parskip}
  }{% else
    \setlength{\parindent}{0pt}
    \setlength{\parskip}{6pt plus 2pt minus 1pt}}
}{% if KOMA class
  \KOMAoptions{parskip=half}}
\makeatother
% Make \paragraph and \subparagraph free-standing
\makeatletter
\ifx\paragraph\undefined\else
  \let\oldparagraph\paragraph
  \renewcommand{\paragraph}{
    \@ifstar
      \xxxParagraphStar
      \xxxParagraphNoStar
  }
  \newcommand{\xxxParagraphStar}[1]{\oldparagraph*{#1}\mbox{}}
  \newcommand{\xxxParagraphNoStar}[1]{\oldparagraph{#1}\mbox{}}
\fi
\ifx\subparagraph\undefined\else
  \let\oldsubparagraph\subparagraph
  \renewcommand{\subparagraph}{
    \@ifstar
      \xxxSubParagraphStar
      \xxxSubParagraphNoStar
  }
  \newcommand{\xxxSubParagraphStar}[1]{\oldsubparagraph*{#1}\mbox{}}
  \newcommand{\xxxSubParagraphNoStar}[1]{\oldsubparagraph{#1}\mbox{}}
\fi
\makeatother


\usepackage{longtable,booktabs,array}
\newcounter{none} % for unnumbered tables
\usepackage{multirow}
\usepackage{calc} % for calculating minipage widths
% Correct order of tables after \paragraph or \subparagraph
\usepackage{etoolbox}
\makeatletter
\patchcmd\longtable{\par}{\if@noskipsec\mbox{}\fi\par}{}{}
\makeatother
% Allow footnotes in longtable head/foot
\IfFileExists{footnotehyper.sty}{\usepackage{footnotehyper}}{\usepackage{footnote}}
\makesavenoteenv{longtable}
\usepackage{graphicx}
\makeatletter
\newsavebox\pandoc@box
\newcommand*\pandocbounded[1]{% scales image to fit in text height/width
  \sbox\pandoc@box{#1}%
  \Gscale@div\@tempa{\textheight}{\dimexpr\ht\pandoc@box+\dp\pandoc@box\relax}%
  \Gscale@div\@tempb{\linewidth}{\wd\pandoc@box}%
  \ifdim\@tempb\p@<\@tempa\p@\let\@tempa\@tempb\fi% select the smaller of both
  \ifdim\@tempa\p@<\p@\scalebox{\@tempa}{\usebox\pandoc@box}%
  \else\usebox{\pandoc@box}%
  \fi%
}
% Set default figure placement to htbp
\def\fps@figure{htbp}
\makeatother


% definitions for citeproc citations
\NewDocumentCommand\citeproctext{}{}
\NewDocumentCommand\citeproc{mm}{%
  \begingroup\def\citeproctext{#2}\cite{#1}\endgroup}
\makeatletter
 % allow citations to break across lines
 \let\@cite@ofmt\@firstofone
 % avoid brackets around text for \cite:
 \def\@biblabel#1{}
 \def\@cite#1#2{{#1\if@tempswa , #2\fi}}
\makeatother
\newlength{\cslhangindent}
\setlength{\cslhangindent}{1.5em}
\newlength{\csllabelwidth}
\setlength{\csllabelwidth}{3em}
\newenvironment{CSLReferences}[2] % #1 hanging-indent, #2 entry-spacing
 {\begin{list}{}{%
  \setlength{\itemindent}{0pt}
  \setlength{\leftmargin}{0pt}
  \setlength{\parsep}{0pt}
  % turn on hanging indent if param 1 is 1
  \ifodd #1
   \setlength{\leftmargin}{\cslhangindent}
   \setlength{\itemindent}{-1\cslhangindent}
  \fi
  % set entry spacing
  \setlength{\itemsep}{#2\baselineskip}}}
 {\end{list}}
\usepackage{calc}
\newcommand{\CSLBlock}[1]{\hfill\break\parbox[t]{\linewidth}{\strut\ignorespaces#1\strut}}
\newcommand{\CSLLeftMargin}[1]{\parbox[t]{\csllabelwidth}{\strut#1\strut}}
\newcommand{\CSLRightInline}[1]{\parbox[t]{\linewidth - \csllabelwidth}{\strut#1\strut}}
\newcommand{\CSLIndent}[1]{\hspace{\cslhangindent}#1}



\setlength{\emergencystretch}{3em} % prevent overfull lines

\providecommand{\tightlist}{%
  \setlength{\itemsep}{0pt}\setlength{\parskip}{0pt}}



 


% --- Fonts -------------------------------------------------------
% Two problems found by running pdffonts over the built book.
%
% 1. All maths was Latin Modern, not EB Garamond. The custom
%    EBGaramondMath.ttf has the glyphs (880 maths alphanumerics, the
%    angle brackets, sum, integral, geqslant) but no OpenType MATH
%    table, and unicode-math needs the table, not the glyphs: it reads
%    fraction-bar thickness, superscript drops and the big-operator
%    and extensible-delimiter variants from it. Without one it warns
%    and keeps its own default, Latin Modern Math. So the mathfont:
%    key has come out of this file, and the maths now uses
%    Garamond-Math, which is a real maths font built to accompany EB
%    Garamond. It ships with TeX Live under the OFL. Set through the
%    mathfont: key below, by family name, which is the form that works
%    in HoPF.tex.
% 2. All headings were Latin Modern Sans Bold, on 62 pages. That is
%    KOMA's default: scrbook sets sectioning in sans, and \setmainfont
%    only replaces the roman family, so the sans family was never
%    touched. disposition is KOMA's parent element for every
%    sectioning command and the table of contents, so this puts them
%    all back into EB Garamond, bold resolving to SemiBold.
\addtokomafont{disposition}{\rmfamily}

% A little extra leading. EB Garamond has a small x-height for its
% body size, so the default 1.2 baseline looks tight at this measure.
\usepackage{setspace}
\setstretch{1.05}

\usepackage{nicefrac}
\usepackage{microtype}
\usepackage{booktabs}
\usepackage{longtable}
\usepackage{etoolbox}

% Rules off entirely, as in iribookstyle.tex / quarto2024book.tex.
% Pandoc writes \toprule\noalign{}, so the commands must survive;
% zeroing the widths is the safe way to hide them.
\setlength\heavyrulewidth{0ex}
\setlength\lightrulewidth{0ex}

% Keep payoff matrices whole. Pandoc emits every table as a longtable,
% which exists to split across pages; a minipage cannot split. The
% \par and \medskip stop the minipage sitting inline with the text.
\BeforeBeginEnvironment{longtable}{%
  \par\medskip\noindent\begin{minipage}{\linewidth}\centering}
\AfterEndEnvironment{longtable}{%
  \end{minipage}\par\medskip}

% Page furniture, from quarto2024book.tex. This is the scrlayer-scrpage
% version rather than the older fancyhdr one in iribookstyle.tex;
% scrlayer-scrpage is KOMA's own and this book is scrbook.
\usepackage[automark]{scrlayer-scrpage}
\clearpairofpagestyles
\cehead{\leftmark}
\cohead{\rightmark}
\ohead{\bfseries \pagemark}
\cfoot{}
\renewcommand{\chaptermark}[1]{\markboth{#1}{}}
\renewcommand{\sectionmark}[1]{\markright{#1}}

% 17pt suited the old wide measure; 12pt is about 1em here.
\setlength{\parindent}{12pt}
\setlength{\parskip}{0ex plus 2ex minus 0ex}

% More air above headings, and a guaranteed minimum. KOMA's defaults
% are -3.5ex plus -1ex minus -.2ex for section and -3.25ex for the
% rest; the negative sign is KOMA's way of saying "no indent on the
% paragraph after the heading", so it has to be kept. Raising the
% first number adds space, and dropping the shrink to 0 stops a
% heading being pushed up against the paragraph before it when the
% page is tight. afterskip is left alone, so the gap below a heading
% does not change.
\RedeclareSectionCommand[beforeskip=-6ex plus -1ex minus -0ex]{section}
\RedeclareSectionCommand[beforeskip=-4.5ex plus -1ex minus -0ex]{subsection}
\RedeclareSectionCommand[beforeskip=-3.5ex plus -1ex minus -0ex]{subsubsection}
\raggedbottom
\widowpenalty=150
\clubpenalty=150

% Keep a list whole and attached to the sentence that introduces it.
% LaTeX's defaults actively encourage breaking there (\@itempenalty
% and \@endparpenalty are both negative out of the box), which is why
% a three-item list was splitting one item / two across a page. The
% longest list in the book is seven items and about 160 words, so
% making them unbreakable costs nothing. \@endparpenalty stays neutral
% rather than 10000, so a break is still allowed after the list.
\makeatletter
\@beginparpenalty=10000
\@itempenalty=10000
\@endparpenalty=0
\makeatother

\makeatletter
\newcommand*\NoIndentAfterEnv[1]{%
  \AfterEndEnvironment{#1}{\par\@afterindentfalse\@afterheading}}
\makeatother
\NoIndentAfterEnv{itemize}
\NoIndentAfterEnv{enumerate}
\NoIndentAfterEnv{description}
\NoIndentAfterEnv{quote}
\NoIndentAfterEnv{equation}
\NoIndentAfterEnv{longtable}

\counterwithout*{footnote}{chapter}
\makeatletter
\@ifpackageloaded{tcolorbox}{}{\usepackage[skins,breakable]{tcolorbox}}
\@ifpackageloaded{fontawesome5}{}{\usepackage{fontawesome5}}
\definecolor{quarto-callout-color}{HTML}{909090}
\definecolor{quarto-callout-note-color}{HTML}{0758E5}
\definecolor{quarto-callout-important-color}{HTML}{CC1914}
\definecolor{quarto-callout-warning-color}{HTML}{EB9113}
\definecolor{quarto-callout-tip-color}{HTML}{00A047}
\definecolor{quarto-callout-caution-color}{HTML}{FC5300}
\definecolor{quarto-callout-color-frame}{HTML}{acacac}
\definecolor{quarto-callout-note-color-frame}{HTML}{4582ec}
\definecolor{quarto-callout-important-color-frame}{HTML}{d9534f}
\definecolor{quarto-callout-warning-color-frame}{HTML}{f0ad4e}
\definecolor{quarto-callout-tip-color-frame}{HTML}{02b875}
\definecolor{quarto-callout-caution-color-frame}{HTML}{fd7e14}
\makeatother
\makeatletter
\@ifpackageloaded{bookmark}{}{\usepackage{bookmark}}
\makeatother
\makeatletter
\@ifpackageloaded{caption}{}{\usepackage{caption}}
\AtBeginDocument{%
\ifdefined\contentsname
  \renewcommand*\contentsname{Table of contents}
\else
  \newcommand\contentsname{Table of contents}
\fi
\ifdefined\listfigurename
  \renewcommand*\listfigurename{List of Figures}
\else
  \newcommand\listfigurename{List of Figures}
\fi
\ifdefined\listtablename
  \renewcommand*\listtablename{List of Tables}
\else
  \newcommand\listtablename{List of Tables}
\fi
\ifdefined\figurename
  \renewcommand*\figurename{Figure}
\else
  \newcommand\figurename{Figure}
\fi
\ifdefined\tablename
  \renewcommand*\tablename{Table}
\else
  \newcommand\tablename{Table}
\fi
}
\@ifpackageloaded{float}{}{\usepackage{float}}
\floatstyle{ruled}
\@ifundefined{c@chapter}{\newfloat{codelisting}{h}{lop}}{\newfloat{codelisting}{h}{lop}[chapter]}
\floatname{codelisting}{Listing}
\newcommand*\listoflistings{\listof{codelisting}{List of Listings}}
\makeatother
\makeatletter
\makeatother
\makeatletter
\@ifpackageloaded{caption}{}{\usepackage{caption}}
\@ifpackageloaded{subcaption}{}{\usepackage{subcaption}}
\makeatother
\usepackage{bookmark}
\IfFileExists{xurl.sty}{\usepackage{xurl}}{} % add URL line breaks if available
\urlstyle{same}
% fallback for those not using the hyperref driver hyperxmp:
\makeatletter
\@ifundefined{xmpquote}{\newcommand{\xmpquote}[1]{#1}}{}
\makeatother
\hypersetup{
  pdftitle={Game Theory and Social Choice},
  pdfauthor={Brian Weatherson},
  hidelinks,
  pdfcreator={LaTeX via pandoc}}


\title{Game Theory and Social Choice}
\usepackage{etoolbox}
\makeatletter
\providecommand{\subtitle}[1]{% add subtitle to \maketitle
  \apptocmd{\@title}{\par {\large #1 \par}}{}{}
}
\makeatother
\subtitle{Lecture Notes}
\author{Brian Weatherson}
\date{Fall 2026}
\begin{document}
\frontmatter
\maketitle

\renewcommand*\contentsname{Table of contents}
{
\setcounter{tocdepth}{2}
\tableofcontents
}

\mainmatter
\bookmarksetup{startatroot}

\chapter*{Preface}\label{preface}
\addcontentsline{toc}{chapter}{Preface}

\markboth{Preface}{Preface}

These are the lecture notes for PHIL 444, Game Theory and Social Choice,
at the University of Michigan.

They cover the game theory half of the course, and they stop where Sen
starts. The second half works through Amartya Sen's \emph{Collective
Choice and Social Welfare} (2017), and I'll do some notes for these, but
I'm not going to do a book-like summary of a book..

Giacomo Bonanno's \emph{Game Theory} (2024) is freely available online
and is worth having open alongside these notes. It covers most of the
same ground in a different voice, and it has 165 solved exercises, which
these notes do not. Bonanno is by profession an economist not a
philosopher (though he has published in philosophy journals, and is
philosophically sophisticated), so he's naturally got a different
emphasis at different points. And occasionally we disagree about the
force of some consideration or other. But apart from differences in
emphasis, the books mostly agree, and if you prefer his way of putting
things (and over the years I've heard lots of students and GSIs say they
learned a lot from his book), you should totally use it.

The notes began life in an older course, so the fit with what we are
doing now is not exact. Where a section of the book is not on the
schedule, it is there as background rather than as something you are
expected to have read. (And the title of the book came from when it was
meant to cover the whole of a course like this, not the first half.)

A PDF of the whole book is linked in the sidebar for anyone who would
rather read it in one piece.

These notes arose from notes for various different courses (mostly at
Michigan, but also at Rutgers and at St Andrews) over the years. I've
received a lot of useful feedback on them, especially from faculty
sitting in on the course at St Andrews, which has been incorporated into
the notes. Before this term I ran the whole thing through Claude, to get
advice on the prose and to see how it all fit together. It had many
suggestions, both about the organisation and the prose, and about what
needed to be added. I accepted a few of these suggestions in total, and
many of them in part. It also chased down many of the references that
I'd been cavalier about; I think I've double checked all of these, but I
might have missed some.

The book is produced in \href{https://quarto.org}{Quarto}, using HTML
and PDF templates that Claude built.

\bookmarksetup{startatroot}

\chapter{Basics of Game Theory}\label{sec-basics}

\section{Nature of Games}\label{sec-nature-of-games}

We'll think of a game as any situation where multiple players have to
make some decisions, and the outcome they get is a function of (perhaps
among other things) the decisions that each make. So understood, that's
a very large class of situations. Indeed, game theory has been used to
analyse any number of real-world interactions, as you might expect from
a definition so broad.

We're going to start our focus on a special class of games, two player
simultaneous move games. The term `two player' probably obvious; it
means that we're looking at games where just two players are involved.
Simultaneous move games are where each player only gets to make one
decision, or move. Each player makes that decision in ignorance of what
the other players have decided. And once the decisions are made, the
payoffs for each player can be revealed.

This feels like a very narrow class of games. After all, most
interactions involve multiple moves by at least one player, or at least
players making moves at different times. But there is a good sense in
which this is not a particularly tight restriction. Think of a game
where the players make multiple moves, like chess. We can imagine a game
between two chess super-computers that have planned, in advance, what
they will do at every possible stage of the game. Call such plans
\textbf{strategies}. If two such computers played, they would in a sense
each just have to make one, simultaneous, move, namely to announce their
strategies. Then it would be relatively trivial to work out the result
of the game from the strategies announced.

In practice, this isn't really possible with current technology. There
are too many possible states of the board, even in a restricted game
like chess, for a computer to have planned for every possible
eventuality. But we can imagine such complete strategies for simple
games like tic-tac-toe, and even for moderately complicated games like
checkers.

For these reasons, when we are discussing simultaneous move games, we
will often refer to the players' options as \emph{strategies}. The
representation of the game as an \(n \times m\) matrix, where one player
(represented on rows) has \(n\) options, and the other player
(represented on columns) has \(m\) options, is sometimes called a
\textbf{normal form} representation of a game and sometimes called a
\textbf{strategic form} representation. And we'll start off with these
matrix representations. When we complicate things by explicitly
representing time in the game, we'll call the representation the
\textbf{extensive form} representation of the game.

\section{Utility}\label{sec-utility}

Here is a representation of an important kind of two player simultaneous
move game.

\begin{longtable}[]{@{}lrcc@{}}
\caption{Prisoners' Dilemma}\label{tbl-pd-standard}\tabularnewline
\toprule\noalign{}
\endfirsthead
\endhead
\bottomrule\noalign{}
\endlastfoot
& & \multicolumn{2}{c@{}}{%
\textbf{Col}} \\
& & X & Y \\
\multirow{2}{*}{\textbf{Row}} & X & 3, 3 & 0, 5 \\
& Y & 5, 0 & 1, 1 \\
\end{longtable}

Each player, Row and Col, has two options: X and Y. The payoffs, given
their moves, are as follows:

\begin{itemize}
\tightlist
\item
  If both play X, both get 3.
\item
  If both play Y, both get 1.
\item
  If one plays X and the other Y, then the one who plays X gets 0, and
  the one who plays Y gets 5.
\end{itemize}

What does it mean to say that they `get' 0, 1, 3 or 5? It means that the
\textbf{utility} they get from the payout is that amount. For any given
game, the results for the players can be specified in two related, but
importantly distinct, ways.

\begin{itemize}
\tightlist
\item
  We can specify their \textbf{payouts}. These are the concrete things
  the players get.
\item
  Alternatively, we can specify the \textbf{utility} each player gets
  from their payout.
\end{itemize}

Utility is a somewhat technical, somewhat theoretical, notion.

The idea is to numerically measure how good an outcome is to an agent.
It is meant to sum up every aspect of outcomes that agents care about.
For instance, if an agent cares a lot about their friend's payout, then
their utility in an outcome will be a function not just of their own
payout, but of their friend's. So if her payout is \$\(x\) and her
friend's payout is \$\(y\), then her utility might equal \(x + 0.5y\).
This would represent that she cares about her friend's payout, but cares
a bit more (well, twice as much) about her own outcome.

The substantive assumption here is that we can sensibly map the goodness
of outcomes onto the real numbers. This assumes that goodness is, in a
sense, one-dimensional. For any two numbers \(x, y\), \(x\) is either
greater than, less than, or equal to \(y\). So one assumption we're
making is that any two outcomes are comparable in this way. We're also
assuming that we can compare the distance between pairs of outcomes. For
any numbers \(x_1, y_1, x_2, y_2\), we can ask which of \(x_1 - y_1\)
and \(x_2 - y_2\) is greater. So if we can map outcomes onto numbers,
then for any four outcomes, we can ask if the difference between the
first two in goodness is more than, less than, or the same as, the
difference in goodness between the last two. This isn't an absurd
assumption, though it can seem a little arbitrary in some cases.

Although we assume that the goodness to an agent of outcomes can be
mapped onto numbers, we do not assume that there is a unique way to do
this mapping. In fact there are two degrees of freedom in choosing a
mapping. The outcome that gets utility zero can be any point is
convenient to use. In practice, it is often convenient to give the
status quo utility zero. But note this is arbitrary, we could use
anything. The scale of the utility function is also arbitrary; we can
multiply the utilities by a positive constant and it won't change their
meaning.

Technically, if one utility function correctly represents an agent, then
so will any other function generated from it by a positive affine
transformation. So if \(u_1\) is a utility function that correctly
represents the agent, then so will the function \(u_2\) defined as
follows.

\begin{equation*}
u_2(X) = au_1(X) + b \text{, where }a\text{ is positive}.
\end{equation*}

What does it mean to say that a function \(u\) correctly represents the
agent? First, it means that the agent prefers \(X\) to \(Y\) just in
case \(u(X) > u(Y)\). Second, it means that for any \(X, Y, Z\) the
agent prefers \(Y\) to a gamble that returns \(X\) with probability
\(p\) and \(Z\) with probability \(1-p\) iff
\(u(Y) > pu(X) + (1-p)u(Z)\).

When it isn't stated otherwise, we will describe the outcomes of games
in terms of the utility each agent gets. Sometimes we will have to just
state what the payout is each agent gets, and the utility function will
be either unknown, or will have to be figured out. But the default is
that we will list the utility of each outcome.

\section{Four Examples of Games}\label{sec-four-examples}

We'll start with a very simple class of games: 2 player simultaneous
move games where each player has 2 choices for their move. As we'll see,
there is plenty of complexity within even this class.

\section{Prisoners' Dilemma}\label{sec-pd}

We've already seen this game, it was Table~\ref{tbl-pd-standard}, but
it's worthwhile having on the table again.

{\def\LTcaptype{none} % do not increment counter
\begin{longtable}[]{@{}lrcc@{}}
\toprule\noalign{}
\endhead
\bottomrule\noalign{}
\endlastfoot
& & \multicolumn{2}{c@{}}{%
\textbf{Col}} \\
& & X & Y \\
\multirow{2}{*}{\textbf{Row}} & X & 3, 3 & 0, 5 \\
& Y & 5, 0 & 1, 1 \\
\end{longtable}
}

The general form of Prisoners' Dilemma is given in
Table~\ref{tbl-pd-abstract}, assuming \(a > b > c > d\), and
\(2b > a + d\).

\begin{longtable}[]{@{}lrcc@{}}
\caption{The general form of a Prisoners' Dilemma
game}\label{tbl-pd-abstract}\tabularnewline
\toprule\noalign{}
\endfirsthead
\endhead
\bottomrule\noalign{}
\endlastfoot
& & \multicolumn{2}{c@{}}{%
\textbf{Col}} \\
& & X & Y \\
\multirow{2}{*}{\textbf{Row}} & X & \(b\), \(b\) & \(d\), \(a\) \\
& Y & \(a\), \(d\) & \(c\), \(c\) \\
\end{longtable}

One characteristic of the game is that for each player, playing \(Y\) is
better than playing \(X\) no matter what the other person does. That is,
playing \(Y\) \textbf{dominates} playing \(X\); it is better no matter
what. But, and this is what gives the game its distinctive character, it
is better for both players if they both play \(X\) than if they both
play \(Y\). (And in a repeated game it is better to always play \(X\)
than to alternate between the outcomes where players make opposite
moves.) It is an interesting way to think about modeling situations
where there are rewards to cooperation, but incentives to individualist
behavior.

\section{Stag Hunt}\label{sec-stag-hunt}

The next game we'll look at is called Stag Hunt. We'll first do a
version with commonly used numbers (in Table~\ref{tbl-sh-standard}),
then the general version.

\begin{longtable}[]{@{}lrcc@{}}
\caption{A familiar version of Stag
Hunt}\label{tbl-sh-standard}\tabularnewline
\toprule\noalign{}
\endfirsthead
\endhead
\bottomrule\noalign{}
\endlastfoot
& & \multicolumn{2}{c@{}}{%
\textbf{Col}} \\
& & X & Y \\
\multirow{2}{*}{\textbf{Row}} & X & 5, 5 & 0, 4 \\
& Y & 4, 0 & 2, 2 \\
\end{longtable}

The general version, shown in Table~\ref{tbl-sh-abstract}, has
\(a > b > c > d\) and \(b + c > a + d\)

\begin{longtable}[]{@{}lrcc@{}}
\caption{The general form of Stag
Hunt}\label{tbl-sh-abstract}\tabularnewline
\toprule\noalign{}
\endfirsthead
\endhead
\bottomrule\noalign{}
\endlastfoot
& & \multicolumn{2}{c@{}}{%
\textbf{Col}} \\
& & X & Y \\
\multirow{2}{*}{\textbf{Row}} & X & \(a\), \(a\) & \(d\), \(b\) \\
& Y & \(b\), \(d\) & \(c\), \(c\) \\
\end{longtable}

In Stag Hunt, the best outcome for all parties is the cooperative one.
But if a player has no idea what the other player will do, and thinks it
is 50/50 whether they will do the cooperative thing, then their best bet
is to not cooperate.

The cute story behind this game comes from the eighteenth century
political philosopher Jean-Jacques Rousseau (Rousseau {[}1755{]} 2004).
Imagine two people are setting out on the day's hunt, and are going to
hunt stag or rabbits. Catching a stag will lead to more meat, but it is
impossible with just one person. One person can catch a lot of rabbits,
especially if no one else is hunting rabbits. In this case, X is joining
the stag hunt, and Y is setting off on a solo rabbit hunt. It isn't as
hard to get cooperative outcomes in this game as in Prisoners' Dilemma,
but it takes some background trust.

\section{Battle of the Sexes}\label{sec-battle-of-sexes}

I don't particularly like the name of this game, but it is common enough
that you probably should know it. Table~\ref{tbl-bs-standard} gives a
commonly used version of the game.

\begin{longtable}[]{@{}lrcc@{}}
\caption{Battle of the Sexes}\label{tbl-bs-standard}\tabularnewline
\toprule\noalign{}
\endfirsthead
\endhead
\bottomrule\noalign{}
\endlastfoot
& & \multicolumn{2}{c@{}}{%
\textbf{Col}} \\
& & X & Y \\
\multirow{2}{*}{\textbf{Row}} & X & 2, 1 & 0, 0 \\
& Y & 0, 0 & 1, 2 \\
\end{longtable}

Table~\ref{tbl-bs-abstract} shows the general version, with
\(a > b > c\),

\begin{longtable}[]{@{}lrcc@{}}
\caption{The abstract form of Battle of the
Sexes}\label{tbl-bs-abstract}\tabularnewline
\toprule\noalign{}
\endfirsthead
\endhead
\bottomrule\noalign{}
\endlastfoot
& & \multicolumn{2}{c@{}}{%
\textbf{Col}} \\
& & X & Y \\
\multirow{2}{*}{\textbf{Row}} & X & \(a\), \(b\) & \(c\), \(c\) \\
& Y & \(c\), \(c\) & \(b\), \(a\) \\
\end{longtable}

The thought is that both parties prefer coordination - making the same
pick - to doing different things. But they differ on what they would
like to coordinate on; Row would like to coordinate on X, and Column
would like to coordinate on Y. The name comes from the version of the
game where Row and Column are husband and wife, and X and Y are
stereotypically male/female things to prefer doing. You can probably
fill in your own example here.

\section{Chicken}\label{sec-chicken}

Our last game, shown in Table~\ref{tbl-chicken-standard}, is not for the
faint of heart.

\begin{longtable}[]{@{}lrcc@{}}
\caption{A familiar form of
Chicken}\label{tbl-chicken-standard}\tabularnewline
\toprule\noalign{}
\endfirsthead
\endhead
\bottomrule\noalign{}
\endlastfoot
& & \multicolumn{2}{c@{}}{%
\textbf{Col}} \\
& & X & Y \\
\multirow{2}{*}{\textbf{Row}} & X & -100, -100 & 1, -1 \\
& Y & -1, 1 & 0, 0 \\
\end{longtable}

The idea comes from a James Dean movie. Two cars are lined up facing
each other at the end of a long straight. They start driving at each
other. The drivers can either stay on the road (that's X), or swerve to
the right (that's Y). If they both stay on the road, they die in a fiery
inferno. If one swerves and the other stays, the one who stays wins the
game, and the other loses. If they both swerve, the game is a draw.

Unlike the previous games we've looked at, there is no pure
\textbf{equilibrium} to this game. That is, there is no point in the
game where no one regrets their choice on finding out what the other
players choose. There is an equilibrium to the game, it turns out, but
getting to it requires a bit of background.

\section{Types of Games}\label{sec-types-of-games}

Let \(u(x, o)\) be the utility that player \(x\) gets in outcome \(o\).
Then a two player game, with players \(p_1, p_2\) is \textbf{zero-sum}
if there is some constant \(c\) such that for any outcome \(o\),
\(u(p_1, o) + u(p_2, o) = c\). That is, however good an outcome is for
\(p_1\), it is equally bad for \(p_2\), and vice versa. Pure
competitions are zero-sum, though real zero-sum interactions are rare in
the real world.

By convention, when a game is zero-sum, we rescale the utilities so the
constant \(c\) is 0. We always can do this - remember any positive
affine transformation of the utility function doesn't change its
meaning. That's why we can talk about zero-sum games, and not just
constant sum games.

You might worry that whether a game is zero-sum or not depends on
completely arbitrary choices about how to scale the utility functions.
Strictly speaking, a game is zero-sum if there is \emph{any} pair of
utility functions that correctly represents the agents, and which makes
the utility each gets from any outcome sum to a constant.

A game is \textbf{cooperative} if in any outcome, each player gets the
same utility. As with the definition of zero-sum, it doesn't matter
whether the game looks cooperative on every way of representing the
utilities of the players, just if it looks cooperative on some
representation.

\section{Dominant Strategies}\label{sec-dominant-strategies}

Say that one player has available to them strategies
\(s_1, \dots, s_n\), and the other players have available to them
collectively strategies \(v_1, \dots, v_m\). In the case of a two player
game, the strategies available to other players will just be strategies.
But in games with more than two players, we'll think of the alternatives
available to other players as \emph{vectors}, with each dimension being
a choice available to one of the other players. So \(v_17\) might be
that player 2 plays X, and player 3 plays Y, and player 4 plays Z, and
player 5 also plays Z. Every one of these vectors specifies the choices
of all other players in the game.

Let \(u_x(s_i, v_k)\) be the utility that \(x\), the player we're
focussing on, gets if she plays \(s_i\) and the others play \(v_k\).
Then strategy \(s_i\) \textbf{strongly dominates} strategy \(s_j\) if
and only if the following condition is satisfied.

\begin{equation*}
\forall m: u_x(s_1, v_m) > u_x(s_2, v_m)
\end{equation*}

That is, no matter what the other players end up playing, \(x\) will do
better playing \(s_1\) than playing \(s_2\).

We'll say that strategy \(s_i\) \textbf{weakly dominates} strategy
\(s_j\) if and only if the following condition is satisfied.

\begin{equation*}
\forall m: u_x(s_1, v_m) \geq u_x(s_2, v_m) \& \exists m: u_x(s_1, v_m) > u_x(s_2, v_m)
\end{equation*}

That is, playing \(s_1\) can't do any worse than playing \(s_2\), no
matter what the others do, and it could do better.

\section{Dominance and Solving Games}\label{sec-dominance-solving}

Here is one simple way to make progress on a game.

\begin{enumerate}
\def\labelenumi{\arabic{enumi}.}
\tightlist
\item
  Assume that no player will play a strongly dominated strategy.
\item
  Assume that every player will make a rational choice given that no
  player will play a dominated strategy.
\end{enumerate}

The first step here lets us solve Prisoners' Dilemma. Since \(Y\)
dominates \(X\), neither player will play \(X\). So the solution to the
game is that both players play \(Y\). Table~\ref{tbl-dominance-example}
is another game that isn't solved by the first step, but is solved by
the second step.

\begin{longtable}[]{@{}lrcc@{}}
\caption{An example of dominance
reasoning}\label{tbl-dominance-example}\tabularnewline
\toprule\noalign{}
\endfirsthead
\endhead
\bottomrule\noalign{}
\endlastfoot
& & \multicolumn{2}{c@{}}{%
\textbf{Col}} \\
& & X & Y \\
\multirow{2}{*}{\textbf{Row}} & X & 2, 1 & 3, 0 \\
& Y & 0, 0 & 1, 1 \\
\end{longtable}

In Table~\ref{tbl-dominance-example} Column has no dominant strategy;
what's best for them to play is a function of what Row plays. But Row
does have a dominant strategy. Whatever Column does, they are better off
playing Y. So we should assume they will play Y. And given they play Y,
Column should also play Y. So the solution to the game again is that
both players play Y.

In later chapters, we'll look at what happens when we iterate this idea.
But first we need an introduction to probability theory.

\bookmarksetup{startatroot}

\chapter{Probability and Decision Theory}\label{sec-decision-theory}

\section{Probability}\label{sec-probability}

We'll start by going over the notion of probability that is used in
decision theory. The best way to think about probability functions is by
thinking about \textbf{measure} functions. Intuitively, a measure
function is a non-negative \textbf{additive} function. That is, it is a
function which only takes non-negative values, and where the value of a
whole is equal to the sum of the value of the (non-overlapping) parts.
In principle the domain of a measure function could be the size of any
set, but we'll only consider functions with finite domains for now.

Such functions are frequently used in the real world. Consider, for
instance, the functions from regions to their population and land mass.
Table~\ref{tbl-uk-geography} lists the (very approximate) population and
land mass for England, Scotland, Wales and Northern Ireland.

\begin{longtable}[]{@{}lcc@{}}
\caption{Some statistics about the United
Kingdom}\label{tbl-uk-geography}\tabularnewline
\toprule\noalign{}
\endfirsthead
\endhead
\bottomrule\noalign{}
\endlastfoot
& Population & Land Mass \\
& (Millions) & (Thousands of mi\(^2\)) \\
England & 50 & 50 \\
Scotland & 5 & 30 \\
Wales & 3 & 8 \\
Northern Ireland & 2 & 5 \\
\end{longtable}

The nice thing about tables like Table~\ref{tbl-uk-geography} is that
they give you a lot of information implicitly. You don't need a separate
table to tell you that the combined population of England and Scotland
is 55 million, or that the combined land mass of England and Wales is
58,000 square miles. Those (approximate) facts follow from the facts
listed on the table, as well as the fact that populations and land
masses are additive. Once you know the population of the parts, you know
the population of the whole.

The next notion we need is that of a \textbf{normalised} measure
function. A normalised measure is a measure where every function takes
the value assigned to the largest whole is 1. We often use this notion
when talking about proportions. Table~\ref{tbl-uk-geography-normal} is a
table listing, again approximately, the portion of the UK population and
land mass in each of the four countries.

\begin{longtable}[]{@{}lcc@{}}
\caption{A normalised verison of
Table~\ref{tbl-uk-geography}.}\label{tbl-uk-geography-normal}\tabularnewline
\toprule\noalign{}
\endfirsthead
\endhead
\bottomrule\noalign{}
\endlastfoot
& Population & Land Mass \\
& (Proportion) & (Proportion) \\
England & 0.84 & 0.54 \\
Scotland & 0.08 & 0.32 \\
Wales & 0.05 & 0.09 \\
Northern Ireland & 0.03 & 0.05 \\
\end{longtable}

The way we turn the regular measure function into a normalised measure
is by dividing each value in Table~\ref{tbl-uk-geography} by the value
of whole. So the land mass of England is (about) 50,000 mi\(^2\), and
the land mass of the UK is (about) 93,000 mi\(^2\), and
\(\nicefrac{50,000}{93,000}\) is about 0.54. So that's why we right 0.54
in the top right cell of the normalised measure table.

The last notion we need is that of a possibility space, or more
formally, a \textbf{field} of propositions. Although these terms might
be unfamiliar, the notion should be very familiar; it is just what we
represent with a truth table. We will identify, for these purposes,
propositions with sets of possible worlds. Then a field of propositions
is a set of propositions such that whenever \(A\) and \(B\) are in the
set, so are \(\neg A\) and \(A \wedge B\). This is, of course, just what
we see with truth tables. Whenever you can represent two propositions on
a truth table, you can represent their conjunction and each of their
negations. Just like with truth-tables, we'll only consider finite
fields. It's possible to extend this to cover any set-sized partition of
possibility space, but we'll ignore such complications for now.

A \textbf{Probability function} is just a normalised measure on a
possibility space. That is, it is a function satisfying these two
conditions, where \(A\) and \(B\) are any propositions, and \(T\) is a
logical truth.

\begin{itemize}
\tightlist
\item
  \(0 \leq \Pr(A) \leq 1 = \Pr(T)\)
\item
  If \(A\) and \(B\) can't both be true, then
  \(\Pr(A \vee B) = \Pr(A) + \Pr(B)\).
\end{itemize}

We normally suppose that any rational agent has \textbf{credences} over
the space of (salient) epistemic possibilities, and these credences form
a probability function.

In any finite field of propositions, there will be some \textbf{atoms}.
These are propositions such that no stronger consistent proposition is
in the field. (Some, but not all, infinite fields are atomic.) We will
use these atoms a lot. I'll use the variable \(w\) to range over atoms,
reminding us that for practical purposes, atoms behave a lot like
\emph{worlds}. They aren't worlds, of course; they only settle the truth
values of a handful of propositions. But they are divisions of
possibility space that are maximally precise for present purposes. That
is, they are being treated as worlds. On a pragmatist approach to
modality, there's no more to being a world than to be treated as a world
in the right circumstances, and thinking about probability encourages, I
think, such a pragmatist approach.

\section{Conditional Probability}\label{sec-conditional-probability}

So far we've talked simply about the probability of various
propositions. But sometimes we're not interested in the absolute
probability of a proposition, we're interested in its
\textbf{conditional} probability. That is, we're interested in the
probability of the proposition \emph{assuming} or \emph{conditional on}
some other proposition obtaining.

It isn't too hard to visualise how conditional probability works if we
think of a possibility space simply as a truth table, so a measure over
a possibility space is simply a measure over on the truth table.
Table~\ref{tbl-tt-first} how we might represent a simple probability
function, defined over the Boolean combinations of three propositions,
if we're thinking about things this way.

\begin{longtable}[]{@{}cccl@{}}
\caption{A truth table, with
measures}\label{tbl-tt-first}\tabularnewline
\toprule\noalign{}
\(p\) & \(q\) & \(r\) & \\
\midrule\noalign{}
\endfirsthead
\toprule\noalign{}
\(p\) & \(q\) & \(r\) & \\
\midrule\noalign{}
\endhead
\bottomrule\noalign{}
\endlastfoot
T & T & T & 0.0008 \\
T & T & F & 0.008 \\
T & F & T & 0.08 \\
T & F & F & 0.8 \\
F & T & T & 0.0002 \\
F & T & F & 0.001 \\
F & F & T & 0.01 \\
F & F & F & 0.1 \\
\end{longtable}

The way to read Table~\ref{tbl-tt-first} is that each row defines a
world. The characteristics of the world are just the truth values of
three salient propositions: \(p, q, r\). If T is shown underneath the
name of a proposition, that means the proposition is true in the world,
and if an F is shown, it mans the proposition is false. So the third
line is the world where \(p\) and \(r\) are true, and \(q\) is false. To
the right of the vertical line is the probability of that world. We can
work out the probability of any proposition that is true at multiple
worlds by adding up the probabilities of the worlds where it is true. So
\(\Pr(q) = 0.0008 + 0.008 + 0.0002 + 0.001 = 0.01\).

Now let's see what happens when we look at conditional probabilities,
i.e., probabilities you get from assuming some particular proposition is
true. If we assume that \(B\) is true, then we should `zero out',
i.e.~assign probability 0, to all the possibilities where \(B\) doesn't
obtain. We're now left with a measure over only the \(B\)-possibilities.
The problem is that it isn't a normalised measure. The values will only
sum to \(\Pr(B)\), not to 1. We need to renormalise. So we divide by
\(\Pr(B)\) and we get a probability back. In a formula, we're left with

\begin{equation*}
Pr(A|B) = \frac{\Pr(A \wedge B)}{\Pr(B)}
\end{equation*}

Assume now that we're trying to find the conditional probability of
\(p\) given \(q\) in the above table. We could do this in two different
ways. First, we could set the probability of any line where \(q\) is
false to 0. So we will get Table~\ref{tbl-tt-second}.

\begin{longtable}[]{@{}cccl@{}}
\caption{What happens to Table~\ref{tbl-tt-first} if we 0 out
\(\neg q\)}\label{tbl-tt-second}\tabularnewline
\toprule\noalign{}
\(p\) & \(q\) & \(r\) & \\
\midrule\noalign{}
\endfirsthead
\toprule\noalign{}
\(p\) & \(q\) & \(r\) & \\
\midrule\noalign{}
\endhead
\bottomrule\noalign{}
\endlastfoot
T & T & T & 0.0008 \\
T & T & F & 0.008 \\
T & F & T & 0 \\
T & F & F & 0 \\
F & T & T & 0.0002 \\
F & T & F & 0.001 \\
F & F & T & 0 \\
F & F & F & 0 \\
\end{longtable}

The numbers don't sum to 1 any more. They sum to 0.01. So we need to
divide everything by 0.01. It's sometimes easier to conceptualise this
as multiplying by \(\nicefrac{1}{\Pr(q)}\), i.e.~by multiplying by 100.
Then we'll end up with Table~\ref{tbl-tt-third}.

\begin{longtable}[]{@{}cccl@{}}
\caption{What happens to Table~\ref{tbl-tt-second} if
renormalise}\label{tbl-tt-third}\tabularnewline
\toprule\noalign{}
\(p\) & \(q\) & \(r\) & \\
\midrule\noalign{}
\endfirsthead
\toprule\noalign{}
\(p\) & \(q\) & \(r\) & \\
\midrule\noalign{}
\endhead
\bottomrule\noalign{}
\endlastfoot
T & T & T & 0.08 \\
T & T & F & 0.8 \\
T & F & T & 0 \\
T & F & F & 0 \\
F & T & T & 0.02 \\
F & T & F & 0.1 \\
F & F & T & 0 \\
F & F & F & 0 \\
\end{longtable}

And since \(p\) is true on the top two lines of
Table~\ref{tbl-tt-third}, the `new' probability of \(p\) is 0.88. That
is, the conditional probability of \(p\) given \(q\) is 0.88. As we were
writing things above, \(\Pr(p | q) = 0.88\).

Alternatively we could just use the formula given above. Just adding up
rows gives us the following numbers.

\begin{align*}
\Pr(p \wedge q) &= 0.0008 + 0.008 = 0.0088 \\
\Pr(q) &= 0.0008 + 0.008 + 0.0002 + 0.001 = 0.01
\end{align*}

Then we can apply the formula.

\begin{align*}
\Pr(p | q) &= \frac{\Pr(p \wedge q)}{\Pr(q)} \\
&= \frac{0.0088}{0.01} \\
&= 0.88
\end{align*}

\section{Bayes Theorem}\label{sec-bayes-theorem}

This is a bit of a digression from what we're really focussing on in
game theory, but it's important for understanding how conditional
probabilities work. It is often easier to calculate conditional
probabilities in the `inverse' direction to what we are interested in.
That is, if we want to know \(\Pr(A | B)\), it might be much easier to
discover \(\Pr(B | A)\). In these cases, we use Bayes Theorem to get the
right result.\footnote{Bayes Theory is named after the
  \href{https://en.wikipedia.org/wiki/Thomas_Bayes}{Reverend Thomas
  Bayes}, an 18th century statistician and nonconformist minister.} I'll
state Bayes Theorem in two distinct ways, then show that the two ways
are ultimately equivalent.

\begin{align*}
\Pr(A|B) &= \frac{\Pr(B|A) \Pr(A)}{\Pr(B)} \\
&= \frac{\Pr(B|A) \Pr(A)}{\Pr(B|A)Pr(A) + \Pr(B|\neg A)Pr(\neg A)}
\end{align*}

These are equivalent because
\(\Pr(B) = \Pr(B|A)Pr(A) + \Pr(B|\neg A)\Pr(\neg A)\). Since this is an
independently interesting result, it's worth going through the proof of
it. First note that

\begin{align*}
\Pr(B|A)\Pr(A) &= \frac{\Pr(A \wedge B)}{\Pr(A)} \Pr(A) \\
&= \Pr(A \wedge B) \\
\\
\Pr(B|\neg A)\Pr(\neg A) &= \frac{\Pr(\neg A \wedge B)}{\Pr\neg (A)} \Pr(\neg A) \\
&= \Pr(\neg A \wedge B)
\end{align*}

Adding those two together we get

\begin{align*}
\Pr(B|A)\Pr(A) + \Pr(B|\neg A)\Pr(\neg A) &= \Pr(A \wedge B) + \Pr(\neg A \wedge B) \\
&= \Pr((A \wedge B) \vee (\neg A \wedge B)) \\
&= \Pr(B)
\end{align*}

The second line uses the fact that \(A \wedge B\) and
\(\neg A \wedge B\) are inconsistent. And the third line uses the fact
that \((A \wedge B) \vee (\neg A \wedge B)\) is equivalent to \(A\). So
we get a nice result, one that we'll have occasion to use a bit in what
follows.

\begin{equation*}
\Pr(B) = \Pr(B|A)\Pr(A) + \Pr(B|\neg A)\Pr(\neg A)
\end{equation*}

So the two forms of Bayes Theorem are the same. We'll often find
ourselves in a position to use the second form.

One kind of case where we have occasion to use Bayes Theorem is when we
want to know how significant a test finding is. So imagine we're trying
to decide whether the patient has disease D, and we're interested in how
probable it is that the patient has the disease conditional on them
returning a test that's positive for the disease. We also know the
following background facts.

\begin{itemize}
\tightlist
\item
  In the relevant demographic group, 5\% of patients have the disease.
\item
  When a patient has the disease, the test returns a position result
  80\% of the time
\item
  When a patient does not have the disease, the test returns a negative
  result 90\% of the time
\end{itemize}

So in some sense, the test is fairly reliable. It usually returns a
positive result when applied to disease carriers. And it usually returns
a negative result when applied to non-carriers. But as we'll see when we
apply Bayes Theorem, it is very unreliable in another sense. So let
\(A\) be that the patient has the disease, and \(B\) be that the patient
returns a positive test. We can use the above data to generate some
`prior' probabilities, i.e.~probabilities that we use prior to getting
information about the test.

\begin{itemize}
\tightlist
\item
  \(\Pr(A) = 0.05\), and hence \(\Pr(\neg A) = 0.95\)
\item
  \(\Pr(B | A) = 0.8\)
\item
  \(\Pr(B | \neg A) = 0.1\)
\end{itemize}

Now we can apply Bayes theorem in its second form.

\begin{align*}
\Pr(A|B) &= \frac{\Pr(B|A) \Pr(A)}{\Pr(B|A)\Pr(A) + \Pr(B|\neg A)\Pr(\neg A)} \\
&= \frac{0.8 \times 0.05}{0.08 \times 0.05 + 0.1 \times 0.95} \\
&= \frac{0.04}{0.04 + 0.095} \\
&= \frac{0.04}{0.135} \\
&\approx 0.296
\end{align*}

So in fact the probability of having the disease, conditional on having
a positive test, is less than 0.3. So in that sense the test is quite
unreliable.

This is actually a quite important point. The fact that the probability
of \(B\) given \(A\) is quite high does not mean that the probability of
\(A\) given \(B\) is equally high. By tweaking the percentages in the
example I gave you, you can come up with cases where the probability of
\(B\) given \(A\) is arbitrarily high, even 1, while the probability of
\(A\) given \(B\) is arbitrarily low.

Confusing these two conditional probabilities is sometimes referred to
as the \emph{Prosecutors' fallacy}, though it's not clear how many
actual prosecutors are guilty of it. The thought is that some
prosecutors start with the premise that the probability of the
defendant's blood (or DNA or whatever) matching the blood at the crime
scene, conditional on the defendant being innocent, is 1 in a billion
(or whatever it exactly is). They conclude that the probability of the
defendant being innocent, conditional on their blood matching the crime
scene, is about 1 in a billion. Because of derivations like the one we
just saw, that is a clearly invalid move.

A more striking real world example comes from the analysis of soccer
carried out by \href{https://en.wikipedia.org/wiki/Charles_Reep}{Charles
Reep} and
\href{https://en.wikipedia.org/wiki/Charles_Hughes_(football_manager)}{Charles
Hughes}. (Reep and Hughes worked separately, but they reached similar
conclusions by similar means.) They watched many games, and tracked each
of the possessions (meaning a period of time where one team has control
of the ball). They counted how many passes took place in each
possession, and whether the possession resulted in a goal. They noted
that the vast majority of the goals came from very short possessions,
ones involving three or fewer passes; Reep put the figure at around 80\%
(Wilson 2008, 141). His own account of the study is Reep and Benjamin
(1968). From this they concluded that soccer teams should be direct;
kicking the ball to forwards as soon as possible to try to score from a
short possession. Hughes became technical director of the Football
Association in England, where he spread this idea of how to play soccer
across the country.

You may be able to spot the error in their reasoning already. The figure
that makes it vivid is one Reep did not dwell on: in the games he
recorded, 91.5\% of moves involved three passes or fewer (Wilson 2008,
141). So the share of goals that came from short possessions was
\emph{lower} than the share of possessions that were short, which is the
opposite of what you would expect if short possessions were the
productive ones. To see why that matters, work through the problem with
some numbers. (Those in Table~\ref{tbl-soccer} are invented, and the
split between medium and long is there to show a trend rather than to
report one, but the overall proportions are close to what Reep found.)
Divide all possessions into short (1-3 passes), medium (4-7 passes) and
long (8+ passes). Imagine that we observed a bunch of games, and
recorded the data in Table~\ref{tbl-soccer}.

\begin{longtable}[]{@{}lcc@{}}
\caption{Some imagined data for how soccer goals are
scored}\label{tbl-soccer}\tabularnewline
\toprule\noalign{}
& Goal & No Goal \\
\midrule\noalign{}
\endfirsthead
\toprule\noalign{}
& Goal & No Goal \\
\midrule\noalign{}
\endhead
\bottomrule\noalign{}
\endlastfoot
Short & 78 & 822 \\
Medium & 12 & 48 \\
Long & 10 & 30 \\
\end{longtable}

As you can see, most goals come from short possessions. But you
shouldn't get excited by finding your team has had a short possession;
over 90\% of short possessions go nowhere. On the other hand, medium to
long possessions often end in goals. In Table~\ref{tbl-soccer-percent}
I've added another column to Table~\ref{tbl-soccer} to make this clear.

\begin{longtable}[]{@{}lccc@{}}
\caption{A version of Table~\ref{tbl-soccer} with success probability
added.}\label{tbl-soccer-percent}\tabularnewline
\toprule\noalign{}
& Goal & No Goal & Success \% \\
\midrule\noalign{}
\endfirsthead
\toprule\noalign{}
& Goal & No Goal & Success \% \\
\midrule\noalign{}
\endhead
\bottomrule\noalign{}
\endlastfoot
Short & 78 & 822 & 8.7 \\
Medium & 12 & 48 & 20 \\
Long & 10 & 30 & 25 \\
\end{longtable}

The longer a possession is, the higher the frequency that it ends in a
goal. I'm not sure in real soccer it is as dramatic as the numbers I
gave you. (And in fact at high levels of soccer not nearly this high a
percentage of possessions end in goals.) But I am confident that (at
least at high levels, at most times in the history of soccer) the broad
pattern I've outlined here is correct. Most goals come from short
possessions, but short possessions lead to fewer goals, on average, than
long possessions. Anyone who has taken a probability class should know
that this suggests aiming for longer possessions rather than shorter
ones, but apparently the England hierarchy didn't take that lesson for
several decades.

\section{Expected Value}\label{sec-expected-value}

One very convenient notion to define using a probability function is
that of an \textbf{expected value}. Say that a \textbf{random variable}
is any function from possibilities to real numbers. We usually designate
random variables with capital letters. So expressions like \(X=x\),
which you'll often see in probability textbooks, means that the random
variable \(X\) takes real value \(x\).

The expected value of a random variable is calculated using the
following formula.

\begin{equation*}
E(X) = \sum_w \Pr(w)X(w)
\end{equation*}

Remember that \(w\) is a variable that ranges over atoms. So the
summation here is over the atoms in the field. \(\Pr(w)\) is the
probability of that atom obtaining, and \(X(w)\) is the value of the
random variable at \(w\). Here's a worked example to make it clearer.

Let our random variable \(X\) be the number of children of the next
person to walk through the door. Let's assume that the department
contains 10 people; 2 of whom have three children, 3 of whom have two
children, 3 of whom have one child, and 2 are childless. Each member of
the department is equally likely to walk through the door, and no one
else is going to. So the probability of any given member of the
department walking through the door is \(\nicefrac{1}{10}\). Then the
expected value of \(X\), i.e., \(E(X)\), is:

\begin{align*}
\nicefrac{1}{10} \times 3 + \nicefrac{1}{10} \times 3 &+ \\
\nicefrac{1}{10} \times 2 + \nicefrac{1}{10} \times 2 + \nicefrac{1}{10} \times 2 &+ \\
\nicefrac{1}{10} \times 1 + \nicefrac{1}{10} \times 1 + \nicefrac{1}{10} \times 1 &+ \\
\nicefrac{1}{10} \times 0 + \nicefrac{1}{10} \times 0 &= 1.5 \\
\end{align*}

Note that the expected value of \(X\) is \emph{not} the value we
`expect', in the ordinary sense, \(X\) to take. We can be pretty sure
that \(X\)'s value will be an integer, so it won't be 1.5, for instance!
That is to say, `expected value' is a term of art.

As well as `absolute' expected values, we can also define conditional
expected values, as follows:

\begin{equation*}
E(X | p) = \sum_w \Pr(w | p)X(w)
\end{equation*}

I won't go through the proof of this, but it turns out that conditional
expected values satisfy a nice theorem. Let \(\{p_1, \dots, p_n\}\) be a
partition of possibility space. Then the following holds.

\begin{equation*}
E(X) = \sum_{i=1}^n \Pr(p_i) E(X | p_i)
\end{equation*}

Here are a few of immediate corrollaries of that result.

\begin{align*}
E(X) -E(Y) &= E(X - Y) \\
&= \sum_{i=1}^n \Pr(p_i) (E(X | p_i) - E(Y | p_i)) \\
&= \sum_{i=1}^n \Pr(p_i) (E((X - Y) | p_i))
\end{align*}

Often we know more about the difference between two expected values than
we do about the absolute value of either. For instance, if we are
comparing the expected values of taking and declining a small bet, we
can figure out which action has a higher expected value without figuring
out the expected value of either, which presumably depends on all sorts
of facts about how our life might (probably) go. But as long as we can
work out the difference in expected value between the two acts
conditional on each \(p_i\), and the probability of each \(p_i\), we can
work out at least which has a higher expected value, and by how much.

\section{Orthodox Decision Theory}\label{sec-orthodox-decision-theory}

Given all that, it is fairly easy to summarise orthodox decision theory.

\begin{itemize}
\tightlist
\item
  Rational agents have a credence function over salient alternatives
  that is identical to some probability function \(\Pr\).
\item
  Rational agents also have a real-valued utility function \(U\) over
  salient outcomes.
\item
  The rational action for such an agent to do is the action that
  maximises \textbf{expected} utility, i.e., maximises the expected
  value of \(U\) relative to that function \(\Pr\).
\end{itemize}

\section{Money Is Not Utility}\label{sec-money-not-utility}

Chapter~\ref{sec-basics} separated a player's payout from their utility,
and said that we would normally specify games using the latter. This
section is about why that separation matters, because the difference
between money and utility plays several different explanatory roles.

Getting \$2x is not twice as good as getting \$x. Getting \$2x is like
getting \$x, and then getting another \$x when you are already \$x
richer, and the second \$x does you less good than the first did. This
what is known as the \textbf{declining marginal utility} of money. The
\textbf{marginal} utility of money is the margin between the utility you
had before and after getting that money. A core thought of economic
theory since the 1870s is that marginal utility is the core explanation
of action; absolute utility is either non-existent or
unimportant.\footnote{Like many things, the importance of marginal
  utility seems to have been independently discovered by many people at
  once. The philosophy connection here is that one of the independent
  discoverers was William Stanley Jevons (1871). I'd like to think he
  came up with the idea while working in Australia in the 1850s, which
  would make this one of the earliest Australian contributions to
  Anglophone philosophy. But as far as I know, there is literally zero
  evidence for that timeline.} And the marginal utility of money falls,
i.e., declines, the more of it you have.

Two claims that sound quite different both fall out of this. You are
better off taking a million dollars than a fifty-fifty shot at two
million. And a million dollars changes your life more than it changes a
billionaire's. The orthodox view (and we'll see reasons that not
everyone buys this when we get to Section~\ref{sec-allais}) is that
these are just two features of the same thing, that money has a
declining marginal utility.

If we graph utility against wealth, declining marginal utility fixes two
things about the shape. The curve slopes upward, because more money is
better than less. And it bends downward; it is concave.

Here is a crude model. Let utility be the square root of net wealth in
dollars, and take three equal steps of \$300.

\begin{itemize}
\tightlist
\item
  Utility at \$100 is \(\sqrt{100} = 10\).
\item
  Utility at \$400 is \(\sqrt{400} = 20\).
\item
  Utility at \$700 is \(\sqrt{700} \approx 26.46\).
\item
  Utility at \$1000 is \(\sqrt{1000} \approx 31.62\).
\end{itemize}

The same \$300 buys 10 units of utility, then 6.46, then 5.17. Put the
other wayround, equal gains in utility take larger and larger amounts of
money.

For wealth above a thousand dollars or so, a logarithm is a more
plausible model than a square root. Let utility be \(\log_{10}\) of net
wealth.

\begin{itemize}
\tightlist
\item
  Utility at \$10,000 is 4.
\item
  Utility at \$100,000 is 5.
\item
  Utility at \$1,000,000 is 6.
\end{itemize}

\emph{Multiplying} your wealth by ten \emph{adds} one to your utility.
The relation between dollars and utility is not just nonlinear, it is
dramatically so.

That has consequences that are easy to miss. Take someone whose net
wealth is \$100,000, so their utility is 5, and offer them a bet that
pays \$900,000 if it wins and costs \$90,000 if it loses, on the flip of
a fair coin.

\begin{align*}
E(U(\text{bet})) &= \nicefrac{1}{2} \times \log_{10}(1{,}000{,}000) + \nicefrac{1}{2} \times \log_{10}(10{,}000) \\
&= \nicefrac{1}{2} \times 6 + \nicefrac{1}{2} \times 4 \\
&= 5
\end{align*}

The bet has exactly the expected utility of doing nothing, so they
should be indifferent, even though its expected dollar return is
\$405,000.

Now shrink the prize to \$800,000, leaving everything else alone.

\begin{align*}
E(U(\text{bet})) &= \nicefrac{1}{2} \times \log_{10}(900{,}000) + \nicefrac{1}{2} \times \log_{10}(10{,}000) \\
&\approx \nicefrac{1}{2} \times 5.954 + \nicefrac{1}{2} \times 4 \\
&\approx 4.977
\end{align*}

That is below the utility of the status quo, so they should decline, and
the expected dollar return is still \$355,000. Expected dollars and
expected utility do not merely differ in size. They can point in
opposite directions. This fact explains why there is an insurance
industry; an insurance contract (from the consumer's perspective)
typically has a negative expected dollar return, and a positive expected
utility return.

\section{The Allais Paradox}\label{sec-allais}

Some people have the intuition that expected utility theory doesn't
sufficiently account for the preference people have for avoiding risk.
After all, we routinely turn down bets with a positive expected dollar
return. The orthodox reply (which I think is completely right, for what
it's worth) is that these preferences over individual bets (when they
are rational) can be explained by the declining marginal utility of
money, plus the view that rationality requires maximising expected
utility. No matter what bets people turn down, if we want to rationalise
them, we can by positing a steeper declining marginal utility.

It turns out, as Maurice Allais (1953) showed\footnote{\href{https://en.wikipedia.org/wiki/Maurice_Allais}{Allais}
  was a French economist. He won the 1988 Nobel, though not primarily
  for the result that he's most famous for in the Anglophone world,
  i.e., this paradox.}, that the theory does have some striking
predictions not for individual betting dispositions, but for
\emph{pairs} of betting dispositions. He came up with a pair of binary
choices, and showed that there was no way for expected utility theory to
explain why one attractive pair of dispositions to choose could be
rationalised.

Here is the first choice.

\begin{itemize}
\tightlist
\item
  \textbf{1A}: \$1,000,000, guaranteed.
\item
  \textbf{1B}: a 1\% chance of \$0, an 89\% chance of \$1,000,000, and a
  10\% chance of \$5,000,000.
\end{itemize}

Many people, I suspect most, take \textbf{1A}.\footnote{The empirical
  evidence on the Allais paradox is actually a bit of a mess; see
  Blavatskyy et al. (2022) for a good overview. As I understand it, in
  these high stakes cases, this part at least is very standard. Of
  course, these are hypothetical choices; we don't run real experiments
  where we give away millions of dollars.} A guaranteed million is not
worth risking for a one in ten shot at five, when the risk is walking
away with nothing. That doesn't maximise expected dollars, but it's
fine; as long as the marginal utility of the first million is
sufficiently greater than the marginal utility of the next four million,
\textbf{1A} could maximise expected utility.

Here is the second.

\begin{itemize}
\tightlist
\item
  \textbf{2A}: an 89\% chance of \$0, and an 11\% chance of \$1,000,000.
\item
  \textbf{2B}: a 90\% chance of \$0, and a 10\% chance of \$5,000,000.
\end{itemize}

Now it seems fine to take \textbf{2B}, and many people do.\footnote{Notably,
  as Blavatskyy et al. (2022) point out, it matters here whether
  \textbf{2A} is framed like it is here, as a `flat' choice, or whether
  it is described as a compound lottery.} Crucially, many people who
take \textbf{1A} take it. Why take \textbf{2B}? Well, you will almost
certainly get nothing either way, so the extra 1\% of risk looks cheap
for the chance at five times the money.

Those two choices are inconsistent with expected utility theory. To see
why, imagine the outcomes are settled by drawing one of a hundred
numbered tickets, as shown in Table~\ref{tbl-allais}.

\begin{longtable}[]{@{}llll@{}}
\caption{The Allais choices, laid out over a hundred
tickets}\label{tbl-allais}\tabularnewline
\toprule\noalign{}
& \textbf{Ticket 1} & \textbf{Tickets 2-11} & \textbf{Tickets 12-100} \\
\midrule\noalign{}
\endfirsthead
\toprule\noalign{}
& \textbf{Ticket 1} & \textbf{Tickets 2-11} & \textbf{Tickets 12-100} \\
\midrule\noalign{}
\endhead
\bottomrule\noalign{}
\endlastfoot
\textbf{1A} & \$1M & \$1M & \$1M \\
\textbf{1B} & \$0 & \$5M & \$1M \\
\textbf{2A} & \$1M & \$1M & \$0 \\
\textbf{2B} & \$0 & \$5M & \$0 \\
\end{longtable}

Look at the last column. Within each pair, the two options pay the same
thing if a ticket from 12 to 100 comes up. So that column cannot bear on
either choice. Delete it, and what is left of \textbf{1A} against
\textbf{1B} is \emph{exactly the same} what is left of \textbf{2A}
against \textbf{2B}.

The algebra says the same thing. Write \(x\) for \(U(\$5\text{M})\),
\(y\) for \(U(\$1\text{M})\) and \(z\) for \(U(\$0)\). Preferring
\textbf{1A} to \textbf{1B} means

\begin{equation*}
y > 0.01z + 0.89y + 0.10x
\end{equation*}

which rearranges to

\begin{equation*}
0.11y - 0.10x - 0.01z > 0
\end{equation*}

And the difference between the two options in the second choice is

\begin{align*}
E(U(2A)) - E(U(2B)) &= (0.11y + 0.89z) - (0.10x + 0.90z) \\
&= 0.11y - 0.10x - 0.01z
\end{align*}

which is the same expression. So anyone who prefers \textbf{1A} to
\textbf{1B} is committed to preferring \textbf{2A} to \textbf{2B},
whatever their utilities for money happen to be. Taking \textbf{1A} and
\textbf{2B} is not a sign of unusual attitudes to risk. It is
inconsistent with expected utility theory. If you think \textbf{1A} and
\textbf{2B} are rational, you think that expected utility theory is
wrong.

\section{Ellsberg}\label{sec-ellsberg}

Around the same time, a similar challenge to expected utility theory was
formulated by
\href{https://en.wikipedia.org/wiki/Daniel_Ellsberg}{Daniel Ellsberg}.
Ellsberg is a fascinating figure; his Wikipedia page lists seven
\emph{films} that his story has a non-trivial part in. We're going to
discuss something he discovered quite early in life, before the work
that made him famous. How related it is to that later work is an
interesting question; Ellsberg himself always said it was.

Like Allais, Ellsberg (1961) discussed a pair of choices. Just as
before, Ellsberg showed that while an expected utility theorist could
explain either choice that the typical subject made, they could not
explain both. The choices concern an urn that has 90 marbles in it. The
subject is told that 30 marbles are red, and the other 60 are a mix of
black and yellow. But they are not told what the mixture is. They will
be offered a choice between two bets, which have the same return, say
\$10 if you win, and \$0 if you lose. All the bets turn on which marble
is drawn from the urn at its next draw.

Here is the first choice.

\begin{itemize}
\tightlist
\item
  \textbf{1A}: You win if the marble is red, and lose otherwise.
\item
  \textbf{1B}: You win if the marble is black, and lose otherwise.
\end{itemize}

Many people, I suspect most, take \textbf{1A}. They prefer knowing they
have a 1/3 chance to not knowing what their odds are. That's actually
fine; theory doesn't disagree.

Here is the second.

\begin{itemize}
\tightlist
\item
  \textbf{2A}: You win if the marble is red or yellow, and lose
  otherwise.
\item
  \textbf{2B}: You win if the marble is black or yellow, and lose
  otherwise.
\end{itemize}

Now most people prefer \textbf{2B}. They prefer knowing that they have a
2/3 chance of winning, to the unknown odds with \textbf{2A}. Ellsberg's
argument is that there is no way for standard expected utility theory to
explain both preferences.

At some level, the problem is the same as in the Allais paradox.
Table~\ref{tbl-ellsberg} is how the table looks if we draw the two of
them together.

\begin{longtable}[]{@{}llll@{}}
\caption{The Ellsberg choices, laid out
together}\label{tbl-ellsberg}\tabularnewline
\toprule\noalign{}
& \textbf{Red} & \textbf{Black} & \textbf{Yellow} \\
\midrule\noalign{}
\endfirsthead
\toprule\noalign{}
& \textbf{Red} & \textbf{Black} & \textbf{Yellow} \\
\midrule\noalign{}
\endhead
\bottomrule\noalign{}
\endlastfoot
\textbf{1A} & Win & Lose & Lose \\
\textbf{1B} & Lose & Win & Lose \\
\textbf{2A} & Win & Lose & Win \\
\textbf{2B} & \$0 & Win & Win \\
\end{longtable}

Here's what both choices have in common. In both choices, there's a
column where you get the same result whether you choose \textbf{A} or
\textbf{B}. That's the right hand column in each. The pair
\textbf{1A}/\textbf{2B} violates this principle.

\begin{tcolorbox}[enhanced jigsaw, arc=.35mm, bottomrule=.15mm, bottomtitle=1mm, breakable, colback=white, colbacktitle=quarto-callout-tip-color!10!white, colframe=quarto-callout-tip-color-frame, coltitle=black, left=2mm, leftrule=.75mm, opacityback=0, opacitybacktitle=0.6, rightrule=.15mm, title=\textcolor{quarto-callout-tip-color}{\faLightbulb}\hspace{0.5em}{Value of Irrelevant Alternatives}, titlerule=0mm, toprule=.15mm, toptitle=1mm]

Consider any choice between lotteries \textbf{A} and \textbf{B} where
there is some possibility where the two options return the same amount,
and you prefer \textbf{A} to \textbf{B}. Changing the lottery simply by
changing that common amount, without making any other changes, shouldn't
change your preferences. All that matters it that you get the same
amount in that outcome, not what that `same amount' is.

\end{tcolorbox}

This is a special case of what is known as the \textbf{Sure Thing
Principle}. Both of these are consequences of the expected utility
framework.\footnote{This arguably understates things. One prominent
  argument for the expected utility framework takes the Sure Thing
  Principle as a premise, and argues that anyone who satisfies it, and
  some other principles, is in effect an expected utility maximiser.
  We're not going to go into this, but if you're interested, look up
  work on \emph{representation theorems}.} And both the Allais and
Ellsberg paradoxes show that what I'm calling \textbf{Value of
Irrelevant Alternatives} is sometimes counterintuitive.

But while the two paradoxes target the same principle of expected
utility theory, they suggest different modifications to it. The Allais
paradox suggests that we are wrong to say that all attitudes to risk are
accounted for by declining marginal utility. The Ellsberg paradox
suggests that we are wrong to act as if everyone has a probability
function in their head. Think about how you'd do a version of the
algebraic argument that \textbf{1A}/\textbf{2B} is inconsistent in the
Ellsberg case.

Here's the most natural way to do this. By assumption, everyone has a
probability for every important possibility in their head. Assume the
chooser assigns probability \(y\) to the drawn marble being yellow. The
probability of \textbf{1A} winning is 1/3, and the probability of
\textbf{1B} winning is 2/3-\(y\), so if they prefer \textbf{1A} to
\textbf{1B}, it follows that \(y\) is greater than 1/3. The probability
of \textbf{2A} winning is 1/3+\(y\), and the probability of \textbf{2B}
winning is 2/3. So \(y\) is less than 1/3, contradicting the claim that
it is greater.

The response that Ellsberg and others gave to this argument is not to
question the algebra, or to question the appeal to expected utility.
It's to question the very assumption that there is such a thing as
\(y\). Or, in some versions, to question the assumption that it is a
\emph{number}. In Chapter~\ref{sec-bayesian}, we'll come back to some
game theoretic reasoning that also presupposes that things like
\emph{the probability that the marble is yellow} exist, and whether we
might want to question that assumption.

\section{For Next Time}\label{sec-ch2-next-time}

Next time we'll start with an application of those ideas that leads to
the most central concept in game theory.

\bookmarksetup{startatroot}

\chapter{Mixed Strategies and Nash
Equilibrium}\label{sec-mixed-and-nash}

\section{Mixed Strategies}\label{sec-mixed-strategies}

In any game, if player \(x\) has strategies \(s_1, s_2\) available, then
we'll also assume they have available the \textbf{mixed strategy} of
playing \(s_1\) with probability \(p\) and \(s_2\) with probability
\(1-p\) (for arbitrary \(p \in (0, 1)\)). Such strategies have an
expected value given strategies for the other players in the game, and
we'll think of that expected value as the return to the player for
playing the strategy.

Mixed strategies are contrasted with \textbf{pure strategies}. These are
the strategies that are usually listed on the game table. They are the
only strategies that we have considered so far.

I'm not sure whether there is a standard terminology for denoting mixed
strategies. We'll write \(\langle s_1, p; s_2, 1-p \rangle\) for the
mixed strategy of playing \(s_1\) with probability \(p\) and \(s_2\)
with probability \(1-p\). It is possible for mixed strategies to involve
more than two pure strategies; in those cases the list of pure
strategies and their probabilities will have to be longer.

If the agent has more than two pure strategies available to them, and
plays \(\langle s_1, p; s_2, 1-p \rangle\), it wouldn't make any
difference if they instead played
\(\langle s_1, p; s_2, 1-p; s_3, 0; \dots; s_n, 0\rangle\). We can treat
any mixed strategy as a mixture of all the pure strategies, just that
some of the strategies have probability 0 of being played. Indeed, we
can treat any strategy as a mixed strategy that perhaps gives some pure
strategies probability 0. So we no longer need to talk about pure
strategies at all; they are just a limit case of mixed strategies where
one pure strategy gets probability 1, and the others all get probability
zero.

Still, even if every mixed strategy can be represented as a mixture of
all pure strategies, there is an important distinction between the pure
strategies that get positive probability, and those that get zero
probability. We'll call the set of pure strategies that get positive
probability the \textbf{support} of the mixed strategy. What we used to
call a pure strategy is a mixed strategy whose support is a singleton.

Now it's a somewhat delicate philosophical question just what this
assumption amounts to. While the assumption is more or less universally
made by game theorists, there is less agreement about just what it comes
to. There are three natural interpretations, which we might call the
objective chance, epistemic and frequentist readings. We'll leave them
there for now; the question of which is right doesn't bite until we can
see what a mixed equilibrium asks of a player, so we'll come back to it
at the end of the chapter.

\section{Best Responses}\label{sec-best-responses}

A strategy \(s\) is a \textbf{best response} for player \(x\) iff there
is some probability distribution \(\Pr\) over the possible strategies of
other players such that no strategy has a higher expected payoff than
\(s\) does, given \(\Pr\).

Let's look at an example of this. So far we've used the same names for
the options that Row and Column have available. But we'll start to look
at asymmetric games, so it helps to use asymmetric terminology. In
Table~\ref{tbl-first-best} we'll use \(U\), for Up, \(M\) for Middle,
and \(D\), for Down, for Row's options, and \(L\) for Left, and \(R\)
for Right, for Column's options.

\begin{longtable}[]{@{}lrcc@{}}
\caption{An example of best
responses.}\label{tbl-first-best}\tabularnewline
\toprule\noalign{}
\endfirsthead
\endhead
\bottomrule\noalign{}
\endlastfoot
& & \multicolumn{2}{c@{}}{%
\textbf{Col}} \\
& & \(L\) & \(R\) \\
\multirow{3}{*}{\textbf{Row}} & \(U\) & 3, 0 & 0, 0 \\
& \(M\) & 2, 0 & 2, 0 \\
& \(D\) & 0, 0 & 3, 0 \\
\end{longtable}

We will just look at things from the perspective of Row. What we want to
show is that all three of the possible moves here are best responses. We
will write \(E(X) = x\) to mean that the expected value of strategy
\(X\) is \(x\).

It is clear that \(U\) is a best response. Set
\(\Pr(L) = 1, \Pr(R) = 0\). Then \(E(U) = 3, E(M) = 2, E(D) = 0\), so
\(U\) is best. It is also clear that \(D\) is a best response. Set
\(\Pr(L) = 0, \Pr(R) = 1\). Then \(E(U) = 0, E(M) = 2, E(D) = 3\), so
\(D\) is best.

The striking thing is that \(M\) can also be a best response. Set
\(\Pr(L) = \Pr(R) = 0.5\). Then \(E(U) = E(D) = 1.5\). But \(E(M) = 2\),
which is greater than 1.5. So if Row thinks it is equally likely that
Col will play either \(L\) or \(R\), then she maximizes expected utility
by playing \(M\). Of course, she doesn't maximize actual utility.
Maximizing actual utility requires making a gamble on which choice \(C\)
will make. That isn't always wise; it might be best to take the safe
option.

It's very plausible that agents should play best responses. If a move is
not a best response, then there is no way that it maximizes expected
utility, no matter what one's views of the other players are. And one
should maximize expected utility.

Any strongly dominated option will never be a best response. That's
because whatever strategy strongly dominates it will have a higher
expected return no matter what the probabilities are. But a weakly
dominated response can be a best response. That's because we've defined
best responses to allow that something can be a best response when
something else ties it. So in Table~\ref{tbl-wd-first}, \(D\) is weakly
dominated by \(U\), but it is a best response given \(\Pr(R) = 1\).

\begin{longtable}[]{@{}lrcc@{}}
\caption{An example of weak
dominance}\label{tbl-wd-first}\tabularnewline
\toprule\noalign{}
\endfirsthead
\endhead
\bottomrule\noalign{}
\endlastfoot
& & \multicolumn{2}{c@{}}{%
\textbf{Col}} \\
& & \(L\) & \(R\) \\
\multirow{2}{*}{\textbf{Row}} & \(U\) & 1, 0 & 0, 0 \\
& \(D\) & 0, 0 & 0, 0 \\
\end{longtable}

What about the other direction? If something is not dominated, does that
mean it is a best response? It turns out the answer to this question is
a little more complicated. Start with Table~\ref{tbl-bad-middle}.

\begin{longtable}[]{@{}lrcc@{}}
\caption{An example where something is dominated by a
mixture.}\label{tbl-bad-middle}\tabularnewline
\toprule\noalign{}
\endfirsthead
\endhead
\bottomrule\noalign{}
\endlastfoot
& & \multicolumn{2}{c@{}}{%
\textbf{Col}} \\
& & \(L\) & \(R\) \\
\multirow{3}{*}{\textbf{Row}} & \(U\) & 3, 0 & 0, 0 \\
& \(M\) & 1, 0 & 1, 0 \\
& \(D\) & 0, 0 & 3, 0 \\
\end{longtable}

In this game, \(M\) is not dominated by either \(U\) or \(D\). After
all, if Col plays \(L\), then \(M\) beats \(D\). And if Col plays \(R\)
then \(M\) beats \(D\).

On the other hand, \(M\) is not a best response. The possible views Row
could have about Col are all characterised by a single variable \(x\),
where \(x\) is the probability (according to Row) of Col playing \(L\).
This variable \(x\) can take any value in \([0, 1]\). But whatever value
it takes, \(M\) is not the best option. Note that the expected return of
\(L\) is \(3x\), and the expected return of \(R\) is \(3(1-x)\). From
that, we can easily compute that:

\begin{itemize}
\tightlist
\item
  If \(0 \leq x < 0.5\) then \(D\) is the best response.
\item
  If \(x = 0.5\) then \(U\) and \(D\) are tied for best response.
\item
  If \(0.5 < x \leq 1\) then \(U\) is the best response.
\end{itemize}

So in no circumstance is \(M\) best. It is not a best response.

But on second thought, perhaps it is dominated. Let's add a row to the
table for the mixed strategy \(\langle U, 0.5; D, 0.5 \rangle\). In
Table~\ref{tbl-bad-mix}, we'll add this to the table, and note its
expected value.

\begin{longtable}[]{@{}lrcc@{}}
\caption{A version of Table~\ref{tbl-bad-middle} with the mixed strategy
made explicit.}\label{tbl-bad-mix}\tabularnewline
\toprule\noalign{}
\endfirsthead
\endhead
\bottomrule\noalign{}
\endlastfoot
& & \multicolumn{2}{c@{}}{%
\textbf{Col}} \\
& & \(L\) & \(R\) \\
\multirow{4}{*}{\textbf{Row}} & \(U\) & 3, 0 & 0, 0 \\
& \(M\) & 1, 0 & 1, 0 \\
& \(D\) & 0, 0 & 3, 0 \\
& \(\langle U, 0.5; D, 0.5 \rangle\) & 1.5, 0 & 1.5, 0 \\
\end{longtable}

And note that \(\langle U, 0.5; D, 0.5 \rangle\) dominates \(M\);
whatever Col plays, that mixed strategy has a higher expected return
than \(M\).

If we change the payout to \(M\) so that it isn't dominated, it becomes
a best response. So recall the version of the game where the payout from
\(M\) from a guaranteed 1 to a guaranteed 2. I've reprinted that as
Table~\ref{tbl-second-best} with the mixed strategy made explicit.

\begin{longtable}[]{@{}lrcc@{}}
\caption{A version of Table~\ref{tbl-first-best} with the mixed strategy
added.}\label{tbl-second-best}\tabularnewline
\toprule\noalign{}
\endfirsthead
\endhead
\bottomrule\noalign{}
\endlastfoot
& & \multicolumn{2}{c@{}}{%
\textbf{Col}} \\
& & \(L\) & \(R\) \\
\multirow{4}{*}{\textbf{Row}} & \(U\) & 3, 0 & 0, 0 \\
& \(M\) & 2, 0 & 2, 0 \\
& \(D\) & 0, 0 & 3, 0 \\
& \(\langle U, 0.5; D, 0.5 \rangle\) & 1.5, 0 & 1.5, 0 \\
\end{longtable}

Now \(M\) is not dominated. Indeed, it dominates
\(\langle U, 0.5; D, 0.5 \rangle\). And as we saw above it is also a
best response. If Row thinks it is 50/50 whether Col plays \(L\) or
\(R\), then \(M\) has the highest expected return.

\section{Iterated Best Responses}\label{sec-iterated-best-responses}

It seems very plausible that one should only play best responses. But
not all best responses are worth playing. Some best responses are pretty
hard to justify when playing against a rational opponent. Consider
Table~\ref{tbl-ibr-first} from Row's perspective.

\begin{longtable}[]{@{}lrcc@{}}
\caption{A dominance example}\label{tbl-ibr-first}\tabularnewline
\toprule\noalign{}
\endfirsthead
\endhead
\bottomrule\noalign{}
\endlastfoot
& & \multicolumn{2}{c@{}}{%
\textbf{Col}} \\
& & \(L\) & \(R\) \\
\multirow{2}{*}{\textbf{Row}} & \(U\) & 5, 5 & 0, -5 \\
& \(D\) & 0, 5 & 2, -5 \\
\end{longtable}

In this game, \(D\) is a best response. It does best if Col chooses
\(R\). But why would Col do that? Col gets 5 for sure if she chooses
\(L\), and -5 for sure if she chooses \(R\). Assuming the weakest
possible rationality constraint on Col, she won't choose a sure loss
over a sure gain. So given that, Row should choose \(U\).

Now we could have shown that with just considering domination; indeed
domination by pure strategies. Table~\ref{tbl-ibr-second} is a little
more complex.

\begin{longtable}[]{@{}lrccc@{}}
\caption{A mixed strategy dominance
example}\label{tbl-ibr-second}\tabularnewline
\toprule\noalign{}
\endfirsthead
\endhead
\bottomrule\noalign{}
\endlastfoot
& & \multicolumn{3}{c@{}}{%
\textbf{Col}} \\
& & \(L\) & \(C\) & \(R\) \\
\multirow{2}{*}{\textbf{Row}} & \(U\) & 1, 3 & 0, 1 & 1, 0 \\
& \(D\) & 0, 0 & 1, 1 & 0, 3 \\
\end{longtable}

In Table~\ref{tbl-ibr-second} neither player has a dominated move, as
long as we restrict ourselves to pure strategies. And, from Row's
perspective, thinking about best responses doesn't help either. \(U\) is
a best response if Col is going to play \(L\) or \(R\) with probability
at least 0.5, and \(D\) is best response if Col is going to play \(m\)
with probability at least 0.5.

But note that \(C\) is not a best response for Col. One of the outer
options has a higher expected value, whatever probability she gives to
Row playing \(U\). Now let's assume, when thinking from Row's
perspective, that Col will play a best response. That is, we're assuming
Col will play either \(L\) or \(R\). Given that, the best thing for Row
to do is to play \(U\).

More carefully, if Row knows that Col is rational, and if Row knows that
rational agents always play best responses, then Row has a compelling
reason to play \(U\). This suggests an important normative status. Call
a \textbf{BRBR} strategy one that is a \textbf{B}est \textbf{R}esponse
to a \textbf{B}est \textbf{R}esponse.

But we can go further than that. What should Col play, given what we
just concluded about Row? She should play \(L\). After all, she knows
that Row will play \(U\), and given that \(L\) has the best return.

The idea that players should play BRBR strategies isn't strong enough to
get us that result. After all, \(D\) is a best response, and \(R\) is a
best response to \(D\). So \(R\) is BRBR. What it isn't is a best
response to something that is a best response to a best response. Rather
than keep counting layers of best responses, let's collapse this into
one notion of being a best response to a best response to a best
response to etc., to infinity. Call any strategy which has this feature
an \textbf{IBR}, or \textbf{I}terated \textbf{B}est \textbf{R}esponse.

\section{Nash Equilibrium}\label{sec-nash-equilibrium}

In a two-player game, for a strategy \(s_0\) to be IBR, it must be the
best response to some strategy \(s_1\), which is the best response to
some strategy \(s_2\), which is the best response to some strategy
\(s_3\), etc. But we are assuming that there are only finitely many
strategy choices available. So how can we get such an infinite chain
going?

The answer, of course, is to repeat ourselves. As long as we get some
kind of loop, we can extend such a chain forever, by keeping on circling
around the loop. And the simplest loop will have just two steps in it.
So consider any pair of strategies \(\langle s_0, s_1 \rangle\) such
that \(s_0\) is the best response to \(s_1\), and \(s_1\) is the best
response to \(s_0\). In that case, each strategy will be BRBRI, since we
can run around that two-step `loop' forever, with each stop on the loop
being a strategy which is a best response to the strategy we previous
stopped at.

When a pair of strategies fit together nicely like this, we say they are
a \textbf{Nash Equilibrium}. More generally, in an \(n\)-player game, we
use the following definition.

\begin{description}
\tightlist
\item[NE]
Some strategies \(s_1, ..., s_n\) in an \(n\)-player game form a
\textbf{N}ash \textbf{E}quilibrium iff for each \(i\), \(s_i\) is a best
response to the strategies played by the other \(n-1\) players. That is,
iff player \(i\) cannot do better by playing any alternative strategy to
\(s_i\), given that the other players are playing what they actually do.
\end{description}

The `Nash' in Nash Equilibrium~is in honour of John Nash, who developed
the concept, and proved some striking mathematical results concerning it
(Nash 1951).

\section{Nash Equilibrium in Simultaneous Move
Games}\label{sec-nash-simultaneous}

Consider Table~\ref{tbl-no-pure-eqm}:

\begin{longtable}[]{@{}lrcc@{}}
\caption{A game with no pure strategy
equilibria}\label{tbl-no-pure-eqm}\tabularnewline
\toprule\noalign{}
\endfirsthead
\endhead
\bottomrule\noalign{}
\endlastfoot
& & \multicolumn{2}{c@{}}{%
\textbf{Col}} \\
& & \(L\) & \(R\) \\
\multirow{2}{*}{\textbf{Row}} & \(U\) & 2, 0 & 0, 2 \\
& \(D\) & 0, 4 & 1, 0 \\
\end{longtable}

Any of the pure strategies is a best response. Indeed, any of them are
iterated best responses. So is all we can say about the game that
anything goes? Most game theorists disagree with that claim. They say
that once we consider mixed strategies, there is a distinctive solution
to the game. Consider the following pair of strategies:

\begin{itemize}
\tightlist
\item
  Row plays \(\langle U, \nicefrac{2}{3}; D, \nicefrac{1}{3} \rangle\)
\item
  Col plays \(\langle L, \nicefrac{1}{3}; R, \nicefrac{2}{3} \rangle\)
\end{itemize}

Given that Row has that strategy, let's work out the expected return (to
Col) of Col's two possible plays.

\begin{align*}
E(L) &= 0 \times \frac{2}{3} + 4 \times \frac{1}{3} \\
&= 0 + \frac{4}{3} \\
&= \frac{4}{3} \\
E(R) &= 2 \times \frac{2}{3} + 0 \times \frac{1}{3} \\
&= \frac{4}{3} + 0 \\
&= \frac{4}{3}
\end{align*}

And we can do the same calculation for the expected return (to Row) of
Row's possible strategies.

\begin{align*}
E(U) &= 2 \times \frac{1}{3} + 0 \times \frac{2}{3} \\
&= \frac{2}{3} + 0\\
&= \frac{2}{3} \\
E(D) &= 0 \times \frac{1}{3} + 1 \times \frac{2}{3} \\
&= 0 + \frac{2}{3} \\
&= \frac{2}{3}
\end{align*}

Both players would be, if they knew what the other player was doing,
indifferent between the two possible strategies that are available to
them.

The expected return of Col's playing the mixed strategy
\(\langle L, x; R, 1-x \rangle\) is \(xE(L) + (1-x)E(R)\). So if
\(E(L) = E(R) = \nicefrac{4}{3}\), then for all \(x \in [0, 1]\), the
expected return of \(\langle L, x; R, 1-x \rangle\) is also
\(\nicefrac{4}{3}\). So all strategies are equally good in response to
Row playing \(\langle U, \nicefrac{2}{3}; D, \nicefrac{1}{3} \rangle\).
That is, they are all best responses. By similar reasoning, all
strategies Row could play in response to Col playing
\(\langle L, \nicefrac{1}{3}; R, \nicefrac{2}{3} \rangle\) are equally
good, so they are all best responses.

So the two strategies
\(\langle U, \nicefrac{2}{3}; D, \nicefrac{1}{3} \rangle\) and
\(\langle L, \nicefrac{1}{3}; R, \nicefrac{2}{3} \rangle\) are best
responses to each other; they form a Nash Equilibrium. Indeed, this is
the only Nash Equilibrium of the game.

We aren't going to prove this, but for any game with bounded payouts,
and finite options for each player, there will be some Nash Equilibrium.
This equilibria will often be a mixed strategy equilibria, but it is
guaranteed to exist.

\section{Finding Nash Equilibria}\label{sec-finding-nash}

It isn't easy in general to find Nash Equilibria. But let's start with
some general approaches we can use.

\subsection{Iterated Dominance}\label{sec-iterated-dominance}

Strongly dominated strategies can't be part of an equilibrium. (I'll
just use `equilibrium' for `Nash Equilibrium' in this section, as is
common in game theory texts.) So strategies that only make sense if they
other player plays a dominated strategy can't be part of an equilibrium
either. This idea can help us find the equilibrium in some games, such
as Table~\ref{tbl-id-first}.

\begin{longtable}[]{@{}lrcc@{}}
\caption{A game that can be solved with
dominance}\label{tbl-id-first}\tabularnewline
\toprule\noalign{}
\endfirsthead
\endhead
\bottomrule\noalign{}
\endlastfoot
& & \multicolumn{2}{c@{}}{%
\textbf{Col}} \\
& & \(L\) & \(R\) \\
\multirow{2}{*}{\textbf{Row}} & \(U\) & 1, 2 & 5, 0 \\
& \(D\) & 0, 0 & 3, 6 \\
\end{longtable}

Neither of Col's strategies is dominated. Col wants to choose whatever
Row chooses. But look at the game from Row's perspective. Here there is
a strongly dominating strategy, namely choosing \(U\). So \(D\) cannot
be part of any equilbria. Indeed, if we assume that Row is rational, and
so won't play strongly dominated strategies, it is really as if the
strategy isn't even there. From that perspective, we can rewrite
Table~\ref{tbl-id-first} as Table~\ref{tbl-id-second}.

\begin{longtable}[]{@{}lrcc@{}}
\caption{The result of deleting a dominated strategy from
Table~\ref{tbl-id-first}}\label{tbl-id-second}\tabularnewline
\toprule\noalign{}
\endfirsthead
\endhead
\bottomrule\noalign{}
\endlastfoot
& & \multicolumn{2}{c@{}}{%
\textbf{Col}} \\
& & \(L\) & \(R\) \\
\textbf{Row} & \(U\) & 1, 2 & 5, 0 \\
\end{longtable}

Now in this game there is a dominating strategy for Col, namely \(L\).
So the only equilibrium of the game is \(\langle U, L \rangle\). And it
seems that those are the only rational strategies for the players,
assuming they each know the other to be rational.

We can use the idea of removing dominating strategies to illustrate some
puzzling features of a couple of other games. I won't do tables for
these games, because they'd be much too big. The first is a location
game that has many applications.

\begin{quote}
Two trucks have to choose where they will sell ice-cream on a particular
beach. There are 11 locations to choose from, which we'll number 0, 1,
\ldots, 9, 10. Spot 0 is at the left end of the beach, Spot 10 is at the
right end of the beach, and the other spots are equally spaced in
between. There are 10 people at each location. Each of them will buy
ice-cream. If one truck is closer, they will buy ice-cream from that
truck. If two trucks are equally close, then 5 of them will buy
ice-cream from one truck, and 5 from the other. Each truck aims to
maximise the amount of ice-cream it sells. Where should the trucks end
up?
\end{quote}

Let's start by looking, in Table~\ref{tbl-ic-tl}, at a fragment of the
payoff matrix. The payoffs are numbers of ice-creams sold. We'll call
the players Row and Col, as usual, and just use the number \(n\) for the
strategy of choosing location \(n\).

\begin{longtable}[]{@{}lrcccccc@{}}
\caption{The top-left corner of the ice cream
game}\label{tbl-ic-tl}\tabularnewline
\toprule\noalign{}
\endfirsthead
\endhead
\bottomrule\noalign{}
\endlastfoot
& & \multicolumn{6}{c@{}}{%
\textbf{Col}} \\
& & 0 & 1 & 2 & 3 & 4 & \ldots{} \\
\multirow{2}{*}{\textbf{Row}} & 0 & 55, 55 & 10, 100 & 15, 95 & 20, 90 &
25, 85 & \ldots{} \\
& 1 & 100, 10 & 55, 55 & 20, 90 & 25, 95 & 30, 80 & \ldots{} \\
\ldots{} & & & & & & & \\
\end{longtable}

I'll leave it as an exercise to confirm that these numbers are indeed
correct. Note that no matter what Col selects, Row is better off picking
1 than 0. If Col picks 0 or 1, she is a lot better off; she sells 45
more ice-creams. And if Col picks a higher number, she is a bit better
off; she sells 5 more ice-creams. So picking 1 dominates picking 0.

Let's look, in Table~\ref{tbl-ic-br} at the opposite corner of the
matrix.

\begin{longtable}[]{@{}lrcccccc@{}}
\caption{The bottom-right corner of the ice cream
game}\label{tbl-ic-br}\tabularnewline
\toprule\noalign{}
\endfirsthead
\endhead
\bottomrule\noalign{}
\endlastfoot
& & \multicolumn{6}{c@{}}{%
\textbf{Col}} \\
& & \ldots{} & 6 & 7 & 8 & 9 & 10 \\
\ldots{} & & & & & & & \\
\multirow{2}{*}{\textbf{Row}} & 9 & \ldots{} & 30, 80 & 25, 85 & 20, 90
& 55, 55 & 100, 10 \\
& 10 & \ldots{} & 25, 85 & 20, 80 & 15, 95 & 10, 100 & 55, 55 \\
\end{longtable}

Again, there is be a pattern. No matter what Col does, Row is better off
picking 9 than 10. In most cases, this leads to selling 5 more
ice-creams. If Col also picks 9 or 10, then picking 9 gives Row a big
advantage. The argument here was obviously symmetric to the argument
about picking 0 or picking 1, so I'll stop concentrating on what happens
when both players select high numbers, and focus from here on the low
numbers.

So there is a clear conclusion to be drawn from what we've said so far.

\begin{itemize}
\tightlist
\item
  Spot 0 is dominated by Spot 1, and Spot 10 is dominated by Spot 9. So
  if Row is rational, she won't pick either of those spots. Since the
  game is symmetric, if Col is rational, she won't pick either of those
  spots either. Moreover, neither 0 nor 10 can be part of any
  equilibrium.
\end{itemize}

Now let's turn, in Table~\ref{tbl-ic-tl-1d}, to the comparison between
Spot 1 and Spot 2, after we conclude that Row won't play 0. Again, we'll
just look at a fragment of the matrix.

\begin{longtable}[]{@{}lrccccccc@{}}
\caption{The top-left corner of the ice cream game, after deleting 0 for
Row}\label{tbl-ic-tl-1d}\tabularnewline
\toprule\noalign{}
\endfirsthead
\endhead
\bottomrule\noalign{}
\endlastfoot
& & \multicolumn{7}{c@{}}{%
\textbf{Col}} \\
& & 0 & 1 & 2 & 3 & 4 & 5 & \ldots{} \\
\ldots{} & & & & & & & & \\
\multirow{2}{*}{\textbf{Row}} & 1 & 100, 10 & 55, 55 & 20, 90 & 25, 95 &
30, 80 & 35, 75 & \ldots{} \\
& 2 & 95, 15 & 90, 20 & 55, 55 & 30, 80 & 35, 75 & 40, 70 & \ldots{} \\
\ldots{} & & & & & & & & \\
\end{longtable}

Again, there's a pattern. If Col selects any number above 2, then Row
sells 5 more ice-creams by picking 2 rather than 1. If Col selects
either 1 or 2, then Row sells 35 more ice-creams by picking 2 rather
than 1. But if Col selects 0, then Row sells 5 \emph{fewer} ice-creams
by picking 2 rather than 1. So picking 2 does \emph{not} dominate
picking 1.

But note the only circumstance when picking 2 is worse than picking 1 is
if Col picks 0. And picking 0 is, for Col, a strongly dominated
strategy. So picking 2 is sure to do better than picking 1 if Col does
not play a strongly dominated strategy. So let's delete that dominated
strategy. Table~\ref{tbl-ic-tl-2d} is what the top left corner of the
game matrix looks like when we do that.

\begin{longtable}[]{@{}lrcccccc@{}}
\caption{The top-left corner of the ice cream game, after deleting 0 for
both players}\label{tbl-ic-tl-2d}\tabularnewline
\toprule\noalign{}
\endfirsthead
\endhead
\bottomrule\noalign{}
\endlastfoot
& & \multicolumn{6}{c@{}}{%
\textbf{Col}} \\
& & 1 & 2 & 3 & 4 & 5 & \ldots{} \\
\multirow{2}{*}{\textbf{Row}} & 1 & 55, 55 & 20, 90 & 25, 95 & 30, 80 &
35, 75 & \ldots{} \\
& 2 & 90, 20 & 55, 55 & 30, 80 & 35, 75 & 40, 70 & \ldots{} \\
\ldots{} & & & & & & & \\
\end{longtable}

And now it looks like picking 1 is dominated. This teaches us another
lesson about the game.

\begin{itemize}
\tightlist
\item
  If we `delete' the option of picking either 0 or 10 for Col, then
  picking 2 dominates picking 1 for Row. In other words, if Row knows
  Col is rational, and hence won't pick a dominated option, then picking
  2 dominates picking 1 for Row, relative to the space of epistemically
  possible moves in the game. For similar reasons, if Row knows Col is
  rational, then picking 8 dominates picking 9. And if Col knows Row is
  rational, then picking either 1 or 9 is dominated (by 2 and 8
  respectively). And for similar reasons, neither 1 nor 9 can be part of
  an equilibrium strategy.
\end{itemize}

Summarising what we know so far,

\begin{itemize}
\tightlist
\item
  If the players are rational, they won't pick 0 or 10.
\item
  If the players know the other player is rational, they also won't pick
  1 or 9.
\end{itemize}

Let's continue down the matrix, now with Table~\ref{tbl-ic-tl-3d}. For
simplicity, we'll leave off columns 0 and 10, since they are dominated,
and we have deleted those as possible options.

\begin{longtable}[]{@{}lrcccccc@{}}
\caption{The top-left corner after deleting 0 and 1 for Row, and 0 for
Column}\label{tbl-ic-tl-3d}\tabularnewline
\toprule\noalign{}
\endfirsthead
\endhead
\bottomrule\noalign{}
\endlastfoot
& & \multicolumn{6}{c@{}}{%
\textbf{Col}} \\
& & 1 & 2 & 3 & 4 & 5 & \ldots{} \\
\multirow{2}{*}{\textbf{Row}} & 2 & 90, 20 & 55, 55 & 30, 80 & 35, 75 &
40, 70 & \ldots{} \\
& 3 & 85, 25 & 80, 30 & 55, 55 & 40, 70 & 45, 65 & \ldots{} \\
\ldots{} & & & & & & & \\
\end{longtable}

Picking 3 doesn't \emph{quite} dominate picking 2. In most
circumstances, Row does better by picking 3 rather than 2. But she does
a little worse if Col picks 1. But wait! We just had an argument that
Col shouldn't pick 1. Or, at least, if Col knows that Row is rational,
she shouldn't pick 1. Let's assume that Col does know that Row is
rational, and Row in turn knows that fact, so she can use it in
reasoning. That means she knows Col won't pick 1. So she can delete it
from consideration too. And once she does, picking 3 dominates picking
2, relative to the reduced space of epistemically possible outcomes.

I haven't belaboured the point as much as in the previous paragraphs,
but hopefully the following conclusion is clear.

\begin{itemize}
\tightlist
\item
  If the players know that the players know that the other player is
  rational, they won't pick 2 or 8. So neither 2 nor 8 can be part of an
  equilibrium.
\end{itemize}

And you can see where this is going. Once we rule out each player
picking 2 or 8, then picking 4 dominates picking 3, so picking 3 should
be ruled out. And picking 7 should be ruled out for symmetric reasons.
But once 3 and 7 are ruled out, picking 5 dominates picking either 4 or
6. So both players will end up picking 5.

And that is, in the standard economic textbooks, a nice explanation of
why we see so much `clustering' of shops. For instance, starting with
work by Harold Hotelling (1929), this has been used to explain why there
are several shops for one kind of good in one neighborhood, rather than
spread over town. Of course the full explanation of clustering is more
complicated, but it is nice to see such a simple model deliver such a
strong outcome.

This process is called the \textbf{Iterative Deletion of Dominated
Strategies}. More precisely, what we've used here is the strategy of
iteratively deleting \textbf{strongly} dominated strategies. This is a
powerful technique for solving games that don't, at first glance, admit
of any easy solution. If only one strategy pair remains after using this
technique, that is the only equilibrium of the game.

We'll return a bit to this point in what follows, but it's worth noting
that in a one-shot version of this game, to get to the conclusion that
each player should play 5, we had to use a number of very strong
assumptions. We had to assume that each player is rational, that they
know the other is rational, that they know that the other knows they are
rational, and so on through six iterations. In principle, to justify
iteratively deleting dominated strategies, we have to assume that there
are arbitrarily many iterations of this idea that each player knows that
the other player knows etc that they are rational.

\subsection{Finding Best Responses}\label{sec-finding-best-responses}

It can be useful to mark explicitly what the best responses are to each
possible play by the other player. In a two-player game, that can
sometimes lead to turning up Nash equilibria. Even when it does not, it
can give you a hint about what the support for a mixed strategy
equilibria could be. To see an example of this, consider
Table~\ref{tbl-complicated}.

\begin{longtable}[]{@{}lrccccc@{}}
\caption{A complicated game}\label{tbl-complicated}\tabularnewline
\toprule\noalign{}
\endfirsthead
\endhead
\bottomrule\noalign{}
\endlastfoot
& & \multicolumn{5}{c@{}}{%
\textbf{Col}} \\
& & \(a\) & \(b\) & \(c\) & \(d\) & \(e\) \\
\multirow{5}{*}{\textbf{Row}} & \(A\) & 9,9 & 5,1 & 6,4 & 3,3 & 1,8 \\
& \(B\) & 3,1 & 7,3 & 1,7 & 3,9 & 8,2 \\
& \(C\) & 4,5 & 8,7 & 2,4 & 4,1 & 2,9 \\
& \(D\) & 3,2 & 1,5 & 3,7 & 7,7 & 9,1 \\
& \(E\) & 5,3 & 8,3 & 7,6 & 3,5 & 8,4 \\
\end{longtable}

We'll start by finding the best response for Row to each of Col's plays.
We'll bold the payoff for Row that does best, assuming Col makes that
choice. Table~\ref{tbl-complicated-bold} shows what it looks like after
we do that.

\begin{longtable}[]{@{}lrccccc@{}}
\caption{A version of Table~\ref{tbl-complicated} with Row's best
responses bolded.}\label{tbl-complicated-bold}\tabularnewline
\toprule\noalign{}
\endfirsthead
\endhead
\bottomrule\noalign{}
\endlastfoot
& & \multicolumn{5}{c@{}}{%
\textbf{Col}} \\
& & \(a\) & \(b\) & \(c\) & \(d\) & \(e\) \\
\multirow{5}{*}{\textbf{Row}} & \(A\) & \textbf{9},9 & 5,1 & 6,4 & 3,3 &
1,8 \\
& \(B\) & 3,1 & 7,3 & 1,7 & 3,9 & 8,2 \\
& \(C\) & 4,5 & \textbf{8},7 & 2,4 & 4,1 & 2,9 \\
& \(D\) & 3,2 & 1,5 & 3,7 & \textbf{7},7 & \textbf{9},1 \\
& \(E\) & 5,3 & \textbf{8},3 & \textbf{7},6 & 3,5 & 8,4 \\
\end{longtable}

Note that under \(b\), we bolded two possible payoffs. That's because
they are tied for best, and best response is compatible with ties. In
Table~\ref{tbl-complicated-both-bold} we also bold Col's best responses.

\begin{longtable}[]{@{}lrccccc@{}}
\caption{A version of Table~\ref{tbl-complicated} with best responses
for both players
bolded.}\label{tbl-complicated-both-bold}\tabularnewline
\toprule\noalign{}
\endfirsthead
\endhead
\bottomrule\noalign{}
\endlastfoot
& & \multicolumn{5}{c@{}}{%
\textbf{Col}} \\
& & \(a\) & \(b\) & \(c\) & \(d\) & \(e\) \\
\multirow{5}{*}{\textbf{Row}} & \(A\) & \textbf{9},\textbf{9} & 5,1 &
6,4 & 3,3 & 1,8 \\
& \(B\) & 3,1 & 7,3 & 1,7 & 3,\textbf{9} & 8,2 \\
& \(C\) & 4,5 & \textbf{8},7 & 2,4 & 4,1 & 2,\textbf{9} \\
& \(D\) & 3,2 & 1,5 & 3,\textbf{7} & \textbf{7},\textbf{7} &
\textbf{9},1 \\
& \(E\) & 5,3 & \textbf{8},3 & \textbf{7},\textbf{6} & 3,5 & 8,4 \\
\end{longtable}

Any cell where both values are bolded is an equilibrium; neither player
can improve their lot by unilaterally changing their play. And we can
see there are three pure strategy equilibria here:
\(\langle Aa \rangle\), \(\langle Dd \rangle\), \(\langle Ec \rangle\).
I'll leave finding the mixed strategy equilibria as an exercise for the
interested reader.

\subsection{Finding Mixed Strategy Equilibria}\label{sec-finding-mixed}

It isn't easy in general to find mixed strategy equilibria. But we can
start by dividing the problem into two parts.

\begin{enumerate}
\def\labelenumi{\arabic{enumi}.}
\tightlist
\item
  Finding the support for each player's part of the equilibria.
\item
  Finding the probability with which they play each part of their
  support.
\end{enumerate}

And it turns out, a bit surprisingly, that it is the first part of the
problem that is harder in general. (Recall that the support of a mixed
strategy is the set of pure strategies that have a positive probability
of being played given that mixed strategy.) And while that isn't an easy
problem to solve in general, it isn't something we need to worry about
in very simple games. In games with just two options per player, then
either there will be a pure strategy equilibria, or the mixed strategy
Nash Equilibria will have all strategies in its support. Even if there
are three strategies available per player, then it usually isn't too
hard to find the support for the mixed strategy equilibria by just
noting that strongly dominated strategies cannot be part of an
equilibria.

We'll start by considering the very abstractly defined game, as in
Table~\ref{tbl-abstract-no-pure}.

\begin{longtable}[]{@{}lrcc@{}}
\caption{An abstract game with no pure strategy
equilibrium.}\label{tbl-abstract-no-pure}\tabularnewline
\toprule\noalign{}
\endfirsthead
\endhead
\bottomrule\noalign{}
\endlastfoot
& & \multicolumn{2}{c@{}}{%
\textbf{Col}} \\
& & \(L\) & \(R\) \\
\multirow{2}{*}{\textbf{Row}} & \(U\) & \textbf{a},b & c,\textbf{d} \\
& \(D\) & e,\textbf{f} & \textbf{0},0 \\
\end{longtable}

Don't make too much of the fact that I've put 0 for both players in the
bottom right. Since utility scales are only defined up to affine
transformation, adding a constant to each possible payout a player gets
literally doesn't change the game. So any game representation can be
changed into an equivalent representation with 0 for each player in an
arbitrarily chosen cell. We'll make that cell be the bottom right, but
it could be anywhere.

But do notice the bolding. That has meaning. In particular, it means
that all of the following must hold.

\begin{itemize}
\tightlist
\item
  \(a > e\), since \(a\) is bolded and \(e\) is not, and those are the
  choices if Col chooses \(L\).
\item
  \(0 > c\), since \(0\) is bolded (for Row) and \(c\) is not, and those
  are the choices if Col chooses \(R\).
\item
  \(d > b\), since \(d\) is bolded and \(b\) is not, and those are the
  choices if Row chooses \(U\).
\item
  \(f > 0\), since \(f\) is bolded and \(0\) is not (for Col), and those
  are the choices if Row chooses \(D\).
\end{itemize}

Since there is no cell where both are bolded, there is no pure strategy
equilibrium. So the equilibria must include at least one mixed strategy.

Recall how we calculated the expected return \(E(M)\) of the mixed
strategy \(\langle x, L; 1-x, R\rangle\), given the expected returns of
\(L\) and \(R\). It is

\begin{equation*}
E(M) = xE(L) + (1-x)E(R)
\end{equation*}

If we let \(k = E(R) - E(L)\), we can rewrite this as

\begin{align*}
E(M) &= xE(L) + (1-x)E(R) \\
&= xE(L) + (1-x)(E(L) + k) \\
&= xE(L) + (1-x)E(L) + k(1-x) \\
&= (x + (1-x))E(L) + k(1-x) \\
&= E(L) + k(1-x)
\end{align*}

A similar bit of reasoning shows that \(E(M) = E(R) + kx\). Now if the
strategy is genuinely mixed, then both \(x\) and \(1-x\) will be greater
than 0. And so if \(k \neq 0\), then one of \(kx\) and \(k(1-x)\) will
be positive. And if one of them is positive, then \(E(M)\) will be less
than one of \(E(L)\) and \(E(R)\), so it won't be a best response.

So the only time when \(\langle x, L; 1-x, R\rangle\) is a best response
for Col is when \(E(L) = E(R)\). In that circumstance, all responses are
best responses. In any other circumstance, one of the pure responses is
best. So,

\begin{itemize}
\tightlist
\item
  If the equilibrium has Col playing a mixed strategy, then Row must
  play a strategy that makes Col's expected returns from \(L\) and \(R\)
  equal.
\item
  And that can only obtain if Row herself plays a mixed strategy.
\end{itemize}

I've done all this from Col's perspective, but we could redo the
argument from Row's perspective to get a similar result; if a mixed
strategy is to be best for Row, then Col must play the mixed strategy
that makes each of Row's pure strategies equally good.

So now we have a recipe, or at least something like one, for finding the
mixed strategy equilibria. Each player must choose a mixed strategy that
makes the expected return of the other player's pure strategies equal.
So let's assume that

\begin{itemize}
\tightlist
\item
  Row plays \(\langle U, x; D, 1-x \rangle\); and
\item
  Col plays \(\langle L, y; R, 1-y \rangle\)
\end{itemize}

Then we need \(x\) to be such that \(E(L) = E(R)\), and \(y\) to be such
that \(E(U) = E(D)\). So we can solve for each of these variables in
turn.

\begin{align*}
E(L) &= E(R) \\
xb + (1-x)f &= xd \\
xb + f - xf &= xd \\
f &= xd + xf - xb \\
f &= x(d + f -b) \\
\frac{f}{d + f - b} &= x
\end{align*}

So Row must play the mixed strategy
\(\langle U, \frac{f}{d + f - b}; D, \frac{d-b}{d + f - b} \rangle\) if
there is to be an equilibrium. If she does, then Col's expected return
from playing either \(L\), or \(R\), or any mixture of them, will be
\(\frac{df - bf}{d + f - b}\).

Let's now solve for \(y\) using the same technique.

\begin{align*}
E(U) &= E(D) \\
ya + (1-y)c &= ye \\
ya + c - yc &= ye \\
c &= ye + yc - ya \\
c &= y(e + c - a) \\
\frac{c}{e + c - a} &= y
\end{align*}

So Col must play the mixed strategy
\(\langle L, \frac{c}{e + c - a}; R, \frac{e-a}{e + c - a} \rangle\) if
there is to be an equilibrium. If she does, then Row's expected return
from playing either \(U\), or \(D\), or any mixture of them, will be
\(\frac{ec}{e + c - a}\).

There is something puzzling about this result. Imagine a version of the
game where Col is a driver, and Row a police officer. Col doesn't care
much about the law, or the safety of other users of the road. Col just
wants to get where they are going as fast as possible, but really
doesn't want a speeding ticket. Row wants to catch speeding motorists,
but it is costly to engage in detection work. Col has a choice, Speed or
Not Speed. Row has a different choice, Detect or Not Detect. The payoffs
are in Table~\ref{tbl-speeding}.

\begin{longtable}[]{@{}lrcc@{}}
\caption{The speeding game}\label{tbl-speeding}\tabularnewline
\toprule\noalign{}
\endfirsthead
\endhead
\bottomrule\noalign{}
\endlastfoot
& & \multicolumn{2}{c@{}}{%
\textbf{Col}} \\
& & \(S\) & \(NS\) \\
\multirow{2}{*}{\textbf{Row}} & \(D\) & 10, -10 & -1, 0 \\
& \(ND\) & 0, 1 & 0, 0 \\
\end{longtable}

The equilibrium of this game is that Col speeds with probability
\(\frac{1}{11}\), and Row detects with probability \(\frac{1}{11}\), and
each has an expected payoff of 0.

Now what happens if the legislature decides that there is too much
speeding, and doubles the fine for speeding. Now the payoff table looks
like Table~\ref{tbl-speeding-adj}.

\begin{longtable}[]{@{}lrcc@{}}
\caption{The speeding game, with raised
fines}\label{tbl-speeding-adj}\tabularnewline
\toprule\noalign{}
\endfirsthead
\endhead
\bottomrule\noalign{}
\endlastfoot
& & \multicolumn{2}{c@{}}{%
\textbf{Col}} \\
& & \(S\) & \(NS\) \\
\multirow{2}{*}{\textbf{Row}} & \(D\) & 10, -20 & -1, 0 \\
& \(ND\) & 0, 1 & 0, 0 \\
\end{longtable}

If nothing else apart from the penalty for speeding, the rational
response by police department Row is to reduce monitoring of speeding,
and spend time on other crimes. After all, the new equilibrium is that
Row detects with probability \(\frac{1}{21}\), and Col still speeds with
probability \(\frac{1}{11}\). If there is a game in a mixed strategy
equilibrium, and the result of any change to the payoffs will be to move
it to a new equilibrium, then the only way to change Col's behavior is
to either create a new dominant strategy, or to change Row's payoffs.
Perhaps what the legislature should do here is to not just increase
speeding fines, but to give the money to police departments - that would
change the equilibrium of the game in a hurry.

\subsection{Combining Strategies}\label{sec-combining-strategies}

We'll work out the equilibria for Table~\ref{tbl-combining}.

\begin{longtable}[]{@{}lrccc@{}}
\caption{The example for
Section~\ref{sec-combining-strategies}}\label{tbl-combining}\tabularnewline
\toprule\noalign{}
\endfirsthead
\endhead
\bottomrule\noalign{}
\endlastfoot
& & \multicolumn{3}{c@{}}{%
\textbf{Col}} \\
& & \(a\) & \(b\) & \(c\) \\
\multirow{3}{*}{\textbf{Row}} & \(A\) & 5,0 & 0,1 & 1,-1 \\
& \(B\) & 0,2 & 1,0 & 2,1 \\
& \(C\) & 3,4 & -1,2 & 6,3 \\
\end{longtable}

As a first step, we'll bold the best responses, and see if any
equilibria turn up that way.

\begin{longtable}[]{@{}lrccc@{}}
\caption{A version of Table~\ref{tbl-combining} with best responses
bolded}\label{tbl-combining-bold}\tabularnewline
\toprule\noalign{}
\endfirsthead
\endhead
\bottomrule\noalign{}
\endlastfoot
& & \multicolumn{3}{c@{}}{%
\textbf{Col}} \\
& & \(a\) & \(b\) & \(c\) \\
\multirow{3}{*}{\textbf{Row}} & \(A\) & \textbf{5},0 & 0,\textbf{1} &
1,-1 \\
& \(B\) & 0,\textbf{2} & \textbf{1},0 & 2,1 \\
& \(C\) & 3,\textbf{4} & -1,2 & \textbf{6},3 \\
\end{longtable}

The bolding in Table~\ref{tbl-combining-bold} doesn't look like it
helps. Row doesn't have any dominated strategies either. But Col does
have a dominated strategy, since \(a\) dominates \(c\). It doesn't
matter that \(c\) might outperform \(b\); there is no reason to play
\(c\) since \(a\) will always do 1 point better. So let's delete that,
as in Table~\ref{tbl-combining-dom}.

\begin{longtable}[]{@{}lrcc@{}}
\caption{A version of Table~\ref{tbl-combining-bold} with a dominated
strategy deleted}\label{tbl-combining-dom}\tabularnewline
\toprule\noalign{}
\endfirsthead
\endhead
\bottomrule\noalign{}
\endlastfoot
& & \multicolumn{2}{c@{}}{%
\textbf{Col}} \\
& & \(a\) & \(b\) \\
\multirow{3}{*}{\textbf{Row}} & \(A\) & \textbf{5},0 & 0,\textbf{1} \\
& \(B\) & 0,\textbf{2} & \textbf{1},0 \\
& \(C\) & 3,\textbf{4} & -1,2 \\
\end{longtable}

We don't have to change the bolding, because a dominated strategy is
never a best response.

Now compare options \(A\) and \(C\). If we rule out Col playing \(c\),
it turns out there is never a reason for Row to play \(C\) over \(A\).
Option \(A\) does better than \(C\) whether Col plays \(a\) or \(b\). So
let's delete option \(C\) as well.

\begin{longtable}[]{@{}lrcc@{}}
\caption{A version of Table~\ref{tbl-combining-dom} with a strategy
deleted by iterated deletion}\label{tbl-combining-dom2}\tabularnewline
\toprule\noalign{}
\endfirsthead
\endhead
\bottomrule\noalign{}
\endlastfoot
& & \multicolumn{2}{c@{}}{%
\textbf{Col}} \\
& & \(a\) & \(b\) \\
\multirow{2}{*}{\textbf{Row}} & \(A\) & \textbf{5},0 & 0,\textbf{1} \\
& \(B\) & 0,\textbf{2} & \textbf{1},0 \\
\end{longtable}

Now we have a familiar kind of game. Each player has two options, and
there is no pure strategy equilibrium. The problem has reduced to the
form of the problem we tackled in the previous subsection. The only
equilibrium will be a mixed strategy equilibrium. And we find it by
making each player leave the other indifferent between their options. I
won't repeat the algebra for finding it, but you can verify yourself
that the following works.

\begin{itemize}
\tightlist
\item
  Row plays \(\langle A, \frac{2}{3}; B, \frac{1}{3} \rangle\).
\item
  Col plays \(\langle a, \frac{1}{6}; b, \frac{5}{6} \rangle\).
\end{itemize}

If you prefer to have the mixed strategies range over all of the options
that were originally available, with the strategies outside the support
getting probability 0, we can write the equilibrium of the game as.

\begin{itemize}
\tightlist
\item
  Row plays \(\langle A, \frac{2}{3}; B, \frac{1}{3}; C, 0 \rangle\).
\item
  Col plays \(\langle a, \frac{1}{6}; b, \frac{5}{6}; c, 0 \rangle\).
\end{itemize}

\section{Why Care about Equilibria?}\label{sec-why-care}

There is a philosophical question about why we should care about Nash
equilibria. Thinking about the mixed strategy equilibria can make these
questions particularly vivid.

Imagine I'm Col in one of the driving games from the last section. Why
should I ever play a mixed strategy? Here is an argument that I should
not.

\begin{itemize}
\tightlist
\item
  If police Row is playing the mixed strategy that brings the game into
  equilibrium, then my expected return from playing either Speed or Not
  Speed is 0. So I don't have any incentive to play the particular mixed
  strategy that brings the game into equilibrium.
\item
  Figuring out how to make it that I speed with probability
  \(\frac{1}{11}\) seems like hard work, when it will have no positive
  return.
\item
  If police Row is not playing the mixed strategy that brings the game
  into equilibrium, then one of the pure strategies open to me will have
  a higher expected return than playing the mixed strategy.
\end{itemize}

So either playing the mixed strategy will involve work for no gain, or
it will actually lower my expected return compared to an available
alternative. This isn't seeming very promising. But we can imagine some
circumstances when the mixed strategy would be more appealing. Here are
three related circumstances where it seems useful.

\begin{enumerate}
\def\labelenumi{\arabic{enumi}.}
\tightlist
\item
  Imagine that I'm too incompetent to make decisions for myself in games
  like this. You have to advise me. But the advice will be public;
  police Row will see it as well. Then the only advice you can
  rationally offer is that I should play the mixed strategy. Any other
  advice will backfire because the police will respond to it in ways
  that make it worthless.
\item
  Imagine that my strategy will leak to the police, no matter what. That
  could be because I'm so incompetent that I need to rely on public
  advice, or it could be because the police have a way of reading my
  moves. And imagine that the police are in the same boat; their
  strategy will leak to me. Then we should each play the equilibrium
  strategy; it is immune to doing worse in light of leaks. That's an
  implausible situation in most cases, but there is one special case
  worth considering.
\item
  Imagine the game will be repeated so many times that eventually our
  strategies will become visible enough just from the record of our
  play. Then we are (eventually) back into the `leaked strategy'
  situation. And in that situation, we should play equilibrium
  strategies.
\end{enumerate}

Personally, I think it's the last possibility that makes the concept of
an equilibrium play so important. Equilibria (in one shot games) are
interesting because they alone are sensible plays in repeated games.

This matches up nicely with our intuitive judgments about games like the
driving game. What would happen, intuitively, if one player deviated
dramatically from the equilibrium in a one-shot game? The answer is that
that player would probably do fine. If Col speeds, then they'll probably
get away with it. But think about what would happen (again intuitively)
if a player repeated non-equilibrium play over a long period of time. I
think it would end up being counterproductive. If Col speeds all the
time, Row will adjust and catch them. If Row detects all the time, Col
will slow down, but Row will end up wasting detecting resources. It's in
repeated games, that non-equilibrium play leads to disaster.

Recall that in Section~\ref{sec-mixed-strategies} we said that it is
hard to sometimes make sense of what it even means to play a mixed
strategy outside of the context of a repeated game. The reasoning of the
last paragraph is why I don't find that worrying. Just like it is hard
to define a mixed strategy in a genuinely one-shot game, it is hard to
motivate such a strategy. Mixed strategies become easier to define in
repeated games, and they become easier to motivate in those games too.

\section{What Is a Mixed Strategy?}\label{sec-what-is-a-mixed-strategy}

Back in Section~\ref{sec-mixed-strategies} we assumed that players have
mixed strategies available, and noted that it is a tricky question what
that assumption amounts to. We can now make that note much more precise.

This is the feature of mixed equilibria that generates the difficulty.
At a mixed equilibrium, every pure strategy in the support has the same
expected return as the mixture itself. We use that fact to discover
equilibria. The trick is to find the probabilities that make the
\emph{other} player indifferent. That means that at the equilibrium, a
player has no reason at all to play the mixture rather than any of the
strategies in its support. Playing \(U\) for certain does exactly as
well as playing \(\langle U, 0.5; D, 0.5 \rangle\), given what the other
player is doing.

For most pure strategy equilibria, abandoning the equilibrium makes you
strictly worse off. (There are exceptions, and these create their own
problems, but the normal case is that deviation has positive cost.) In a
mixed equilibrium nothing pushes you towards the equilibrium mixture.
The only thing your mixture accomplishes is keeping the other player
indifferent, which is an odd thing for a self-interested agent to be
aiming at.

So the usual story about equilibrium, that each player is doing the best
they can for themselves, does not explain why anyone would mix in the
proportions the equilibrium specifies. So what does explain it? The
different interpretations suggest different answers to this question.
We'll look at the three most important interpretations, though note that
there are others, and that each of these three admits some internal
variation.

\subsection{Objective Chance}\label{objective-chance}

On the first of our three interpretations, playing a mixed strategy is
using a randomising device, and the probabilities are objective chances.

This is the original reading. Neumann and Morgenstern (1944) introduced
mixtures precisely so that an opponent could not find out what you were
going to do, and they meant you're using a randomising device. This is
the natural reading for games where one side is trying to keep their
strategy hidden from the other. (Think about poker, or sports where
deception is important, or a tax authority deciding who to audit.) In
those cases, being predictable is what you most want to avoid.
Professional athletes do come close to the predicted mixtures; there is
a small empirical literature on tennis serves and penalty kicks.

There are a few of reasons to not like this interpretation. One is that
it's a wide open empirical question whether there really are any
randomising devices. The most common interpretations of modern physics
are indeterministic, but some are deterministic. Another problem is that
if there is indeterminism at the quantum level in physics, most people
we describe as playing a mixed strategy are not actually using devices
which are indeterministic in this quantum sense. And a third is that it
doesn't actually explain the puzzle we started with. Even if you have a
randomising device, and you have reason to use it, you need a further
reason to set it to the equilibrium probabilities. In equilibrium, you
don't have such a reason.

There is a weaker version of the same thought that avoids the first two
of these problems. What matters, one might say, is not that I actually
use a randomiser, but that you cannot tell the difference between what I
am doing and a chance process. On this version I might be running
something perfectly deterministic; if you cannot model it, then as far
as the game is concerned it is a mixture. Notice that this has quietly
relocated the probability from my machinery to your ignorance, which
makes it a step towards the epistemic reading rather than a version of
this one. And it still doesn't solve the explanatory problem. But it
least it doesn't make game theory sensitive to how quantum mechanics
ends up.

\subsection{Frequentist}\label{frequentist}

On the second interpretation the mixture is not in anyone's head. It is
a long-run frequency. This is similar to the \emph{frequentist}
interpretation of probability, most associated with Hans Reichenbach
(1949). Playing a mixture of 1/3 \(U\) and 2/3 \(D\) just means that if
the game is repeated enough times, you will average towards 1/3 of your
plays being \(U\), and 2/3 being \(D\). (There need to be some extra
clauses to, to prevent playing \(U\) exactly every third play. Just what
those constraints should be, and why they are justified inside the
frequentist framework, are quite delicate questions, but not ones we'll
deal with here.)

The immediate problem is that you might not be playing the game
repeatedly. We know from other games that when a game is repeated,
sometimes that changes the strategy which one should play. A completely
orthodox view is that one should defect in a one-shot Prisoners'
Dilemma, but cooperate (until the other party defects) in a repeated
Prisoners' Dilemma. The frequentist interpretation simply assumes that
what we say about repeated games applies to the one-shot game, and that
assumption is not, in general, true.

There is a variant that avoids the problem by moving from one player to
many. On the \textbf{population} reading, the mixture describes a
population rather than a person: some fraction of players are
\(U\)-types who always play \(U\), the rest are \(D\)-types, and nobody
randomises at all. The equilibrium becomes a fact about how the
population is composed, and the interesting question is which
compositions are stable. This is the entry point to evolutionary game
theory, which we are not doing in this course. Skyrms (1996) is where to
go if you want to follow it up.

There is also a variant which is a more radical departure from the
textbooks, but I think not uncommon if you talk to economists
informally. It's that the presupposition of the question is wrong: just
like the puzzle says, there is \emph{zero} reason to mix in a one-shot
game. On this view, the frequentist is right about the repeated game,
and there is nothing systematic to say about the one-shot game. (Perhaps
the rationalizability view, which we'll come back to in
Section~\ref{sec-rationalizability}, is right.) Let's set that aside,
and stick with the assumption that the orthodox game theoretic account
is right, and there is a reason to mix. This pushes us towards the last
of the interpretations.

\subsection{Epistemic}\label{epistemic}

The third interpretation gives up on the idea that anyone randomises. It
says the probabilities in mixed strategies are all about beliefs not
about actions.

Aumann and Brandenburger (1995) develop the strongest version of this
idea. Row plays some pure strategy. The mixture attributed to Row is
Col's \emph{belief} about what Row will do. Row's indifference is then
not a puzzle about Row's motivation at all; it is a fact about Col's
uncertainty.

This doesn't just change what a mixed strategy is. It changes what an
equilibrium is. On this view, an equilibrium is a series of beliefs
about the parties that have the following two features. First, they are
all true. Second, on the assumption that everyone shares these beliefs,
everyone is maximizing expected ytility. On this view, the theorem to
prove is that mutual knowledge of rationality and of the conjectures
about what people are believing forces those conjectures into the
pattern that Nash equilibrium describes. It also dissolves the
indifference puzzle rather than answering it: there was never a decision
to explain, because nobody was choosing a mixture.

To the extent that contemporary philosophers talk about rational choice
game theory (as distinct from evolutionary game theory), this seems to
be the most popular interpretation. It's certainly the interpretation
which minimizes the gap between game theory and decision theory. But
while this view is not without adherents in economics, indeed, Aumann
won a Nobel Prize in part for his work on it, it is far from orthodoxy
in economics.

\subsection{Where This Leaves Us}\label{where-this-leaves-us}

The three readings disagree about where the probability lives: in a
device, in a sequence of plays, or in someone's beliefs. They also
disagree about what an equilibrium is a claim about. The first makes it
a claim about individual behaviour, the second about a series or a
population, the third about states of information.

There is a fourth answer, due to Harsanyi, which says the randomness is
real but sits in the other player's ignorance of your exact payoffs. We
cannot state it properly yet, because it needs machinery we have not
built. We will come back to it in Chapter~\ref{sec-bayesian}.

You do not have to choose between these interpretations. You do have to
notice that ``the players mix 50-50'' is not yet a claim with a settled
meaning.

\bookmarksetup{startatroot}

\chapter{Other Equilibrium Concepts}\label{sec-other-equilibria}

The primary concept we'll use in this course is Nash Equilibrium. But
there are many many other solution concepts for games. Many of them have
practical applications. Many of them have a case to be the one true
constraint on rational agents playing games.

\section{Rationalizability}\label{sec-rationalizability}

We looked at this briefly in Section~\ref{sec-iterated-best-responses},
under the label of `Iterated Best Response'. But it's worth stating the
idea more carefully, and saying a bit about why it is an important
concept.

A class of strategies is rationalizable if:

\begin{itemize}
\tightlist
\item
  A class of strategies is rationalizable if every member of the class
  is a best response, assuming that the probability that every other
  player plays some strategy from the class is 1;
\item
  A strategy is rationalizable if it is a member of some class of
  rationalizable strategies.
\end{itemize}

This is strictly weaker than the notion of a Nash equilibrium. Note that
for any Nash equilibrium, the class consisting of just the strategies
that form the equilibrium will be rationalizable, so every member of the
class will be rationalizable. But some strategies that are not part of
any Nash equilibrium are rationalizable. Consider, for example,
Table~\ref{tbl-rat-no-pure} (which is similar to some we've looked at in
the previous chapter).

\begin{longtable}[]{@{}lrcc@{}}
\caption{A game with no pure strategy
equilibria}\label{tbl-rat-no-pure}\tabularnewline
\toprule\noalign{}
\endfirsthead
\endhead
\bottomrule\noalign{}
\endlastfoot
& & \multicolumn{2}{c@{}}{%
\textbf{Col}} \\
& & \(a\) & \(b\) \\
\multirow{2}{*}{\textbf{Row}} & \(A\) & 5,0 & 0, 1 \\
& \(B\) & 0,1 & 1, 0 \\
\end{longtable}

There is a unique Nash equilibrium for the game. Row plays
\(\langle A, 0.5; B, 0.5 \rangle\), and Col plays
\(\langle a, \frac{1}{6}; b, \frac{5}{6} \rangle\). But any strategy
whatsoever, including all mixtures, is rationalizable. For any strategy
you like, for either player, there is some opposing strategy it is a
best response to.

Now compare the following two claims about rationality and game playing.

\begin{enumerate}
\def\labelenumi{\arabic{enumi}.}
\tightlist
\item
  Rational players always play some part of a Nash equilibrium strategy.
\item
  Rational players always play rationalizable strategies.
\end{enumerate}

These two theories make different assessments of
Table~\ref{tbl-rat-no-pure}. Principle 1 says that rationality rules out
all but 1 strategy for each player. Principle 2 says anything goes.
Which of them get it right?

If the game was being repeated many times, the answer is easy. The pure
strategy \(A\) will lead to disaster for Row, since after a while, Col
will pick up what's going on and respond by playing \(b\). In a repeated
game, randomizing as Principle 1 requires is the right thing to do.

But in a one shot game it is harder to say what is correct. There is a
nice argument that Principle 1 is self-undermining. Assume that the only
rational thing to do is to play the Nash Equilibrium strategy. Then Row
will flip a coin to decide what to play. So Col has the same expected
return either way, given Row's strategy. Moreover, since Col knows what
rationality is, Col knows that Row will flip a coin. So now Col is just
facing a decision problem; not a game. And in that decision problem, all
responses have equal expected utility. So all responses are equally
rational. So Col can rationally do anything. But that's just what
Principle 2 says.

The argument of the last paragraph is perfectly general. In any game
where the only Nash equilibrium is a mixed strategy equilibrium, the
following principles cannot all be true.

A. Only strategies that are part of Nash equilibria are rational, and
all players know this. B. All players are rational, and all players know
this. C. Any strategy that is known to have the same expected return as
a rational strategy can be rationally played.

By (A) and (B), the other player will play a mixed strategy. But as
we've seen, this will mean that every strategy in the support for one's
own mixed strategy will have equal expected return. And by (C), that
means that the pure strategies are just as rational as the mixed
strategies. But that contradicts the assumption in (A) that only mixed
strategies are rational.

Conventional wisdom in game theory texts is that (C) fails. Principle 2
says that (A) fails. I'm not going to take a stand on this, just to note
that there is a clash between the idea that Nash equilibria have a
distinctive role to play in rational strategy choice, and the principle
that rational agents maximize expected utility.

In games where the players only have two choices to pick from, the
difference between the Principles comes down to whether it is obligatory
to randomize, or whether it is permissible to play one of the pure
strategies from the support of the mixed strategy Nash equilibrium. This
not a trivial question to even have intuitions about, because it is tied
up with the metaphysics of what a mixed strategy is. But in three way
choices there can be more separation between the two principles.
Consider either Table~\ref{tbl-rat-payout-1} or
Table~\ref{tbl-rat-payout-2}.

\begin{longtable}[]{@{}lrccc@{}}
\caption{A three-way game with a poor pure strategy
equilibrium}\label{tbl-rat-payout-1}\tabularnewline
\toprule\noalign{}
\endfirsthead
\endhead
\bottomrule\noalign{}
\endlastfoot
& & \multicolumn{3}{c@{}}{%
\textbf{Col}} \\
& & \(a\) & \(b\) & \(c\) \\
\multirow{3}{*}{\textbf{Row}} & \(A\) & 3,0 & 0,3 & 0,2 \\
& \(B\) & 0,3 & 3,0 & 0,2 \\
& \(C\) & 2,0 & 2,0 & 1,1 \\
\end{longtable}

\begin{longtable}[]{@{}lrccc@{}}
\caption{A version of Table~\ref{tbl-rat-payout-1} with the pure
strategy equilibrium raised to 2}\label{tbl-rat-payout-2}\tabularnewline
\toprule\noalign{}
\endfirsthead
\endhead
\bottomrule\noalign{}
\endlastfoot
& & \multicolumn{3}{c@{}}{%
\textbf{Col}} \\
& & \(a\) & \(b\) & \(c\) \\
\multirow{3}{*}{\textbf{Row}} & \(A\) & 3,0 & 0,3 & 0,2 \\
& \(B\) & 0,3 & 3,0 & 0,2 \\
& \(C\) & 2,0 & 2,0 & 2,2 \\
\end{longtable}

In both games, there is a unique Nash equilibrium, \(Cc\). But all
strategies are rationalizable. Note that \(A\) is a best response to
\(a\), which is a best response to \(B\), which is a best response to
\(b\), which is a best response to \(A\). Once you have a cycle like
this, every member of the cycle is rationalizable. But there is no mixed
strategy Nash equilibrium with those pure strategies as supports. For
reasons that should be familiar by now, the only mixed strategy
equilibrium that would be available is for each player to play a 50/50
split between their first two options. But if one player plays that
mixed strategy, the other player has a return of 1.5 if they
reciprocate, and a return of 2 if they play the third option. So there
is only one equilibrium, pure or mixed, for the game, while all options
are rationalizable.

Ideally, we would have clear intuitions here about what strategies are
rationally permissible, and this would let us decide between Principle 1
and Principle 2. Unfortunately, I don't think the intuitions about the
cases are at all clear. For what it is worth, I think intuitively that
playing one of \(A\) or \(B\) is rational when the alternative is the
equilibrium that returns 1 to each player, but not when the alternative
is a guaranteed 2. So these cases don't ultimately settle the question.
But they do make more vivid what is at stake. Could it ever be rational,
under conditions of common knowledge of rationality, to play \(A\) or
\(B\) in either of these games?

\section{Correlated Equilibrium}\label{sec-correlated-equilibrium}

To get Nash equilibrium, we have to assume that every player has a way
of letting their play be chosen by a random process. But we assume that
these random choices are independent. If Row has a 40\% chance of
playing Up, and Column has a 30\% chance of playing Left, then the
probability of the Up-Left outcome is 12\%. (That is, 40\% times 30\%.)
The idea behind correlated equilibrium is that the players may have a
way to coordinate on a random strategy.

Start with a game, Table~\ref{tbl-chicken-correlated} where we might
want such a coordination. Imagine that Row and Col are drivers,
approaching an intersection at right angles. Each of them has a choice
to go (written as G/g), or stop (written as S/s). They would each prefer
to get through the intersection faster. But the other person going and
them stopping isn't terrible; there is no crash and the intersection is
cleared. And both going is a disaster. So we can approximate the game
table as this. (In reality it's more like a many move game, but this is
a first approximation.)

\begin{longtable}[]{@{}lrcc@{}}
\caption{A version of
Chicken}\label{tbl-chicken-correlated}\tabularnewline
\toprule\noalign{}
\endfirsthead
\endhead
\bottomrule\noalign{}
\endlastfoot
& & \multicolumn{2}{c@{}}{%
\textbf{Col}} \\
& & \(g\) & \(s\) \\
\multirow{2}{*}{\textbf{Row}} & \(G\) & 0,0 & 10, 9 \\
& \(S\) & 9,10 & 8, 8 \\
\end{longtable}

Table~\ref{tbl-chicken-correlated} has two pure strategy equilibria,
\(Gs\) and \(Sg\). It also has a rather annoying mixed strategy
equilibria, where each player plays goes with probability
\(\nicefrac{2}{11}\), and stays otherwise, and they crash about 3\% of
the time.

Those are all the mixed strategy equilibria. Two of them are asymmetric,
and intuitively unfair, and it isn't easy to see how they would get
chosen. The only symmetric equilibrium is much worse than either pure
strategy equilibria; it even allows for the possibility of crashes.

The correlated equilibria are designed to avoid those problems. Imagine
that instead of choosing their own strategy at random, each player could
see a common randomising device. Maybe we could have a coin that was
flipped in public, and Row goes if it lands heads, and Col goes if it
lands tails. That would give each player an expected return of 9.5. And,
crucially, once each player sees the (common) signal, each player has no
incentive to defect. If Row knows that Col will go if the coin comes up
tails, and Row sees that the coin comes up tails, and that Col sees
this, Row will prefer to Stay than Go.

This is not completely unlike how we actually solve these problems in
real life. Traffic lights are ways of correlating our actions without
having one party always getting what they want. I'm simplifying over a
lot of details here, but the main thing to stress is that a traffic
light (a very public signal) is a way of introducing a new kind of
solution to the game drivers face.

In this case, the coin flip is the best symmetric strategy that we can
have. But that is to some extent a feature of the values. I won't try to
do a just so story to make the following table make sense, but imagine
that players faced Table~\ref{tbl-correlated-not-chicken}. (Perhaps it
comes about if they really dislike waiting at lights while the other
person goes through.)

\begin{longtable}[]{@{}lrcc@{}}
\caption{A traffic game that isn't quite
Chicken}\label{tbl-correlated-not-chicken}\tabularnewline
\toprule\noalign{}
\endfirsthead
\endhead
\bottomrule\noalign{}
\endlastfoot
& & \multicolumn{2}{c@{}}{%
\textbf{Col}} \\
& & \(g\) & \(s\) \\
\multirow{2}{*}{\textbf{Row}} & \(G\) & 0,0 & 10, 7 \\
& \(S\) & 7,10 & 9, 9 \\
\end{longtable}

The Nash equilibria don't change much with this change. And the
correlated equilibrium of the previous problem still works. That is, the
solution: flip a coin in public, if Heads, Row goes, if Tails, Col goes,
is still an equilibrium. Each player has an incentive to do their part
in the equilibrium once they see the coin. But there is a yet better
equilibrium available.

Imagine an outside judge makes three cards: \(Gs\); \(Sg\) and \(Ss\).
Each card has a strategy for Row (in Capital letters) and a strategy for
Col (in small letters). The judge selects a card, and then tells each
player the strategy that is recommended for them, but not what's
recommended for the other player. Each player still has an incentive to
follow the judge's recommendations. If the judge says go, they know the
other will stop, so they will want to go and get the best result. If the
judge says stop, it is 50/50 whether the judge has told the other player
go or stop. And in that circumstance, going is a bad gamble. It has an
expected return of 5; stopping has an expected return of 8. So both
parties will do what the judge says.

And if the parties agree in advance to this process, which is fair and
symmetric, their expected return is even higher than if they agree to
the coin flip. If they flip a coin to decide who will go, their expected
return is 8.5. Under this strategy, their expected return is
\(8\nicefrac{2}{3}\). That's because they have equal chances of getting
7, 9 or 10 as their return.

Let's take stock. A \textbf{correlated equilibrium} is a way of randomly
choosing instructions for each player to follow such that each player is
motivated to follow the instructions assuming (a) the probabilities of
different instructions being chosen are common knowledge, and (b) each
other player will follow the instructions. Every correlated equilibrium
will consist entirely of rationalizable strategies, but some will not be
Nash equilibria. Every Nash equilibria will be a correlated equilibrium.

Given a correlated equilibrium, we can work out the expected return to
each player. We do this, as always, by working out their probability of
getting each outcome, multiplying the probability by the utility of that
outcome, and adding up. We can also find the correlated equilibrium that
maximises expected returns, but that takes some more time. Let's look
again at Table~\ref{tbl-correlated-not-chicken}.

First, let's work out the mixed strategy equilibrium for the game. We
want to find a probability for Row to Go that makes Col indifferent
between going and staying. Let that probability be \(x\). Then we have
all the following lines being equivalent.

\begin{align*}
E(g) &= E(s) \\
10(1-x) &= 7x + 9(1-x) \\
10 - 10x &= 7x + 9 - 9x \\
10 - 10x &= 9 - 2x \\
10 &= 9 + 8x \\
1 &= 8x \\
x &= \frac{1}{8}
\end{align*}

So the mixed strategy equilibrium is that each player goes with
probability \(\nicefrac{1}{8}\). That leaves them in the crash outcome
in \(\frac{1}{64}\) of cases. Despite that, the strategies have a
surprisingly high expected return of 8.75.

Note one other consequence of the calculation we just did. If the
probability of the other player going is \(\nicefrac{1}{8}\), the other
player will be indifferent between their options. If it is higher than
\(\nicefrac{1}{8}\), the other player will maximise expected return by
stopping; if it is lower, the other player will maximise expected return
by going.

Now look at the simplest correlated equilibrium. A coin is flipped, and
if it is Heads, Row Goes, and Col stops, while if it is Tails, Row Stops
and Col goes. It is easy to confirm that each player will conform to
this instruction; if you know what the other player is doing, you do the
opposite. But this isn't a great equilibrium. Each player gets 10 with
probability 0.5, and 7 with probability 0.5, for an expected return of
8.5. That's worse even than the mixed strategy equilibrium.

If the judge selects one of the three cards \(Gs\); \(Sg\) and \(Ss\) at
random, then each player will follow the instructions. That's obvious if
they are told to go. If they are told to stop, the probability that the
other player was told to go is 0.5. Just to confirm that,

\begin{align*}
&Pr(\text{Other player told to go}|\text{I'm told to stop}) \\
= &\frac{Pr(\text{Other player told to go and I'm told to stop})}{Pr(\text{I'm told to stop})} \\
= &\frac{Pr(\text{Other player told to go and I'm told to stop})}{Pr(\text{They told go and I'm told stop}) + Pr(\text{They told go and I'm told stop})} \\
= &\frac{\frac{1}{3}}{\frac{1}{3} + \frac{1}{3}} \\
= &0.5
\end{align*}

And if the probability that the other player goes is greater than
\(\nicefrac{1}{8}\), a player will stop. But as we noticed above, the
expected return of this strategy is 8\(\nicefrac{2}{3}\), which is still
lower than the expected return of the mixed strategy. Can we do better?

It turns out that we can. The expected return of each player is just
half the sum of the expected returns of each player (in any symmetric
equilibrium). And the sum of the expected returns is maximised when we
maximise the probability of ending up in the bottom left cell, where the
players combine for 18 utils, compared to just 17 utils in the other two
non-crash cells. So we want to make the probability of ending up there
as high as possible. But we can't just tell the players to go to that
cell; they won't have an incentive to cooperate.

What we want to do is to maximise the probability of getting the \(Ss\)
instruction, consistent with each player being motivated to follow that
instruction. And that will happen as long as each player thinks the
probability the other player goes (conditional on them being told to
stop) is at least \(\nicefrac{1}{8}\). So let's assume that the judge
makes the \(Gs\) instruction with probability \(x\), the \(Sg\)
instruction with probability \(x\), and the \(Ss\) instruction with
probability \(1 - 2x\). Now do the above conditional probability
calculation with those values.

\begin{align*}
&Pr(\text{Other player told to go}|\text{I'm told to stop}) \\ 
= &\frac{Pr(\text{Other player told to go and I'm told to stop})}{Pr(\text{I'm told to stop})} \\
= &\frac{Pr(\text{Other player told to go and I'm told to stop})}{Pr(\text{They told go and I'm told stop}) + Pr(\text{They told go and I'm told stop})} \\
= &\frac{x}{x + (1-2x)}
\end{align*}

And we want that last value to be \(\nicefrac{1}{8}\). A little
calculation shows that this will be true if \(x = \nicefrac{1}{9}\). So
consider the following instructions:

\begin{itemize}
\tightlist
\item
  The judge makes the \(Gs\) instruction with probability
  \(\nicefrac{1}{9}\).
\item
  The judge makes the \(Sg\) instruction with probability
  \(\nicefrac{1}{9}\).
\item
  The judge makes the \(Ss\) instruction with probability
  \(\nicefrac{7}{9}\).
\end{itemize}

If a player is told to go, they will know the other player is stopping
and will go. If a player is told to stop, they will think the
probability the other player is going is \(\nicefrac{1}{8}\), so they
will be indifferent between stopping and going, so they will not have an
incentive to deviate from judge's advice.

If the equilibrium is followed, each player expects the following
probability distribution over their returns:

\begin{itemize}
\tightlist
\item
  They get 10 with probability \(\nicefrac{1}{9}\);
\item
  They get 7 with probability \(\nicefrac{1}{9}\); and
\item
  They get 9 with probability \(\nicefrac{7}{9}\)
\end{itemize}

That means their expected return is \(8\nicefrac{8}{9}\), which is a
little larger than the mixed strategy equilibrium. And that's the
highest expected return of any correlated equilibrium.

The possibility of attractive correlated equilibria shows us something
about the value of \textbf{cheap talk}. Cheap talk is the general term
we use for any kind of costless signalling. There are all sorts of
interesting results that we can derive about costly signalling, and
we'll look at that later in the course. But it is somewhat intuitive
that cheap talk is never valuable. After all, it is of no use in games
like Prisoners Dilemma.

The existence of games with nice correlated equilibria shows that this
intuitive finding is actually false. Players who are allowed cheap talk
can work out a correlated equilibrium that they both prefer to
strategies that are available without cheap talk.

\section{Coordination Games}\label{sec-coordination-games}

Many games have multiple equilibria. In such cases, it is interesting to
work through the means by which one equilibrium rather than another may
end up being chosen. Consider, for instance, the following three games
(two of which we've seen before).

\begin{longtable}[]{@{}lrcc@{}}
\caption{Simple
coordination}\label{tbl-simple-coordination}\tabularnewline
\toprule\noalign{}
\endfirsthead
\endhead
\bottomrule\noalign{}
\endlastfoot
& & \multicolumn{2}{c@{}}{%
\textbf{Col}} \\
& & \(a\) & \(b\) \\
\multirow{2}{*}{\textbf{Row}} & \(A\) & 1,1 & 0,0 \\
& \(B\) & 0,0 & 1,1 \\
\end{longtable}

\begin{longtable}[]{@{}lrcc@{}}
\caption{Asymmetric
coordination}\label{tbl-asymmetric-coordination}\tabularnewline
\toprule\noalign{}
\endfirsthead
\endhead
\bottomrule\noalign{}
\endlastfoot
& & \multicolumn{2}{c@{}}{%
\textbf{Col}} \\
& & \(a\) & \(b\) \\
\multirow{2}{*}{\textbf{Row}} & \(A\) & 2,1 & 0,0 \\
& \(B\) & 0,0 & 1,2 \\
\end{longtable}

\begin{longtable}[]{@{}lrcc@{}}
\caption{Stag Hunt}\label{tbl-coord-stag}\tabularnewline
\toprule\noalign{}
\endfirsthead
\endhead
\bottomrule\noalign{}
\endlastfoot
& & \multicolumn{2}{c@{}}{%
\textbf{Col}} \\
& & \(a\) & \(b\) \\
\multirow{2}{*}{\textbf{Row}} & \(A\) & 5,5 & 0,4 \\
& \(B\) & 4,0 & 2,2 \\
\end{longtable}

In each case, both \(Aa\) and \(Bb\) are Nash equilibria.

The top game is a purely cooperative game; the players have exactly the
same payoffs in every outcome. It is a model for some real-world games.
The two players, Row and Col, have to meet up, and it doesn't matter
where. They could either go to location \(A\) or \(B\). If they go to
the same location, they are happy; if not, not.

In such an abstract presentation of game, it is hard to see how we could
select an equilibrium out of the two possibilities. In practice, it
often turns out not to be so hard. Thomas Schelling (1960) noted that a
lot of real-life versions of this game have what he called `focal
points'. (These are sometimes called Schelling points now, in his
honour.) A focal point is a point that stands out from the other options
in some salient respect, and it can be expected that other players will
notice that it stands out. Here's a nice example devised by the linguist
Christopher Potts (2008).

\begin{quote}
\(A\) and \(B\) have to select, without communicating, one of the
following nine figures. They each get a reward iff they select the same
figure.

\pandocbounded{\includegraphics[keepaspectratio]{index_files/mediabag/figures/svg/focal-point-1.pdf}}
\end{quote}

It's not obvious what they should do there, but compare the situation
our players face in the following game.

\begin{quote}
\(A\) and \(B\) have to select, without communicating, one of the
following nine figures. They each get a reward iff they select the same
figure.

\pandocbounded{\includegraphics[keepaspectratio]{index_files/mediabag/figures/svg/focal-point-2.pdf}}
\end{quote}

We could run experiments to test this, but intuitively, players will do
better at the second game than the first. That's because in the second
game, they have a focal point to select; one of the options stands out
from the crowd.

Schelling tested this by asking people where they would go if they had
to meet a stranger in a strange city, and had no way to communicate. The
answers suggested that meetups would be much more common than you might
expect. People in Paris would go to the Eiffel Tower; people in New York
would go to (the information booth at) Grand Central Station, and so on.
(This is all very contingent. I don't think Grand Central would still
work as a focal point in New York. Or maybe it would work as a focal
point for people who regularly caught east-bound regional trains from
New York.)

The second game we started this section with is often called `Battle of
the Sexes'. The real world example of it that is usually used is that
Row and Col are a married couple trying to coordinate on a night out.
But for whatever reason, they are at the stage where they simply have to
go to one of two venues. Row would prefer that they both go to \(A\),
while Col would prefer that they both go to \(B\). But the worst case
scenario is that they go to different things. I won't go over that game
much more, because the previous section told you most of what you need
to know about it. The game has two asymmetric pure strategy Nash
equilibria. It has a symmetric, but quite unfortunate, mixed strategy
equilibria. And it has a nice correlated equilibrium, where they flip a
coin to decide what to do. That's probably what real couples would do,
demonstrating again the importance of correlated equilibria.

The third coordination game I started the section with is called Stag
Hunt. The idea behind the game goes back to Rousseau ({[}1755{]} 2004),
although you can find similar examples in Hume ({[}1739{]} 2003). The
general idea is that \(A\) is some kind of cooperative action, and \(B\)
is a `defecting' action. If players make the same selection, then they
do well. If they both cooperate, they do better than if they both
defect. But if they don't cooperate, it is better to defect than to
cooperate.

Rousseau's example was of a group of hunters out hunting a stag, each of
whom sees a hare going by. If they leave the group, they will catch the
hare for sure, and have something to eat, but the group hunting the stag
will be left with a much lower probability of stag catching. Hume was
more explicit than Rousseau that these games come up with both two
player and many player versions. Here is a nice many player version.
I've adapted it from a game Ben Polak (2007) plays with his own class;
the threshold is his, the payoffs are simplified.

\begin{tcolorbox}[enhanced jigsaw, arc=.35mm, bottomrule=.15mm, bottomtitle=1mm, breakable, colback=white, colbacktitle=quarto-callout-tip-color!10!white, colframe=quarto-callout-tip-color-frame, coltitle=black, left=2mm, leftrule=.75mm, opacityback=0, opacitybacktitle=0.6, rightrule=.15mm, title=\textcolor{quarto-callout-tip-color}{\faLightbulb}\hspace{0.5em}{The Investment Game}, titlerule=0mm, toprule=.15mm, toptitle=1mm]

Everyone in the group has a choice between Investing and Not Investing.
The payoff to anyone who doesn't invest is 0. If 90\% or more of the
group Invests, then everyone who Invests gains \$1. If less than 90\% of
the group Invests, then everyone who invests loses \$1.

\end{tcolorbox}

Again, there are two Nash equilibria: everyone invests, and everyone
does not invest. Hume's hypothesis was that in games like this, the
cooperative move (in this case Investing) would be more likely in small
groups than large groups. This is basically what happened in Polak's
classes at Yale.

As well as the two pure strategy Nash equilibria in Stag Hunt, there is
a mixed strategy Nash equilibria. We find it, as always, by finding the
mixed strategy that makes the other player indifferent between their
choices. So assume that Row plays \(A\) with probability \(x\). Then the
expected return for Col from her two options is

\begin{align*}
E(a) &= x \times 5 + (1-x) \times 0 \\
 &= 5x \\
E(b) &= x \times 4 + (1-x) \times 2 \\
 &= 4x + 2 - 2x \\
 &= 2 + 2x
\end{align*}

So there will be equality only if \(5x = 2 + 2x\), i.e., only if
\(x = \nicefrac{2}{3}\). If that happens, \(C\)'s expected return from
any strategy will be \(\nicefrac{10}{3}\). Since the game is completely
symmetric, a very similar argument shows that if \(C\) plays the mixed
strategy of \(a\) with probability \(\nicefrac{2}{3}\), and \(b\) with
probability \(\nicefrac{1}{3}\), then Row will be indifferent between
her choices. So each player cooperating with probability
\(\nicefrac{2}{3}\) is a Nash equilibrium.

Unlike it Meeting or Battle of the Sexes, the mixed strategy equilibrium
is not worse for each player than either pure strategy equilibria. But
it isn't a great outcome either. In Stag Hunt, the mixed equilibrium is
better than the defecting equilibrium, but worse than the cooperating
equilibrium.

Say an equilibrium is \textbf{Pareto-preferred} to another equilibrium
iff every player would prefer the first equilibrium to the second. An
equilibrium is \textbf{Pareto-dominant} iff it is Pareto-preferred to
all other equilibria. The cooperative equilibrium is Pareto-dominant in
Stag Hunt; neither of the other games have Pareto-dominant equiilbria.

Consider each of the above games from the perspective of a player who
does not know what strategy their partner will play. Given a probability
distribution over the other player's moves, we can work out which
strategy has the higher expected return, or whether the various
strategies have the same expected returns as each other. For any
strategy \(S\), and possible strategy \(S^\prime\) of the other player,
let \(f\) be a function that maps \(S\) into the set of probabilities
\(x\) such that if the other player plays \(S^\prime\) with probability
\(x\), \(S\) has the highest expected return. In two strategy games, we
can remove the relativisation to \(S^\prime\), since for the purposes
we'll go on to, it doesn't matter which \(S^\prime\) we use. In somewhat
formal language, \(f(S)\) is the \textbf{basin of attraction} for \(S\);
it is the range of probability functions that points towards \(S\) being
played.

Let \(m\) be the usual (Lesbegue) measure over intervals; all that you
need to know about this measure is that for any interval \([x, y]\),
\(m([x, y]) = y - x\); i.e., it maps a continuous interval onto its
length. Say \(m(f(S))\) is the risk-attractiveness of \(S\).
Intuitively, this is a measure of how big a range of probability
distributions over the other person's play is compatible with \(S\)
being the best thing to play.

Say that a strategy \(S\) is risk-preferred iff \(m(f(S))\) is larger
than \(m(f(S^*))\) for any alternative strategy \(S^*\) available to
that agent. Say that an equilibrium is \textbf{risk-dominant} iff it
consists of risk-preferred strategies for all players.

For simple two-player, two-option games like we've been looking at, all
of this can be simplified a lot. An equilibrium is risk-dominant iff
each of the moves in it are moves the players would make if they
assigned probability \(\nicefrac{1}{2}\) to each of the possible moves
the other player could make.

Neither the simple Meeting game nor Battle of the Sexes have a
risk-dominant equilibrium. An asymmetric Meeting game, however, does.
Table~\ref{tbl-coord-favor} is one example.

\begin{longtable}[]{@{}lrcc@{}}
\caption{A coordination game with a favored
outcome}\label{tbl-coord-favor}\tabularnewline
\toprule\noalign{}
\endfirsthead
\endhead
\bottomrule\noalign{}
\endlastfoot
& & \multicolumn{2}{c@{}}{%
\textbf{Col}} \\
& & \(a\) & \(b\) \\
\multirow{2}{*}{\textbf{Row}} & \(A\) & 2,2 & 0,0 \\
& \(B\) & 0,0 & 1,1 \\
\end{longtable}

In this case \(Aa\) is both Pareto-dominant and risk-dominant. Given
that, we might expect that both players would choose \(A\). (What if an
option is Pareto-dominated, and risk-dominated, but focal? Should it be
chosen then? Perhaps; the Eiffel Tower isn't easy to get to, or
particularly pleasant once you're there, but seems to be chosen in the
Paris version of Schelling's meetup game because of its focality.)

The real interest of risk-dominance comes from Stag Hunt. In that game,
cooperating is Pareto-dominant, but defecting is risk-dominant. The
general class of Stag Hunt games is sometimes defined as the class of
two-player, two-option games with two equilibria, one of which is
Pareto-dominant, and the other of which is risk-dominant. In those
games, players have a choice between risk and reward. And at the same
time they have to choose between a socially useful choice, and a
socially harmful choice. The risky choice is the most socially useful
one. But note that the socially useful choice is not necessarily a
self-sacrifice. Assuming other players are being cooperative, each
player does better by being cooperative than by defecting.

It's at this point that the analogy to society formation becomes, for
some philosophers, quite vivid. The thought is that we don't require
coercion to get people to stay living in civilised society. Assuming
that a civilisation is up and running, each member of the society has a
preference for staying in the society to leaving it. But some
combination of coercion and public spiritedness might be required to get
a society running. If the players are unsure whether the others will
behave socially or anti-socially, the rational thing to do is to be
anti-social. And in the long run, that's a tragic outcome. I'll leave it
to other classes, however, to figure out how much social interactions
look like iterated Prisoners Dilemma, and how much they look like
iterated Stag Hunt.

\bookmarksetup{startatroot}

\chapter{Games and Time}\label{sec-games-and-time}

So far we've looked just at games where each player makes just one move,
and they make it simultaneously. That might not feel like most games
that you know about. It isn't how we play, for instance, chess. Instead,
most games involve players making multiple moves, and these moves taking
place in time. Here is one simple such game. It's not very interesting;
you won't have fun games nights playing it. But it is useful to study.

\begin{quote}
\textbf{Five}\\
There are two players, who we'll call \(A\) and \(B\). First \(A\)
moves, then \(B\), then finally \(A\) moves again. Each move involves
announcing a number, 1 or 2. \(A\) wins if after the three moves, the
numbers announced sum to 5. \(B\) wins otherwise.
\end{quote}

This is a simple zero-sum game. The payoff is either 1 to \(A\) and 0 to
\(B\), or 1 to \(B\) and 0 to \(A\). For simplicity, we'll describe this
as \(A\) winning or \(B\) winning. We'll soon be interested in games
with draws, which are a payoff of to each player. But for now we're
looking at games that someone wins.

Before we go on, it's worth thinking about how you would play this game
from each player's perspective. The formal approach we'll eventually
take is pretty similar, I think, to the way one thinks about the game.

\section{Normal Form and Extensive Form}\label{sec-normal-and-extensive}

Figure~\ref{fig-five} is the \textbf{extensive form} representation of
\textbf{Five}. We will use \(W\) to denote that \(A\) wins, and \(L\) to
denote that \(B\) wins.

\begin{figure}

\centering{

\pandocbounded{\includegraphics[keepaspectratio,alt={Extensive form of the game Five}]{index_files/mediabag/figures/svg/five.pdf}}

}

\caption{\label{fig-five}Five}

\end{figure}%

To read what's going on here, start at the node that isn't filled in,
which in this case is the node at the bottom of the tree. The first move
is made by \(A\). She chooses to play either 1 or 2. If she plays 1, we
go to the left of the chart, and now \(B\) has a choice. She plays
either 1 or 2, and now again \(A\) has a choice. As it turns out, \(A\)
has the same choice to make whatever \(B\) does, though there is nothing
essential to this. This is just a two-player game, though there is also
nothing essential to this. We could easily have let it be the case that
a third player moved at this point. Or we could have made it that who
moved at this point depended on which move \(B\) had made. The extensive
form representation allows for a lot of flexibility in this respect.

At the top of the diagram is a record of who won. In the more general
form, we would put the payoffs to each player here. But when we have a
simple zero-sum game, it is easier to just record the payoffs to one
player. You should check that the description of who wins and loses in
each situation is right. In each case, \(A\) wins if 2 is played twice,
and 1 once.

The extensive form representation form is very convenient for turn
taking games. But you might think that it is much less convenient for
games where players move simultaneously, as in Prisoners' Dilemma. But
it turns out that we can represent that as well, through a small
notational innovation. Figure~\ref{fig-pd-tree} is the game tree for
Prisoners' Dilemma.

\begin{figure}

\centering{

\pandocbounded{\includegraphics[keepaspectratio,alt={Prisoners Dilemma in extensive form, with an information set}]{index_files/mediabag/figures/svg/pd-tree.pdf}}

}

\caption{\label{fig-pd-tree}Prisoners' Dilemma}

\end{figure}%

The crucial thing here is the dashed line between the two nodes where P2
(short for Player 2) moves. What this means is that P2 doesn't know
which of these nodes she is at when she moves. We normally assume that a
player knows what moves are available to her. So we normally only put
this kind of dashed line in when it connects two nodes at which the same
player moves, and the available moves are the same. When we put in
notation like this, we say that the nodes that are connected form an
\textbf{information set}. If we don't mark anything, we assume that a
node is in a degenerate information set, one that only contains itself.

Strictly speaking, you could regard this tree as a representation of a
game where Player 1 goes first, then Player 2 moves, but Player 2 does
not know Player 1's move when she moves. But that would be reading, I
think, too much into the symbolic representation. What's crucial is that
we can represent simultaneous move games in extensive form.

There is a philosophical assumption in this notational convention that I
just want to mention. We won't have time to cover it further here,
though it is worth thinking about. It is usual to use dashed lines, like
I've done, or circles or ovals to represent that for all Player 2 knows,
she is at one of the connected nodes. But drawing circles around a set
of possibilities is only a good way to represent Player 2's uncertainty
if we assume quite a lot about knowledge. In particular, it is only good
if we assume that which nodes are epistemically open for Player 2 is
independent of which node, within that group, she is at. In formal
terms, it amounts to assuming that each player knows exactly what they
know, and what they don't know.\footnote{The assumption that if someone
  knows something, they know that they know it, is famously criticised
  by Timothy Williamson (2000). This has been incredibly controversial
  in philosophy. The assumption that if someone does not know something,
  they know that they do not know it, is much more widely rejected in
  philosophy. The most careful statement of why we should reject it that
  I know is by Lloyd Humberstone (2015, 380--401).} This is not, in
general, a true assumption. Normal people have some false beliefs. Those
false beliefs are not pieces of knowledge, but the holders of the false
belief do not in general know that they do not amount to knowledge. A
more philosophically defensible game theory would not assume that
everyone knows what it is that they know. But for now we'll follow
orthodoxy and assume we can use representations like these dashed lines
to represent uncertainty, even though that requires a lot more
self-knowledge than most people have.

These representations using graphs are knows as \textbf{extensive form}
representations of games. The representations using tables that we've
used previously are known as \textbf{normal form} or \textbf{strategic
form} representations. The idea is that any game really can be
represented as game of one simultaneous move. The `moves' the players
make are the selections of \textbf{strategies}. A strategy in this sense
is a plan for what to do in any circumstance whatsoever.

Let's do this for the \textbf{Five}. A strategy for \(A\) has three
variables. It must specify what she does at the first move, what she
does at the second move if \(B\) plays 1, and what she does at the
second move if \(B\) plays 2. (We're assuming here that \(A\) will play
what she decides to play at the first move. It's possible to drop that
assumption, but it results in much more complexity.) So we'll describe
\(A\)'s strategy as \(\alpha \beta \gamma\), where \(\alpha\) is her
move to begin with, \(\beta\) is what she does if \(B\) plays 1, and
\(\gamma\) is what she does if \(B\) plays 2. A strategy for \(B\) needs
to only have two variables: what to do if \(A\) plays 1, and what to do
if \(A\) plays 2. So we'll notate her strategy as \(\delta \epsilon\),
where \(\delta\) is what she does if \(A\) plays 1, and \(\epsilon\) is
what she does if \(A\) plays 2. So \(A\) has 8 possible strategies, and
\(B\) has 4 possible strategies. Let's record the giant table listing
the outcomes if they play each of those strategies.

{\def\LTcaptype{none} % do not increment counter
\begin{longtable}[]{@{}lrcccc@{}}
\toprule\noalign{}
\endhead
\bottomrule\noalign{}
\endlastfoot
& & \multicolumn{4}{c@{}}{%
\(B\)} \\
& & 11 & 12 & 21 & 22 \\
\multirow{8}{*}{\(A\)} & 111 & \(L\) & \(L\) & \(L\) & \(L\) \\
& 112 & \(L\) & \(L\) & \(W\) & \(W\) \\
& 121 & \(L\) & \(L\) & \(L\) & \(L\) \\
& 122 & \(L\) & \(L\) & \(W\) & \(W\) \\
& 211 & \(L\) & \(W\) & \(L\) & \(W\) \\
& 212 & \(L\) & \(L\) & \(L\) & \(L\) \\
& 221 & \(W\) & \(W\) & \(W\) & \(W\) \\
& 222 & \(W\) & \(L\) & \(W\) & \(L\) \\
\end{longtable}
}

There is something quite dramatic about this representation. We can see
what \(A\) should play. If her strategy is 221, then whatever strategy
\(B\) plays, \(A\) wins. So she should play that; it is a (weakly)
dominant strategy. This isn't completely obvious from the extended form
graph.

Here's a related fact. Note that there are only 8 outcomes of the
extended form game, but 32 cells in the table. Each outcome on the tree
is represented by multiple cells of the table. Let's say we changed the
game so that it finishes in a draw, represented by \(D\), if the numbers
picked sum to 3. That just requires changing one thing on the graph; the
\(L\) in the top-left corner has to be changed to a \(D\). But it
requires making many changes to the table.

{\def\LTcaptype{none} % do not increment counter
\begin{longtable}[]{@{}lrcccc@{}}
\toprule\noalign{}
\endhead
\bottomrule\noalign{}
\endlastfoot
& & \multicolumn{4}{c@{}}{%
\(B\)} \\
& & 11 & 12 & 21 & 22 \\
\multirow{8}{*}{\(A\)} & 111 & \(D\) & \(D\) & \(L\) & \(L\) \\
& 112 & \(D\) & \(D\) & \(W\) & \(W\) \\
& 121 & \(L\) & \(L\) & \(L\) & \(L\) \\
& 122 & \(L\) & \(L\) & \(W\) & \(W\) \\
& 211 & \(L\) & \(W\) & \(L\) & \(W\) \\
& 212 & \(L\) & \(L\) & \(L\) & \(L\) \\
& 221 & \(W\) & \(W\) & \(W\) & \(W\) \\
& 222 & \(W\) & \(L\) & \(W\) & \(L\) \\
\end{longtable}
}

In part because of this fact, i.e., because every change to the value
assignment in the extensive form requires making many changes in the
values on the normal form, it isn't a coincidence that there's a row
containing nothing but \(W\). The following claim about our game can be
proved.

\begin{quote}
Assign \(W\) and \(L\) in any way you like to the eight outcomes of the
extended form game. Then draw table that is the normal form
representation of the game. It will either have a row containing nothing
but \(W\), i.e., a winning strategy for \(A\), or a column containing
nothing but \(L\), i.e., a winning strategy for \(B\).
\end{quote}

We will prove this in a little bit, but first we will look at how to
`solve' games like \textbf{Five}.

\section{Backwards Induction}\label{sec-backwards-induction}

The way to think through games like \textbf{Five} is by working from top
to bottom. \(A\) moves last. Once we get to the point of the last move,
there is no tactical decision making needed. \(A\) knows what payoff she
gets from each move, and she simply will take the highest payoff
(assuming rationality).

So let's assume she does that. Let's assume, that is, that \(A\) does
play her best strategy. Then we know three things about what will happen
in the game.

\begin{itemize}
\item
  If \(A\) plays 1, and \(B\) plays 2, \(A\) will follow with 2.
\item
  If \(A\) plays 2, and \(B\) plays 1, \(A\) will follow with 2.
\item
  If \(A\) plays 2, and \(B\) plays 2, \(A\) will follow with 1.
\end{itemize}

Moreover, once we make this assumption there is, in effect, one fewer
step in the game. Once \(B\) moves, the outcome is determined. So let's
redraw the game using that assumption, and just listing payoffs after
the second move. This will be Figure~\ref{fig-five-2}

\begin{figure}

\centering{

\pandocbounded{\includegraphics[keepaspectratio,alt={Five, with the last move assumed}]{index_files/mediabag/figures/svg/five-2.pdf}}

}

\caption{\label{fig-five-2}Five with last move assumed}

\end{figure}%

Now we can make the same assumption about \(B\). Assume that \(B\) will
simply make her best move in this (reduced) game. Since \(B\) wins if
\(A\) loses, \(B\)'s best move is to get to \(L\). This assumption then,
gives us just one extra constraint.

\begin{itemize}
\tightlist
\item
  If \(A\) plays 1, \(B\) will follow with 1.
\end{itemize}

And, once again, we can replace \(B\)'s actual movement with a payoff of
the game under the assumption that \(B\) makes the rational move. This
gives us an even simpler representation of the game that we see in
Figure~\ref{fig-five-3}.

\begin{figure}

\centering{

\pandocbounded{\includegraphics[keepaspectratio,alt={Five, with the last two moves assumed}]{index_files/mediabag/figures/svg/five-3.pdf}}

}

\caption{\label{fig-five-3}Five with last two moves assumed}

\end{figure}%

And from this version of the game, we can draw two more conclusions,
assuming \(A\) is rational.

\begin{itemize}
\item
  \(A\) will play 2 at the first move.
\item
  \(A\) will win.
\end{itemize}

Let's put all of that together. We know \(A\) will start with 2, so her
strategy will be of the form \(2 \beta \gamma\). We also know that \(B\)
doesn't care which strategy she chooses at that point, so we can't make
any further reductions. But we do know that if \(A\) plays 2 and \(B\)
plays 1, \(A\) will follow with 2. So \(A\)'s strategy will be of the
form \(2 2 \gamma\). And we know that if \(A\) plays 2 and \(B\) plays
2, then \(A\) will play 1. So \(A\)'s strategy will be 221, as we saw on
the table.

Note that the use of backwards induction here hides a multitude of
assumptions. We have to assume each player is rational, and each player
knows that, and each player knows that, and so on for at least as many
iterations as there are steps in the game. If we didn't have those
assumptions, it wouldn't be right to simply replace a huge tree with a
single outcome. Moreover, we have to make those assumptions be very
modally robust.

We can see this with a slight variant of the game. Let's say this time
that the left two outcomes are \(D\). So the graph looks like
Figure~\ref{fig-five-rev}.

\begin{figure}

\centering{

\pandocbounded{\includegraphics[keepaspectratio,alt={The revised game Five-double-prime}]{index_files/mediabag/figures/svg/five-rev.pdf}}

}

\caption{\label{fig-five-rev}Five'\,'}

\end{figure}%

Now we assume that \(A\) makes the optimal move at the last step, so we
can replace the top row of outcomes with the outcome that would happen
if \(A\) moves optimally. This gets us Figure~\ref{fig-five-rev-2}.

\begin{figure}

\centering{

\pandocbounded{\includegraphics[keepaspectratio,alt={Five-double-prime, with the last move assumed}]{index_files/mediabag/figures/svg/five-rev-2.pdf}}

}

\caption{\label{fig-five-rev-2}Five'\,' with last move assumed}

\end{figure}%

Now assume that \(A\) plays 1 on the first round, then \(B\) has to
move. From Figure~\ref{fig-five-rev-2} it looks like \(B\) has an easy
choice to make. If she plays 1, she gets a draw, if she plays 2, then
\(A\) wins, i.e., she loses. Since drawing is better than losing, she
should play 1 and take the draw.

But why think that playing 2 will lead to \(A\) winning? The argument
that it did depending on assuming \(A\) is perfectly rational. And
assuming \(B\) is in a position to make this choice, that seems like an
unlikely assumption. After all, if \(A\) were perfectly rational, she'd
have chosen 2, and given herself a chance to force a win.

Now you might think that even if \(A\) isn't perfectly rational, it
still is crazy to leave her with an easy winning move. And that's
probably a sufficient reason for \(B\) to accept the draw, i.e., play 1.
But the argument that \(B\) should regard playing 2 as equivalent to
choosing defeat seems mistaken. \(B\) knows that \(A\) isn't perfectly
rational, and she shouldn't assume perfect rationality from here on out.

We will come back to this point, a lot, in subsequent discussions of
backwards induction. Note that it is a point that doesn't really arise
in the context of normal form games. There we might wonder about whether
common knowledge of rationality is a legitimate assumption at the
\emph{start} of the game. But once we've settled that, we don't have a
further issue to decide about whether it is still a legitimate
assumption at later stages of the game.

\section{Value of Games}\label{sec-value-of-games}

Consider games with the following characteristics.

\begin{itemize}
\tightlist
\item
  \(A\) and \(B\) take turns making moves. We will call each point at
  which they make a move, or at which the game ends, a \textbf{node} of
  the game.
\item
  At each move, each player knows what moves have been previously made.
\item
  At each move, the players have only finitely many possible moves open.
\item
  The players' preferences over the outcomes are opposite to one
  another. So if \(A\) prefers outcome \(o_1\) to \(o_2\), then \(B\)
  prefers \(o_2\) to \(o_1\), and if \(A\) is indifferent between
  \(o_1\) and \(o_2\), then \(B\) is indifferent between them as well.
\item
  \(A\)'s preferences over the outcomes are complete; for any two
  outcomes, she either prefers the first, or prefers the second, or is
  indifferent between them.
\item
  There is a finite limit to the total number of possible moves in the
  game.
\end{itemize}

The finiteness assumptions entail that there are only finitely many
possible outcomes. So we can order the outcomes by \(A\)'s preferences.
(Strictly speaking, we want to order sets of outcomes that are
equivalence classes with respect to the relation that \(A\) is
indifferent between them.) Assign the outcome \(A\) least prefers the
value 0, the next outcome the value 1, and so on.

Now we recursively define the \textbf{value} of a node as follows.

\begin{itemize}
\tightlist
\item
  The value of a \textbf{terminal node}, i.e., a node at which the game
  ends, is the payoff at that node.
\item
  The value of any node at which \(A\) makes a choice is the greatest
  value of the nodes between which \(A\) is choosing to move to.
\item
  The value of any node at which \(B\) makes a choice is the least value
  of the nodes between which \(B\) is choosing to move to.
\end{itemize}

Finally, we say that the value of the game is the value of the initial
node, i.e., the node at which the first choice is made. We can prove a
number of things about the value of games. The proof will make crucial
use of the notion of a \textbf{subgame}. In any extensive form game, a
\textbf{subgame} of a game is the game we get by treating any perfect
information node as the initial node of a game, and including the rest
of the game `downstream' from there.

By a perfect information node, I mean a node such that when it is
reached, it is common knowledge that it is reached. That's true of all
the nodes in most of the games we're looking at, but it isn't true in,
for instance, the extensive form version of Prisoners' Dilemma we looked
at. Nodes that are in non-degenerate information sets, i.e., information
sets that contain other nodes, can't trigger subgames. That's because we
typically assume that to play a game, players have to know what game
they are playing.

Note that a subgame is really a game, just like a subset is a set. Once
we're in the subgame, it doesn't matter a lot how we got there. Indeed,
any game we represent is the consequence of some choices by the agent;
they are all subgames of the game of life.

\begin{itemize}
\tightlist
\item
  The value of a game is the value of one of the terminal nodes.
\end{itemize}

We prove this by induction on the length of games. (The length of a game
is the \emph{maximum} number of moves needed to reach a terminal node.
We've only looked so far at games where every path takes the same number
of moves to reach a conclusion, but that's not a compulsory feature of
games.)

If the game has zero moves, then it has just one node, and its value is
the value of that node. And that node is a terminal node, so the value
is the value of a terminal node.

Now assume that the claim is true for any game of length \(k\) or less,
and consider an arbitrary game of length \(k + 1\) The first node of the
game consists of a choice about which path to go down. So the value of
the initial node is the value of one of the subsequent nodes. Once that
choice is made, we are in a subgame of length \(k\), no matter which
choice is made. By the inductive hypothesis, the value of that subgame
is the value of one of its terminal nodes. So the value of the game,
which is the value of one of the immediate subsequent nodes to the
initial node, is the value of one of its terminal nodes.

\begin{itemize}
\tightlist
\item
  \(A\) can guarantee that the outcome of the game is at least the value
  of the game.
\end{itemize}

Again, we prove this by induction on the length of games. It is trivial
for games of length 0. So assume it is true for all games of length at
most \(k\), and we'll prove it for games of length \(k+1\). The initial
node of such a game is either a move by \(A\) or a move by \(B\). We
will consider these two cases separately.

Assume that \(A\) makes the first move. Then the value of the initial
node is the maximum value of any immediate successor node. So \(A\) can
select to go to the node with the same value as the value of the game.
Then we're in a subgame of length \(k\). By the inductive assumption, in
that game \(A\) can guarantee that the outcome is at least the value of
the subgame. And since the value of that node is the value of the
subgame, so it is also the value of the initial node, i.e., the value of
the initial game. So by choosing that node, and starting that subgame,
\(A\) can guarantee that the outcome is at least the value of the game.

Now assume that \(B\) makes the first move. \(B\) can choose the node
with the least value of the available choices. Then, as above, we'll be
in a subgame in which (by the inductive hypothesis) \(B\) can guarantee
has an outcome which is at most its value. That is, \(B\) can guarantee
the outcome of the game is at most the value of the initial node of the
subgame. And since \(B\) can guarantee that that subgame is played,
\(B\) can guarantee that the game has an outcome of at most its value.

\begin{itemize}
\tightlist
\item
  \(B\) can guarantee that the outcome of the game is at most the value
  of the game.
\end{itemize}

The proof of this exactly parallels the previous proof, and the details
are left as an exercise.

Let's note a couple of consequences of these theorems.

First, assume that the rationality of each player is common knowledge,
and that it is also common knowledge that this will persist throughout
the game. Then the kind of backwards induction argument we used is
discussing \textbf{Five} will show that the outcome of the game will be
the value of the game. That's because if \(A\) is rational and knows
this much about \(B\), the outcome won't be lower than the value, and if
\(B\) is rational and knows this much about \(A\), the outcome won't be
greater than the value.

Second, these theorems have many applications to real-life games. Both
chess and checkers, for instance, satisfy the conditions we listed. The
only condition that is not immediately obvious in each case is that the
game ends in finite time. But the rules for draws in each game guarantee
that is true. (Exercise: Prove this!) Since these games end with White
win, Black win or draw, the value of the game must be one of those three
outcomes.

In the case of checkers, we know what the value of the game is. It is a
draw. This was proved by the Chinook project at the University of
Alberta (Schaeffer et al. 2007). We don't yet know what the value of
chess is, but it is probably also a draw. Given how many possible moves
there are in chess, and in principle how long games can go on, it is
hard to believe that chess will be `solved' any time soon. But advances
in chess computers may be able to justify to everyone's satisfaction a
particular solution, even if we can't prove that is the value of chess.

\section{Subgame Perfect Equilibrium}\label{sec-subgame-perfect}

Consider the following little game. The two players, call them Player I
and Player II, have a choice between two options, call them Good and
Bad. The game is a sequential move game; first Player I moves then
Player II moves. Each player gets 1 if they choose Good and 0 if they
choose Bad. We will refer to this game as \textbf{Easy Game}. Its game
tree is Figure~\ref{fig-easy-game}.

\begin{figure}

\centering{

\pandocbounded{\includegraphics[keepaspectratio,alt={The Easy Game tree}]{index_files/mediabag/figures/svg/easy-game.pdf}}

}

\caption{\label{fig-easy-game}Easy Game}

\end{figure}%

A strategy for Player I in \textbf{Easy Game} is just a choice of one
option, Good or Bad. A strategy for Player II is a little more
complicated. She has to choose both what to do if Player I chooses Good,
and what to do if Player II chooses Bad. We'll write her strategy as
\(\alpha \beta\), where \(\alpha\) is what she does if Player I chooses
Good, and \(\beta\) is what she does if Player II chooses Bad. (We will
often use this kind of notation in what follows. Wherever it is
potentially ambiguous, I'll try to explain it. But the notation is very
common in works on game theory, and it is worth knowing.)

The most obvious Nash equilibrium of the game is that Player I chooses
Good, and Player II chooses Good whatever Player I does. But there is
another Nash equilibrium, as looking at the strategic form of the game
reveals. We'll put Player I on the row and Player II on the column.
(We'll also do this from now on unless there is a reason to do
otherwise.)

{\def\LTcaptype{none} % do not increment counter
\begin{longtable}[]{@{}lrcccc@{}}
\toprule\noalign{}
\endhead
\bottomrule\noalign{}
\endlastfoot
& & \multicolumn{4}{c@{}}{%
Player II} \\
& & gg & gb & bg & bb \\
\multirow{2}{*}{Player I} & G & 1, 1 & \textbf{1, 1} & 1, 0 & 1, 0 \\
& B & 0, 1 & 0, 0 & 0, 1 & 0, 0 \\
\end{longtable}
}

Look at the cell that I've bolded, where Player I plays Good, and Player
II plays Good, Bad. That's a Nash equilibrium. Neither player can
improve their outcome, given what the other player plays. But it is a
very odd Nash equilibrium. It would be very odd for Player II to play
this, since it risks getting 0 when they can guarantee getting 1.

It's true that Good, Bad is weakly dominated by Good, Good. But as we've
already seen, and as we'll see in very similar examples to this soon,
there are dangers in throwing out \emph{all} weakly dominated
strategies. Many people think that there is something else wrong with
what Player II does here.

Consider the sub-game that starts with the right-hand decision node for
Play\-er II. That isn't a very interesting game; Player I has no
choices, and Player II simply has a choice between 1 and 0. But it is a
game. And note that Good, Bad is not a Nash equilibrium~of that game.
Indeed, it is a \emph{strictly} dominated strategy in that game, since
it involves taking 0 when 1 is freely available.

Say that a strategy pair is a \textbf{subgame perfect equilibrium} when
it is a Nash equilibrium, and it is a Nash equilibrium~of every sub-game
of the game. The pair Good and Good, Bad is not subgame perfect, since
it is not a Nash equilibrium~of the right-hand subgame.

When we solve extensive form games by backwards induction, we not only
find Nash equilibria, but subgame perfect equilibria. Solving this game
by backwards induction would reveal that Player II would choose Good
wherever she ends up, and then Player I will play Good at the first
move. And the only subgame perfect equilibrium~of the game is that
Player I plays Good, and Player II plays Good, Good.

The stakes here don't seem to high, since the ruled out Nash equilibrium
has the same payout as the subgame perfect equilibrium. But that need
not be the case. Sometimes the equilibrium that is not subgame perfect
is an \textbf{incredible threat}. We mean \emph{incredible} here in the
literal sense of not credible. In the spirit of this, let's consider a
game I'll call \textbf{Incredible}. Its game tree is
Figure~\ref{fig-incredible}.

\begin{figure}

\centering{

\pandocbounded{\includegraphics[keepaspectratio,alt={The game Incredible, showing an incredible threat}]{index_files/mediabag/figures/svg/incredible.pdf}}

}

\caption{\label{fig-incredible}Incredible}

\end{figure}%

Using backwards induction, we can see that the only subgame perfect
equilibrium of the game is \(Rbg\). That is, Player I plays \(R\), and
Player II is disposed to play \(b\) if Player I plays Left, and \(g\) if
Player plays \(R\). But writing out the able reveals another Nash
equilibria.

{\def\LTcaptype{none} % do not increment counter
\begin{longtable}[]{@{}lrcccc@{}}
\toprule\noalign{}
\endhead
\bottomrule\noalign{}
\endlastfoot
& & \multicolumn{4}{c@{}}{%
Player II} \\
& & \(gg\) & \(gb\) & \(bg\) & \(bb\) \\
\multirow{2}{*}{Player I} & \(L\) & 4,3 & 4,3 & 1,4 & \textbf{1,4} \\
& \(R\) & 3,2 & 0,0 & \textbf{3,2} & 0,0 \\
\end{longtable}
}

I've bolded both the Nash equilibria. And note the one in the top right
corner of the table. If Player II is committed to playing \(bb\), then
the best response Player I can make is to play \(L\).

But look how implausible this is in real time. It means that if we get
to the right hand node, Player II will choose to get 0, when 2 is
available. Of course, we won't actually get to that node. But it
requires Player II to have dispositions that make no sense. This is the
intuition behind subgame perfect equilibrium. We require that players
not just make choices which make sense, but are disposed to do so at all
times.

In the rest of this chapter, we'll look at reasons why this might be too
strong a constraint on rational players.

\section{Problems with Backwards
Induction}\label{sec-problems-backwards-induction}

When each player only moves once, it is very plausible that each player
should play their part of a subgame perfect equilibrium~solution to the
game. But this is less compelling when players have multiple moves. We
can start to see why by looking at a game that the philosopher Robert
Stalnaker (1996, 1998) has used in a few places to argue against
backwards induction reasoning. Again, it is an extensive form game. We
will call it \textbf{Cooperative Caterpillar}.

\begin{figure}

\centering{

\pandocbounded{\includegraphics[keepaspectratio,alt={The Cooperative Caterpillar}]{index_files/mediabag/figures/svg/caterpillar.pdf}}

}

\caption{\label{fig-caterpillar}Cooperative Caterpillar}

\end{figure}%

The game starts in the upper-left corner. There are up to three moves,
each of them Across or Down. As soon as one player moves Down, the game
ends. It is a common interest game; each player gets the payout at the
terminal node.

Since there is only one path to each decision node, a strategy merely
has to consist of a set of plans for what to play if that node is
reached. Alice's strategy will consist of two capitalised moves, and
Bob's strategy will consist of one lower-case move.

If we apply backwards induction, we get that the unique solution of the
game is \(\langle A_1A_2, a \rangle\). Alice would play \(A_2\) if it
gets that far. Given that, Bob would be better off playing \(a\) than
\(d\) if he gets to play. And given that, Alice is better off playing
\(A_1\) than \(D_1\).

But there are many other Nash equilibria of the game. One of them is
\(\langle D_1A_2, d \rangle\). Given that Player I is playing \(D_1\),
Player II can play anything without changing her payout; it will be 2
come what may. Given that Player II is playing \(d\), Player I is best
off playing \(D_1\) and taking 2, rather than just getting the 1 that
comes from leaving Player II to make a play.

Could this equilibrium~be one that rational players, who know each other
to be rational, reach? Stalnaker argues that it could. Assume each
player knows that the other player is rational, and is playing that
strategy. Given what Player I knows, her strategy choice is clearly
rational. She takes 2 rather than the 1 that she would get by playing
\(A_1\), and she is disposed to take 3 rather than 0 if she gets to the
final decision node. So her actual choice is rational, and her
disposition is to make another rational choice if we reach the end of
the game.

Things are a little trickier for Player II. You might think it is
impossible for rational Player II, who knows Player I to be rational, to
move \(d\). After all, if Player I is rational, then she'll play
\(A_2\), not \(D_2\). And if she plays \(A_2\), it is better to play
\(a\) than \(d\). So it looks hard to justify Player II's move. But
looks can be deceiving. In fact there isn't anything wrong with Player
II's move, as long as he has the right beliefs to justify it. It's very
important to distinguish the following two conditionals.

\begin{itemize}
\tightlist
\item
  If Player I has the choice, she chooses \(A_2\) over \(D_2\).
\item
  If Player I were to have the choice, she would choose \(A_2\) over
  \(D_2\).
\end{itemize}

Player II knows that the first, indicative, conditional is true. And
indeed it is true. But he doesn't know that the second, subjunctive,
conditional is true. After all, if Player I were to have the choice
between \(A_2\) and \(D_2\), she would have, \emph{irrationally}, chosen
\(A_1\) over \(D_1\). And if she had chosen irrationally once, it's
possible that she would choose irrationally again.

Here's an analogy that may help explain what's going on. The following
set of beliefs is consistent.

\begin{itemize}
\tightlist
\item
  Any perfectly rational being gets all of the questions on their
  algebra exam right.
\item
  Alice is perfectly rational.
\item
  If Alice had got the second-hardest question on the algebra exam
  wrong, she would have got the hardest question on the algebra exam
  wrong as well.
\end{itemize}

Player II's beliefs are like that. He believes Player I is perfectly
rational. He also believes that if Player I were to make an irrational
move, she would continue to make irrational moves. That's consistent
with belief in perfect rationality, and nothing about the game setup
rules out such a belief. He also believes that playing \(A_1\) would be
irrational. That's correct, given what Player I knows about Player II.
Given all those beliefs, playing \(d\), if he had the chance, would be
rational.

Stalnaker argues that many game theorists have tacitly confused
indicative and subjunctive conditionals in reasoning about games like
\textbf{Cooperative Caterpillar}. There are other games where similar
confusions arguably take place.

\section{The Centipede}\label{sec-centipede}

In \textbf{Cooperative Caterpillar}, backwards induction reasoning gets
the players to the outcome that many people would say is intuitive.
Stalnaker's objection is that the reasoning doesn't go through in
theory, but at least you might say it works in practice. Once we make
the game less cooperative, backwards induction reasoning ceases to be
either intuitive, or to match what we see in experimental settings.

The game is due to Rosenthal (1981), and is usually called the
\textbf{centipede}, because that is what the diagram looks like once you
draw enough legs on it.

Two players alternate. At each turn a player can take the larger share
of a growing pot, ending the game, or pass, in which case the pot grows
and the other player faces the same choice. Table~\ref{tbl-centipede} is
a short version, with four decision nodes.

\begin{longtable}[]{@{}lccccc@{}}
\caption{A short version of
centipede}\label{tbl-centipede}\tabularnewline
\toprule\noalign{}
\endfirsthead
\endhead
\bottomrule\noalign{}
\endlastfoot
\textbf{Node} & 1 & 2 & 3 & 4 & end \\
\textbf{Mover} & Alice & Bob & Alice & Bob & \\
\textbf{Take} & 2, 0 & 1, 3 & 4, 2 & 3, 5 & \\
\textbf{Pass} & on & on & on & on & 6, 4 \\
\end{longtable}

Read that as follows. If Alice takes at node 1, she gets 2 and Bob gets
0. If she passes, Bob chooses at node 2, where taking gets him 3 and
gets Alice 1. And so on. If everyone passes at every node, they end at
6, 4.

Notice two things about the payoffs. First, the total grows: 2 at the
first node, 4 at the second, 6 at the third, 8 at the fourth, 10 at the
end. Passing always makes the pie bigger. Second, at each node the mover
does better by taking now than by taking one round later. Alice gets 2
by taking at node 1 and only 1 if Bob takes at node 2.

Now look what happens when we step through the backwards induction
reasoning. At node 4, Bob compares taking, which gets him 5, with
passing, which gets him 4. So Bob takes. Given that, at node 3 Alice
compares taking, which gets her 4, with passing, which leads to Bob
taking and gets her 3. So Alice takes. Given that, at node 2 Bob
compares 3 with the 2 he gets when Alice takes at node 3. So Bob takes.
And given that, at node 1 Alice compares 2 with the 1 she gets when Bob
takes at node 2. So Alice takes at node 1.

The unique subgame perfect equilibrium is that Alice takes immediately,
and the players walk away with 2 and 0. They never reach the part of the
tree where the money is. Moreover, we could have let the game go on
arbitrarily long, and this would have been the same result. As long as
there's a final node, where one player knows the game will end no matter
what they do, backwards induction reasoning can get started, and it
results in Alice taking at node 1.

This is really strange. No individual step is strange. Bob really should
take rather than pass at the last step. It is fairly plausible that
Alice should take rather than pass at the second last step. But if we
add up enough of these plausible steps, we get something quite
implausible.

The implausibility is not just that they end up with a result that is
worse, for both of them, than what they could have got. One-shot
prisoners' dilemma has that structure, and we mostly think that it's
tragic, rather than irrational, that the players end up in a
Pareto-inferior outcome.

Rather, the implausibility is that it feels like there is a way for them
to have done better. By passing in the early rounds, they can give each
other good reason to think that they will keep on passing, and hence a
reason to pass themselves. Think about what happens too the backwards
induction reasoning if anyone deviates. (This will be just like what
Stalnaker says goes wrong in \textbf{Cooperative Caterpillar}.) Suppose
Alice passes at node 1. Bob now has to decide what to do at a node that
backwards induction says should never be reached. The argument for
taking at node 2 ran through the claim that Alice will take at node 3,
and that claim was derived from the assumption that Alice is rational.
But Alice has just done the thing that assumption ruled out.

So what should Bob believe? There are at least three answers, and each
has been defended.

\begin{itemize}
\tightlist
\item
  Alice made a mistake, a slip of the hand, and is otherwise rational.
  Then Bob should still expect her to take at node 3, and should take
  now.
\item
  Alice is not rational after all. Then Bob has no idea what she will do
  at node 3, and the backwards induction argument for taking at node 2
  has lost its premise.
\item
  Alice is rational, and is signalling that she intends to pass again,
  in the hope that Bob will pass too and they will both do better. This
  is forward induction reasoning of the kind we will see in the
  \textbf{Money Burning Game} (Section~\ref{sec-money-burning}), and on
  this reading Bob should pass.
\end{itemize}

The centipede is therefore not just a puzzle about whether backwards
induction gives the right answer. It is a puzzle about what the
assumption of common knowledge of rationality is even saying, given that
the reasoning it licenses runs through nodes that it says will never be
reached. Stalnaker's discussion of \textbf{Cooperative Caterpillar} is
designed to show that we can see the same mistake in backwards induction
reasoning even in cases where it gets the intuitive right answer, and
that hence the first option here is mistaken.

There is an empirical footnote. When people play centipede games in the
laboratory, they mostly do not take at the first node (McKelvey and
Palfrey 1992). They pass for a while and take somewhere in the middle.
That does not settle anything about what is rational, but it does mean
the theory is predicting something that does not happen, which is at
least a reason to look at the theory again.

\section{Money Burning Game}\label{sec-money-burning}

Elchanen Ben-Porath and Eddie Dekel (1992) suggested this variant to the
Battle of the Sexes game. The variation is in two steps. First, we
change the payouts for the basic game to those in
Table~\ref{tbl-bpd-bos}. (Note that I'm using a lowercase letter for one
of Row's options here; this is to distinguish the \(d\) of down from the
\(D\) of Don't burn.)

\begin{longtable}[]{@{}lrcc@{}}
\caption{The version of Battle of the Sexes being
used}\label{tbl-bpd-bos}\tabularnewline
\toprule\noalign{}
\endfirsthead
\endhead
\bottomrule\noalign{}
\endlastfoot
& & \multicolumn{2}{c@{}}{%
Player II} \\
& & \(l\) & \(r\) \\
\multirow{2}{*}{Player I} & \(u\) & 4, 1 & 0, 0 \\
& \(d\) & 0, 0 & 1, 4 \\
\end{longtable}

Then we give Player I the option of publicly burning 2 utils before
playing this game. We will use \(D\) for Don't burn, and \(B\) for burn.
So actually each player has four choices. Player I has to choose both
\(D\) or \(B\), then \(u\) or \(d\). Player 2 has to choose whether to
play \(l\) or \(r\) in each of the two possibilities: first, when \(D\)
is played, second, when \(B\) is played. We'll write \(lr\) for the
strategy of playing \(l\) if \(D\), and \(r\) if \(B\), and so on for
the other strategies.

\begin{longtable}[]{@{}lrcccc@{}}
\caption{The strategy table for the money burning
game}\label{tbl-bpd-mb}\tabularnewline
\toprule\noalign{}
\endfirsthead
\endhead
\bottomrule\noalign{}
\endlastfoot
& & \multicolumn{4}{c@{}}{%
\textbf{Player II}} \\
& & \(ll\) & \(lr\) & \(rl\) & \(rr\) \\
\multirow{4}{*}{\textbf{Player I}} & \(Du\) & 4, 1 & 4, 1 & 0, 0 & 0,
0 \\
& \(Dd\) & 0, 0 & 0, 0 & 1, 4 & 1, 4 \\
& \(Bu\) & 2, 1 & -2, 0 & 2, 1 & -2, 0 \\
& \(Bd\) & -2, 0 & -1, 4 & -2, 0 & -1, 4 \\
\end{longtable}

Now we can analyse this game a couple of ways. First, we can go through
eliminating weakly dominated strategies in Table~\ref{tbl-bpd-mb}. Note
that \(Du\) strictly dominated \(Bd\), so we can eliminate it. If \(Bd\)
is out, then \(ll\) weakly dominates \(lr\), and \(rl\) weakly dominates
\(rr\). So we can eliminate \(lr\) and \(rr\). Now \(Bu\) weakly
dominates \(Dd\), relative to the remaining options, so we can eliminate
it. Given that just \(Du\) and \(Bu\) remain, \(ll\) weakly dominates
all options for Player II, so it is all that is left. And given that
\(ll\) is all that remains, \(Du\) strongly dominates \(Bu\). So the
iterative deletion of weakly dominated strategies leaves us with
\(\langle Du, ll \rangle\) as the unique solution of the game.

Alternatively, we can think the players reason as follows. (The
following reasoning is sometimes called `forward induction' reasoning.)
Playing \(Bd\) is obviously irrational for Player I, since its maximum
return is -1, and playing \(Du\) or \(Dd\) guarantees getting at least
0. So any rational player who plays \(B\) must be playing \(Bu\). That
is, if Player I is rational and plays \(B\), she will play \(Bu\).
Moreover, this is common knowledge among the players. So if Player I
plays \(B\), she will recognise that Player II will know that she is
playing \(Bu\). And if Player II knows Player I is playing \(u\) after
burning, she will play \(l\), since that returns her 1 rather than 0. So
Player I knows that playing \(B\) leads to the 2,1 state; i.e., it
returns her 2. But now it is irrational to play \(Dd\), since that gets
at most 1, and playing \(B\) leads to a return of 2. So Player I will
play \(Bu\) or \(Du\). And since the reasoning that leads to this is
common knowledge, Player II must play \(ll\), since playing \(l\) is her
best response to a play of \(u\), and she knows Player I will play
\(u\). But if Player II is going to play \(ll\), Player I doesn't need
to burn; she can just play \(Du\).

There is a crucial conditional in the middle of that argument; let's
isolate it.

\begin{itemize}
\tightlist
\item
  If Player I is rational and plays \(B\), she will play \(Bu\).
\end{itemize}

But that's not what is needed to make the argument work. What Player I
needs, and in fact needs to be common knowledge, is the following
subjunctive.

\begin{itemize}
\tightlist
\item
  If Player I is rational then, if she were to play \(B\), she would
  play \(Bu\).
\end{itemize}

But in fact there's no reason to believe that. After all, if the forward
induction argument is right, then it is \emph{irrational} to burn the
money. So if Player I were to play \(B\), i.e., burn the money, then she
would be doing something irrational. And the fact that Player I is
actually rational is consistent with the counterfactual that if she were
to do one irrational thing, it would be because she is following an
irrational strategy. Note that on the assumption that \(Du\) is optimal,
then if Player I finds herself having played \(B\), it isn't clear that
playing \(d\) is irrational; it isn't like it is dominated by \(u\).

So it isn't true that if Player I were to burn the utils, Player II
would know that she is playing \(Bu\), and react accordingly by playing
\(l\). And if that's not true, then there's no reason to think that
\(Bu\) is better than \(Dd\). And if there's no reason to think that,
there's no reason to be confident Player 2 will play \(ll\).

Robert Stalnaker (1998), in his commentary on the game, goes further and
provides a positive model where the players start with a common belief
in rationality, and in what each other will do, but in which the players
play \(Bu\) and \(rl\). I'll leave it as an exercise to work out what
counterfactuals the agents have to believe to make this rational, but
suffice to say that it is a situation where each side has true beliefs
about the other, and maximises expected utility given their beliefs.

There is another, relatively simple, equilibrium~to the game. Player I
will play \(Dd\), and Player II will play \(rr\). Player II is certain
that Player I will play \(d\), and will keep that belief even if Player
I irrationally burns 2 utils to start with. Given that certainty, it is
rational for Player I to play \(Dd\), so Player II's belief is
consistent with believing Player I to be rational. And since Player I is
playing \(Dd\), playing \(rr\) is perfectly rational for Player II;
indeed it results in her best possible outcome. Moreover, given that
Player II is playing \(rr\), it makes sense for Player I to play \(d\)
even if they were, irrationally, to burn the utils. So even though
playing \(Bd\) would be irrational, since it is strictly dominated, if
they were to (irrationally) play \(D\), they \emph{should} follow that
by playing \(d\). There's an important point about the scope of the
claim that playing \(Bd\) is irrational. It \emph{doesn't} imply that if
the player were to irrationally play \(B\), it would be irrational to
follow with a play of \(d\). If Player I was convinced that Player II
was playing \(rr\), then it would be rational to follow \(B\) with
\(d\).

\section{Iterated Prisoners' Dilemma}\label{sec-iterated-pd}

Let's return to Prisoners Dilemma, using \(C\) and \(D\) for cooperate
and defect. We'll call the players I and II, and use uppercase letters
for Player I's strategies, and lower case letters for Player II's
strategies. We'll assume that each player knows that the other players
are perfectly rational in Stalnaker's sense. (That is, they have true
beliefs about the other, and they maximise expected utility given their
beliefs.)

There are a couple of arguments that the players must end up at the
equilibrium~where every player defects on every move. One of these is an
argument by backwards induction.

At the last move of the game, defecting dominates cooperating. Both
players know this. At the second-last move, you might have thought
antecedently that there was some benefit to cooperating. After all, it
might induce cooperation in the other player. But the only benefit one
could get from cooperating was cooperation at the next (i.e., last)
move. And no rational player will cooperate on the last move. So, if
you're playing with a rational player, there's no benefit to cooperating
on the second-last move. But if that's right, and everyone knows it,
then there is no benefit to cooperating on the third-last move. The only
benefit would be if that would induce cooperation, and it couldn't,
since we just proved that any rational player would defect on the
second-last move. And so on for any finite length game that you like.

Alternatively, we could look at the game from a strategic perspective.
As we've seen in games like \textbf{Burning}, sometimes there are Nash
equilibria that don't appear in backwards induction reasoning. But this
isn't the case in finitely iterated Prisoners' Dilemma. The only Nash
equilibrium~is that both players defect in every circumstance.

This is a little tricky to check since the strategic form of iterated
Prisoners' Dilemma gets very complex very quickly. Let's just consider
the three round version of the game. Already each player has 128
strategies to choose from. The choice of a strategy involves making 7
distinct choices:

\begin{enumerate}
\def\labelenumi{\arabic{enumi}.}
\tightlist
\item
  What to do in the first round.
\item
  What to do in the second round if the other player cooperates in the
  first round.
\item
  What to do in the second round if the other player defects in the
  first round.
\item
  What to do in the third round if the other player cooperates in each
  of the first two rounds.
\item
  What to do in the third round if the other player cooperates in the
  first round then defects in the second round.
\item
  What to do in the third round if the other player defects in the first
  round then cooperates in the second round.
\item
  What to do in the third round is the other player defects in each fo
  the first two rounds.
\end{enumerate}

Since these 7 choices are distinct, and the player has 2 choices at each
point, there are \(2^7 = 128\) possible strategies. So the strategic
form of the table involves \(128 \times 128 = 16384\) cells. Needless to
say, we \emph{won't} be putting that table here. (Though it isn't too
hard to get a computer to draw the table for you. The four round game,
where each player has \(2^{15}\) choices, and there are over one billion
cells in the decision table, requires more computing power!)

We'll write a strategy as \(\langle x_1x_2x_3x_4x_5x_6x_7 \rangle\),
where \(x_i\) is 0 if the player's answer to the \(i\)'th question above
is to cooperate, and 1 if it is to defect. So
\(\langle 0000000 \rangle\) is the strategy of always cooperating,
\(\langle 1111111 \rangle\) is the strategy of always defecting,
\(\langle 0010101 \rangle\) is `tit-for-tat', the strategy of
cooperating on the first move, then copying the other player's previous
move, and so on.

Let \(\langle x_1x_2x_3x_4x_5x_6x_7 \rangle\) be any strategy that
doesn't always involve defection on the final round, i.e., a strategy
where \(x_4 + x_5 + x_6 + x_7 < 4\). It is easy enough to verify that
such a strategy is weakly dominated by
\(\langle x_1x_2x_31111 \rangle\). In some plays of the game, the
defection on the final round leads to getting a better outcome. In other
plays of the game, \(\langle x_1x_2x_3x_4x_5x_6x_7 \rangle\) and
\(\langle x_1x_2x_31111 \rangle\) have the same play. For instance,
\(\langle 0001000 \rangle\) will do just as well as
\(\langle 0001111 \rangle\) if the opponent plays any strategy of the
form \(\langle 000x_4x_5x_6x_7 \rangle\), but it can never do better
than \(\langle 0001111 \rangle\). So if each player is perfectly
rational, we can assume their strategy ends with \(1111\).

That cuts each player's choices down to 8. Let's do that table. We'll
use \(C\) and \(c\) rather than 0, and \(D\) and \(d\) rather than 1, so
it is easier to distinguish the two players' strategies. When we label
the table, we'll leave off the trailing \(D\)s and \(d\)'s, since we
assume players are defecting on the last round. We'll also leave off the
3 units the players each get in the last round. (In other words, this
will look a lot like the table for a \emph{two} round iterated
Prisoners' Dilemma, but it is crucial that it is actually a three-round
game.)

{\def\LTcaptype{none} % do not increment counter
\begin{longtable}[]{@{}lrcccccccc@{}}
\toprule\noalign{}
\endhead
\bottomrule\noalign{}
\endlastfoot
& & \multicolumn{8}{c@{}}{%
\textbf{Player II}} \\
& & \(ccc\) & \(ccd\) & \(cdc\) & \(cdd\) & \(dcc\) & \(dcd\) & \(ddc\)
& \(ddd\) \\
\multirow{8}{*}{\textbf{Player I}} & \(CCC\) & 6, 6 & 6, 6 & 3, 8 & 3, 8
& 3, 8 & 3, 8 & 0, 10 & 0, 10 \\
& \(CCD\) & 6, 6 & 6, 6 & 3, 8 & 3, 8 & 5, 5 & 5, 5 & 1, 6 & 1, 6 \\
& \(CDC\) & 6, 6 & 8, 3 & 4, 4 & 4, 4 & 3, 8 & 3, 8 & 0, 10 & 0, 10 \\
& \(CDD\) & 6, 6 & 8, 3 & 4, 4 & 4, 4 & 5, 5 & 5, 5 & 1, 6 & 1, 6 \\
& \(DCC\) & 6, 6 & 5, 5 & 8, 3 & 5, 5 & 4, 4 & 1, 6 & 4, 4 & 1, 6 \\
& \(DCD\) & 6, 6 & 5, 5 & 8, 3 & 5, 5 & 6, 1 & 2, 2 & 6, 1 & 2, 2 \\
& \(DDC\) & 6, 6 & 6, 1 & 10, 0 & 6, 1 & 4, 4 & 1, 6 & 4, 4 & 1, 6 \\
& \(DDD\) & 6, 6 & 6, 1 & 10, 0 & 6, 1 & 6, 1 & 2, 2 & 6, 1 & 2, 2 \\
\end{longtable}
}

In this game \(CCC\) is \emph{strongly} dominated by \(CDD\), and
\(DCC\) is strongly dominated by \(DDD\). Similarly \(ccc\) and \(dcc\)
are strongly dominated. So let's delete them.

{\def\LTcaptype{none} % do not increment counter
\begin{longtable}[]{@{}lrcccccc@{}}
\toprule\noalign{}
\endhead
\bottomrule\noalign{}
\endlastfoot
& & \multicolumn{6}{c@{}}{%
\textbf{Player II}} \\
& & \(ccd\) & \(cdc\) & \(cdd\) & \(dcd\) & \(ddc\) & \(ddd\) \\
\multirow{6}{*}{\textbf{Player I}} & \(CCD\) & 6, 6 & 3, 8 & 3, 8 & 5, 5
& 1, 6 & 1, 6 \\
& \(CDC\) & 8, 3 & 4, 4 & 4, 4 & 3, 8 & 0, 10 & 0, 10 \\
& \(CDD\) & 8, 3 & 4, 4 & 4, 4 & 5, 5 & 1, 6 & 1, 6 \\
& \(DCD\) & 5, 5 & 8, 3 & 5, 5 & 2, 2 & 6, 1 & 2, 2 \\
& \(DDC\) & 6, 1 & 10, 0 & 6, 1 & 1, 6 & 4, 4 & 1, 6 \\
& \(DDD\) & 6, 1 & 10, 0 & 6, 1 & 2, 2 & 6, 1 & 2, 2 \\
\end{longtable}
}

Note that each of these is a best response. In particular,

\begin{itemize}
\tightlist
\item
  \(CCD\) is a best response to \(dcd\).
\item
  \(CDC\) is a best response to \(ccd\).
\item
  \(CDD\) is a best response to \(ccd\) and \(dcd\).
\item
  \(DCD\) is a best response to \(ddc\).
\item
  \(DDC\) is a best response to \(cdc\) and \(cdd\).
\item
  \(DDD\) is a best response to \(cdc\), \(cdd\), \(ddc\) and \(ddd\).
\end{itemize}

Now the only Nash equilibrium~of that table is the bottom-right corner.
But there are plenty of strategies that, in a strategic version of the
game, would be consistent with common belief in perfect rationality. The
players could, for instance, play \(CDC\) and \(cdc\), thinking that the
other players are playing \(ccd\) and \(CCD\) respectively. Those
beliefs would be false, but they wouldn't be signs that the other
players are irrational, or that they are thinking the other players are
irrational.

But you might suspect that the strategic form of the game and the
extensive form are crucially different. The very fact that there are
Nash equilibria that are not subgame perfect equilibria suggests that
there are tighter constraints on what can be played in an extensive game
consistent with rationality and belief in rationality. Stalnaker argues,
however, that this isn't right. In particular, any strategy that is
(given the right beliefs) perfectly rational in the strategic form of
the game is also (given the right beliefs and belief updating
dispositions) perfectly rational in the extensive form. We'll illustrate
this by working more slowly through the argument that the game play
\(\langle CDC, cdc \rangle\) is consistent with common belief in perfect
rationality.

The first thing you might worry about is that it isn't clear that
\(CDC\) is perfectly rational, since it is looks to be weakly dominated
by \(CDD\), and perfect rationality is inconsistent with playing weakly
dominated strategies. But in fact \(CDC\) isn't weakly dominated by
\(CDD\). It's true that on the six columns represented here, \(CDC\)
never does better than \(CDD\). But remember that each of these
strategies is short for a \emph{three}-round strategy; they are really
short for \(CDCDDDD\) and \(CDDDDDD\). And \(CDCDDDD\) is not weakly
dominated by \(CDDDDDD\); it does better, for example, against
\(dcccddd\). That is the strategy is defecting the first round,
cooperating the second, then defecting on the third round unless the
other player has cooperated on each of the first two rounds. Since
\(CDC\) does cooperate each of the first two rounds, it gets the
advantage of defecting against a cooperator on the final round, and ends
up with 8 points, whereas \(CDD\) merely ends up with 6.

But why would we think Player II might play \(dcccddd\)? After all, it
is itself a weakly dominated strategy, and perfectly rational beings
(like Player II) don't play weakly dominated strategies. But we're not
actually thinking Player II \emph{will} play that. Remember, Player I's
assumption is that Player II will play \(ccddddd\), which is not a
weakly dominated strategy. (Indeed, it is tit-for-tat-minus-one, which
is a well-known strategy.) What Player I also thinks is that if she's
wrong about what Player II will play, then it is possible that Player II
is not actually perfectly rational. That's consistent with believing
Player II is actually perfectly rational. The assumption of common
belief in perfect rationality is not sufficient for the assumption that
one should believe that the other player is perfectly rational \emph{no
matter what surprises happen}. Indeed, that assumption is barely
coherent; some assumptions are inconsistent with perfect rationality.

One might be tempted by a weaker assumption. Perhaps a player should
hold on to the belief that the other player is perfectly rational unless
they get evidence that is \emph{inconsistent} with that belief. But it
isn't clear what could motivate that. In general, when we are surprised,
we have to give up something that isn't required by the surprise. If we
antecedently believe \(p \wedge q\), and learn \(\neg (p \wedge q)\),
then what we've learned is inconsistent with neither \(p\) nor \(q\),
but we must give up one of those beliefs. Similarly here, it seems we
must be prepared to give up a belief in the perfect rationality of the
other player in some circumstances when that is not entailed by our
surprise. (Note that even if you don't buy the argument of this
paragraph, there still isn't a reason to think that perfectly rational
players can't play \(CDD\). But analysing that possibility is beyond the
scope of these notes.)

What happens in the extensive form version of the game? Well, each
player cooperates, thinking this will induce cooperation in the other,
then each player is surprised by a defection at round two, then at the
last round they both defect because it is a one-shot Prisoners' Dilemma.
Nothing seems irrational there. We might wonder why neither defected at
round one. Well, if they believed that the other player was playing
tit-for-tat-minus-one, then it is better to cooperate at round one (and
collect 8 points over the first two rounds) than to defect (and collect
at most 6). And playing tit-for-tat-minus-one is rational if the other
person is going to defect on the first and third rounds, and play
tit-for-tat on round two. So as long as Player I thinks that Player II
thinks that Player I thinks that she is going to defect on the first and
third rounds, and play tit-for-tat on round two, then her cooperation at
round one is rational, and consistent with believing that the other
player is rational.

But note that we had to attribute an odd belief to Player II. We had to
assume that Player II is \emph{wrong} about what Player I will play. It
turns out, at least for Prisoners' Dilemma, that this is crucial. If
both players are certain that both players are perfectly rational, and
have no false beliefs, then the only strategy that can be rationalised
is permanent defection. The proof (which I'm leaving out) is in
Stalnaker (1996).

This \emph{isn't} because the no false beliefs principle suffices for
backwards induction reasoning in general. In \textbf{Cooperative
Caterpillar}, we can have a model of the game where both players are
perfectly rational, and have correct beliefs about what the other player
will do, and both those things are common belief, and yet the backwards
induction solution is not played. In that game \(A\) believes, truly,
that \(B\) will play \(d\), and \(B\) believes, truly, that \(A\) will
play \(A_1D_2\). And each of these moves is optimal given (true!)
beliefs about the other player's strategy.

But Prisoners' Dilemma is special. It isn't that in a game played
between players who know each other to have no false beliefs and be
perfectly rational that we must end up at the bottom-right corner of the
strategic table. But we must end up in a game where every player defects
every time. It could be that Player I thinks Player II is playing either
\(dcd\) or \(ddd\), with \(ddd\) being much more, and on that basis she
decides to play \(dcd\). And Player II could have the converse beliefs.
So we'll end up with the play being \(\langle DCD, dcd \rangle\). But of
course that means each player will defect on the first two rounds, and
then again on the last round since they are perfectly rational. In such
a case the requirement that each player have no false beliefs won't even
be enough to get us to Nash equilibrium  since
\(\langle DCD, dcd \rangle\) is not a Nash equilibrium. But it is enough
to get permanent defection.

\bookmarksetup{startatroot}

\chapter{Bayesian Games}\label{sec-bayesian}

\section{Two Kinds of Ignorance}\label{sec-two-kinds-of-ignorance}

In the games of Chapter~\ref{sec-games-and-time}, players were sometimes
ignorant about what had already happened. When a player has to move
without knowing which node of the tree they are at, we represent that
with an information set, and we say the game is one of \textbf{imperfect
information}. Prisoners' Dilemma in extensive form is like this. The
second player moves knowing the rules, knowing the payoffs, knowing
everything about the game, and not knowing only what the first player
just did.

There is a different thing a player might not know. They might not know
what the game is. They might not know what the other player's payoffs
are, and so not know what the other player is trying to achieve. Games
like this are games of \textbf{incomplete information}.

The distinction matters and it is easy to lose, so it is worth putting
the two side by side. In a game of imperfect information you know the
whole payoff table and you don't know where you are in the tree. In a
game of incomplete information you don't know the payoff table.

Almost every interesting application of game theory outside the
classroom involves incomplete information. When you bargain over a used
car you do not know how much the seller values it. When a firm decides
whether to enter a market it does not know the incumbent's costs. When
two states approach a crisis, neither knows how much the other is
willing to lose. Writing down a payoff table and assuming both players
can see it is a serious idealisation.

But we don't obviously have a theory for these cases. Our solution
concepts all take a game as input, and a game is a payoff table. If the
players don't know the payoff table, we cannot even say what game they
are playing, and if we cannot say that, we cannot say what an
equilibrium of it would be. Worse, there is a regress waiting. If Row
doesn't know Col's payoffs, Row has beliefs about them. Col then has
beliefs about Row's beliefs about Col's payoffs. And Row has beliefs
about those. It looks like modelling ignorance about payoffs requires an
infinite hierarchy of beliefs about beliefs, which is not something we
can write on a page.

\section{Types}\label{sec-types}

John Harsanyi (1967) found a clever way through this. Games of
incomplete information can be turned into games of imperfect
information, which we already know how to handle.

The trick is to let the ignorance be about a move rather than about the
game. We add a new player, \textbf{Nature}, who moves first, and whose
move settles what the payoffs are. Nature's move is not a choice; it is
a draw from a probability distribution that is common knowledge. The
players then observe part of what Nature did, and not all of it.

What each player observes about Nature's move is called their
\textbf{type}. A type settles everything about a player that the other
players might be unsure of. In most of our examples the type just is the
player's payoff function, but nothing turns on that. A type could encode
how patient a player is, or how good they are at chess, or what they
believe about the other player. It could encode something completely
distinct; in a scenario where one player can observe whether it is
raining and the other can't, it's sometimes useful to classify the first
player's type as \emph{rain} or \emph{sun}. (If you insist, you could
call their type \emph{sees rain} or \emph{sees sun}, but we often use
the externally individuated types.)

So a \textbf{Bayesian game} has:

\begin{itemize}
\tightlist
\item
  a set of players;
\item
  for each player, a set of actions;
\item
  for each player, a set of possible types;
\item
  a probability distribution over the profiles of types, called the
  \textbf{prior};
\item
  for each player, a payoff function that depends on the actions of
  everyone and on their own type.
\end{itemize}

We assume the prior is common knowledge. This is the \textbf{common
prior assumption}, and without it we need a more complicated theory. One
consequence of it is that the players start out agreeing about how
likely each type profile is, and disagree only because they have seen
different draws. That assumption is not always plausible in the cases
the these tools are applied to, and it's often worth questioning on an
application. For now, notice one useful theoretical simplification that
this assumption gets us. Instead of thinking about a potentially
infinite hierarchy of beleifs about beliefs, we can just use the
conditional probabilities a player gets from the prior, given their own
type, and everyone can compute those.

\section{Bayesian Nash Equilibrium}\label{sec-bne}

Once the game has been transformed this way, we can reuse Nash
equilibrium almost unchanged. There is one adjustment to make first.

In a Bayesian game, a player knows their own type before acting. So a
strategy cannot just be an action; it has to say what the player does
for each type they might turn out to have. A \textbf{strategy} in a
Bayesian game is a function from that player's types to that player's
actions.

This is a common mistake, so it's worth belaboring the point. If Col has
two types, and Col has two moves regardless of their type, then Col has
\emph{four} strategies, not two: do X whatever happens, do Y whatever
happens, do X if the first type and Y if the second, or do Y if the
first and X if the second. Col will only ever perform one action in any
play of the game. But the strategy has to be defined for the type Col
didn't turn out to have, because Row's reasoning depends on what Col
would have done in that case.

A \textbf{Bayesian Nash equilibrium} is then a profile of strategies,
one for each player, such that each type of each player is choosing an
action that maximises their expected payoff, given the strategies of the
others and given the beliefs about the others' types that they get by
conditioning the prior on their own type.

Equivalently, and this is the cleaner way to say it: a Bayesian Nash
equilibrium is a Nash equilibrium of the imperfect information game that
Harsanyi's transformation produces. In a sense, this isn't a new
solution concept; it's a new term for an old concept.

\section{An Example}\label{sec-bayesian-example}

Take Battle of the Sexes from Section~\ref{sec-battle-of-sexes}, and
make one change. Row still wants to meet Col. But Row is no longer sure
that Col wants to meet Row.

Col has two types. With probability \(\nicefrac{1}{2}\) Col is
\emph{sociable}, and wants to meet, with the payoffs we saw before. With
probability \(\nicefrac{1}{2}\) Col is \emph{elusive}, and would rather
be wherever Row is not. Col knows which type Col is. Row does not.

If Col is sociable, the game (shown in Table~\ref{tbl-familiar-bos}) is
the one we already know.

\begin{longtable}[]{@{}lrcc@{}}
\caption{Familiar Battle of the
Sexes}\label{tbl-familiar-bos}\tabularnewline
\toprule\noalign{}
\endfirsthead
\endhead
\bottomrule\noalign{}
\endlastfoot
& & \multicolumn{2}{c@{}}{%
\textbf{Col}} \\
& & X & Y \\
\multirow{2}{*}{\textbf{Row}} & X & 2, 1 & 0, 0 \\
& Y & 0, 0 & 1, 2 \\
\end{longtable}

If Col is elusive, Row's payoffs are unchanged, but Col's change.

\begin{longtable}[]{@{}lrcc@{}}
\caption{A variant of @\#tbl-familiar-bos where Col does not want to
meet}\label{tbl-changed-bos}\tabularnewline
\toprule\noalign{}
\endfirsthead
\endhead
\bottomrule\noalign{}
\endlastfoot
& & \multicolumn{2}{c@{}}{%
\textbf{Col}} \\
& & X & Y \\
\multirow{2}{*}{\textbf{Row}} & X & 2, 0 & 0, 2 \\
& Y & 0, 1 & 1, 0 \\
\end{longtable}

In Table~\ref{tbl-changed-bos}, look first at Col, who has it easy,
because Col knows which table is `real'.

Sociable Col wants to match Row. If Row plays X, sociable Col gets 1
from X and 0 from Y, so plays X. If Row plays Y, sociable Col gets 0
from X and 2 from Y, so plays Y.

Elusive Col wants to mismatch. If Row plays X, elusive Col gets 0 from X
and 2 from Y, so plays Y. If Row plays Y, elusive Col gets 1 from X and
0 from Y, so plays X.

Now Row, who has to average over the two. Suppose Col's strategy is
\emph{play X if sociable, Y if elusive}. Then Col plays X half the time
and Y half the time, from Row's point of view, and Row's expected
payoffs are

\[
E(X) = \frac{1}{2} \times 2 + \frac{1}{2} \times 0 = 1
\]

\[
E(Y) = \frac{1}{2} \times 0 + \frac{1}{2} \times 1 = \frac{1}{2}
\]

So Row plays X. And we should check this is consistent: given that Row
plays X, sociable Col does want to play X, and elusive Col does want to
play Y. Nothing unravels, so we have a Bayesian Nash equilibrium.

Two things about this are worth spelling out in more detail.

The first is that Col's two types do different things. The possibility
of this happening is why the type system is interesting. Externally, Row
is facing one person. But from the inside, from Row's perspective, they
are facing a distribution over opponents, and has to pick an action that
is good on average against it.

The second is that Row's behaviour looks, from the outside, exactly like
behaviour in an ordinary game. Row picks X and sticks with it. The
incomplete information is not visible in what Row does. But it is
visible in why Row does it, and in what Row would have done had the
probabilities been different. Suppose Col were sociable with probability
\(p\) rather than \(\nicefrac{1}{2}\). Then Row's expected payoff from X
is \(2p\) and from Y is \(1-p\), so the equilibrium we just found
survives exactly when \(p \geqslant \frac{1}{3}\). Below that it falls
apart, though something strange starts to happen. There is no pure
strategy equilibrium when \(p\) is below \(frac{1}{3}\). Whatever Row
does, if Col knows it, they will avoid Row, which Row doesn't want. So
Row has to mix, playing X a third of the time, and Row and Col still
find each other sometimes, by luck rather than design.

\section{Purification}\label{sec-purification}

We can now settle a debt from
Section~\ref{sec-what-is-a-mixed-strategy}.

The problem there was that mixed equilibria seem to require players to
randomise, and it is not clear why anyone would. At a mixed equilibrium
every pure strategy in the support does exactly as well as the mixture,
so nothing recommends the mixture over any of them. The only thing a
player's mixture accomplishes is keeping the \emph{other} player
indifferent, which is a strange thing for a self-interested agent to be
aiming at.

Harsanyi's response is that the randomness is really in ignorance of the
payoffs, not in what the other player does (Harsanyi 1973). Here's the
idea. Imagine that we're in a game with a mixed strategy equilibrium and
no pure strategy equilibrium. And each player has a small amount of
private information. They know their own payoffs. For instance, in one
outcome they might know that their payoff is 2, while the other player
only knows that it is \(2 + \delta\), where \(\delta\) is randomly drawn
from \([-\epsilon, \epsilon]\), for arbitrarily small \(\epsilon\). Now
we're playing a Bayesian game; the player's type is the value of
\(\delta\).

In the Bayesian equilibrium of the perturbed game, almost every type has
a strict best reply. A player whose private draw tips them slightly
towards one action simply takes it. Almost no one randomises; almost
every type plays a pure strategy. But an opponent who cannot see the
draw cannot predict the action, and from the opponent's point of view
the player's behaviour is a probability distribution over actions.

Harsanyi's theorem (1973) is that as \(\epsilon\) shrinks to zero, that
of pure strategies converges to the mixed equilibrium of the original
game. The mixed equilibrium is the limit of pure strategy Bayesian Nash
equilibria of nearby games of incomplete information.

If that is right, the interpretive problem dissolves. Nobody was ever
randomising. Each player was doing the strictly best thing their own
private situation called for, and the mixture was never a description of
anyone's decision. It was a description of what the other player could
work out.

Whether it is right is a further question, and there are two obvious
worries. The first is that this all rests on reasoning from a sequence
to what happens in the limit. You might worry that this is illegitimate.
Relatedly, you might worry that we can't tell much about a game by
looking at different games, even if they are neighbors. For what it's
worth, the idea that we should look at nearby games has been really
influential. We're not going to focus on it, but the formal statements
of many of the concepts we'll somewhat informally introduce in the rest
of this chapter uses the same idea.

\section{Perfect Bayesian Equilibrium}\label{sec-pbe}

The core idea behind subgame perfect equilibrium~was that we wanted to
eliminate equilibria that relied on `incredible threats'. That is, there
are some equilibria such that the first player makes a certain move only
because if they make a different move, the second player to move would,
on their current strategy, do something that makes things worse for both
players. But there's no reason to think that the second player would
actually do that.

The same kind of consideration can arise in games where there aren't any
subgames. For instance, consider the following game, which we'll call
\textbf{Perfect}.

\begin{figure}

\centering{

\pandocbounded{\includegraphics[keepaspectratio,alt={The game Perfect}]{index_files/mediabag/figures/svg/perfect.pdf}}

}

\caption{\label{fig-perfect}Perfect}

\end{figure}%

The game here is a little different to one we've seen so far. First
Player 1 makes a move, either \(L\), \(M\) or \(R\). If her move is
\(R\), the game ends, with the (2, 2) payout. If she moves \(L\) or
\(M\), the game continues with a move by Player 2. But crucially, Player
2 does not know what move Player 1 made, so she does not know which node
she is at. So there isn't a subgame here to start. Player 2 chooses
\(l\) or \(r\), and then the game ends.

There are two Nash equilibria to \textbf{Perfect}. The obvious
equilibrium~is \(Ll\). In that equilibrium  Player 1 gets her maximal
payout, and Player 2 gets as much as she can get given that the
rightmost node is unobtainable. But there is another Nash equilibrium
available, namely \(Rr\). In that equilibrium, Player 2 gets her maximal
payout, and Player 1 gets the most she can get given that Player 2 is
playing \(r\). But note what a strange equilibrium~it is. It relies on
the idea that Player 2 would play \(r\) were she to be given a move. But
that is absurd. Once she gets a move, she has a straight choice between
1, if she plays \(l\), and 0, if she plays \(r\). So obviously she'll
play \(l\).

This is just like the examples that we used to motivate subgame perfect
equilibrium, but that concept doesn't help us here. So we need a new
concept. The core idea is that each player should be modelled as a
Bayesian expected utility maximiser. More formally, the following
constraints are put on players.

\begin{enumerate}
\def\labelenumi{\arabic{enumi}.}
\tightlist
\item
  At each point in the game, each player has a probability distribution
  over where they are in the game. These probability distributions are
  correct about the other players' actions. That is, if a player is
  playing a strategy \(S\), everyone has probability 1 that they are
  playing \(S\). If \(S\) is a mixed strategy, this might involve having
  probabilities between 0 and 1 in propositions about which move the
  other player will make, but players have correct beliefs about other
  players' strategies.
\item
  No matter which node is reach, each player is disposed to maximise
  expected utility on arrival at that node.
\item
  When a player had a positive probability of arriving at a node, on
  reaching that node they update by conditionalisation.
\item
  When a player gave 0 probability to reaching a node (e.g., because the
  equilibrium~being played did not include that node), they have some
  disposition or other to form a set of consistent beliefs at that node.
\end{enumerate}

The last constraint is very weak, but it does enough to eliminate the
equilibrium~\(Rr\). The constraint implies that when Player 2 moves, she
must have some probability distribution \(\Pr\) such that there's an
\(x\) such that \(\Pr(L) = x\) and \(\Pr(M) = 1-x\). Whatever value
\(x\) takes, the expected utility of \(l\) given \(\Pr\) is 1, and the
expected utility of \(r\) is 0. So being disposed to play \(r\) violates
the second condition. So \(Rr\) is not a Perfect Bayesian equilibrium.

It's true, but I'm not going to prove it, that all Perfect Bayesian
equilibria are Nash equilibria. It's also true that the converse does
not hold, and this we have proved; \textbf{Perfect} is an example.

\section{Signaling Games}\label{sec-signaling-games}

Concepts like Perfect Bayesian equilibrium~are useful for the broad
class of games known as signaling games. In a signaling game, Player 1
gets some information that is hidden from player 2. Many applications of
these games involve the information being \emph{de se} information, so
the information that Player 1 gets is referred to as her \emph{type}.
But that's inessential; what is essential is that only one player gets
this information. Player 1 then has a choice of move to make. There is a
small loss of generality here, but we'll restrict our attention to games
where Player 1's choices are independent of her type, i.e., of the
information she receives. Player 2 sees Player 1's move (but not,
remember, her type) and then has a choice of her own to make. Again with
a small loss of generality, we'll restrict attention to cases where
Player 2's available moves are independent of what Player 1 does.

We'll start, in game \textbf{Simple Signal} with a signaling game where
the parties' interests are perfectly aligned.

\begin{figure}

\centering{

\pandocbounded{\includegraphics[keepaspectratio,alt={A simple signalling game}]{index_files/mediabag/figures/svg/simple-signal.pdf}}

}

\caption{\label{fig-simple-signal}Simple Signal}

\end{figure}%

As above, we use an empty disc to signal the initial node of the game
tree. In this case, it is the node in the centre of the tree. The first
move is made by Nature, again denoted as \(N\). Nature assigns a type to
Player 1; that is, she makes some proposition true, and reveals it to
Player 1. Call that proposition \(q\). We'll say that Nature moves left
if \(q\) is true, and right otherwise. We assume the probability (in
some good sense of probability) of \(q\) is \(p\), and this is known
before the start of the game. After Nature moves, Player 1 has to choose
\(U\)p or \(D\)own. Player 2 is shown Player 1's choice, but not
Nature's move. That's the effect of the horizontal dashed lines. If the
game reaches one of the upper nodes, Player 2 doesn't know which one it
is, and if it reaches one of the lower nodes, again Player 2 doesn't
know which it is. Then Player 2 has a make a choice, here simply denoted
as \(l\)eft or \(r\)ight.

In any game of this form, each player has a choice of four strategies.
Player 1 has to choose what to do if \(q\) is true, and what to do if
\(q\) is false. We'll write \(\alpha \beta\) for the strategy of doing
\(\alpha\) if \(q\), and \(\beta\) if \(\neg q\). Since \(\alpha\) could
be \(U\) or \(D\), and \(\beta\) could be \(U\) or \(D\), there are four
possible strategies. Player 2 has to choose what to do if \(U\) is
played, and what to do if \(D\) is played. We'll write \(\gamma \delta\)
for the strategy of doing \(\gamma\) if \(U\) is played, and \(\delta\)
if \(D\) is played. Again, there are four possible choices. (If Player 2
knew what move Nature had made, there would be four degrees of freedom
in her strategy choice, so she'd have 16 possible strategies.
Fortunately, she doesn't have that much flexibility!)

Any game of this broad form is a signaling game. signaling games differ
in (a) the interpretation of the moves, and (b) the payoffs.
\textbf{Simple Signal} has a very simple payoff structure. Both players
get 1 if Player 2 moves \(l\) iff \(q\), and 0 otherwise. If we think of
\(l\) as the formation of a belief that \(q\), and \(r\) as the
formation of the opposite belief, this becomes a simple communication
game. The players get a payoff iff Player 2 ends up with a true belief
about \(q\), so Player 1 is trying to communicate to Player 2 whether
\(q\) is true. This kind of simple communication game was used by David
Lewis (1969) to show that game theoretic approaches could be fruitful in
the study of meaning. The game is perfectly symmetric if
\(p = \nicefrac{1}{2}\); so as to introduce some asymmetries, I'll work
through the case where \(p = \nicefrac{3}{5}\).

\textbf{Simple Signal} has a dizzying array of Nash equilibria, even
given this asymmetry introducing assumption. They include the following:

\begin{itemize}
\tightlist
\item
  There are two \textbf{separating} equilibria, where what Player 1 does
  depends on what Nature does. These are \(\langle UD, lr \rangle\) and
  \(\langle DU, rl \rangle\). These are rather nice equilibria; both
  players are guaranteed to get their maximal payout.
\item
  There are two \textbf{pooling} equilibria, where what Player 1 does is
  independent of what Nature does. These are \(\langle UU, ll \rangle\)
  and \(\langle DD, ll \rangle\). Given that she gets no information
  from Player 1, Player 2 may as well guess. And since
  \(\Pr(q) > \nicefrac{1}{2}\), she should guess that \(q\) is true;
  i.e., she should play \(l\). And given that Player 2 is going to
  guess, Player 1 may as well play anything. So these are also
  equilibria.
\item
  And there are some \textbf{babbling} equilibria. For instance, there
  is the equilibrium~where Player 1 plays either \(UU\) or \(DD\) with
  some probability \(r\), and Player 2 plays \(ll\).
\end{itemize}

Unfortunately, these are all Perfect Bayesian equilibria too. For the
separating and babbling equilibria, it is easy to see that
conditionalising on what Player 1 plays leads to Player 2 maximising
expected utility by playing her part of the equilibrium. And for the
pooling equilibria, as long as the probability of \(q\) stays above
\(\nicefrac{1}{2}\) in any `off-the-equilibrium-path' play (e.g.,
conditional on \(D\) in the \(\langle UU, ll \rangle\) equilibrium),
Player 2 maximises expected utility at every node.

That doesn't mean we have nothing to say. For one thing, the separating
equilibria are Pareto-dominant in the sense that both players do better
on those equilibria than they do on any other. So that's a
non-coincidental reason to think that they will be the equilibria that
arise. There are other refinements on Perfect Bayesian equilibria that
are more narrowly tailored to signaling games. We'll introduce them by
looking at a couple of famous signaling games.

\section{Signaling without
Cooperation}\label{sec-signaling-without-cooperation}

Economists have been interested for several decades in game-theoretic
representations of college that give make going to college essentially a
signal in a game. For example, consider the following variant of the
signaling game, \textbf{College Signal}. It has the following intended
interpretation:

\begin{itemize}
\tightlist
\item
  Player 1 is a student and potential worker, Player 2 is an employer.
\item
  The student is either type high or low\footnote{For some purposes we
    interpret `high' as high intelligence, other times as `high' labor
    market value, other times as something else. I'll use use `high' and
    `low', and leave off what this might mean until we get to the
    question of whether this is really a model of anything in the real
    world} with probability \(p\) of being high. Nature reveals the type
  to the \emph{student}, but only the probability to the employer, so
  \(p\) is that the student is high.
\item
  The student has the choice of going to college (\(U\)) or the beach
  (\(D\)).
\item
  The employer has the choice of hiring the student (\(l\)) or rejecting
  them (\(r\)).
\end{itemize}

In \textbf{College Signal} we make the following extra assumptions about
the payoffs.

\begin{itemize}
\tightlist
\item
  Being hired is worth 4 to the student.
\item
  Going to college rather than the beach costs the high student 1, and
  the low student 5, since college is much harder work for low types.
\item
  The employer gets no benefit from hiring college students as such.
\item
  Hiring a high student pays the employer 1, hiring a low student costs
  the employer 1, and rejections have no cost or benefit.
\end{itemize}

The resulting game tree is shown in Figure~\ref{fig-college-signal}. The
first move, by nature, sends high types to the left, and low types to
the right. In this game, low students never prefer to go to college,
since even the lure of a job doesn't make up for the pain of actually
having to study. So a rational strategy for Player 1 will never be of
the form \(\alpha U\), since for low students, college is a dominated
option, being dominated by \(\alpha D\). But whether high students
should go to college is a trickier question. That is, it is trickier to
say whether the right strategy for Player 1 is \(UD\) or \(DD\), which
are the two strategies consistent with eliminating the strictly
dominated strategies. (Remember that strictly dominated strategies
cannot be part of a Nash equilibrium.)

\begin{figure}

\centering{

\pandocbounded{\includegraphics[keepaspectratio,alt={The college signalling game}]{index_files/mediabag/figures/svg/college-signal.pdf}}

}

\caption{\label{fig-college-signal}College Signal}

\end{figure}%

First, consider the case where \(p < \nicefrac{1}{2}\). In that case, if
Player 1 plays \(DD\), then Player 2 gets \(2p-1\) from playing either
\(ll\) or \(rl\), and 0 from playing \(lr\) or \(rr\). So she should
play one of the latter two options. If she plays \(rr\), then \(DD\) is
a best response, since when employers aren't going to hire, students
prefer to go to the beach. So there is one \textbf{pooling} equilibrium,
namely \(\langle DD, rr \rangle\). But what if Player 2 plays \(lr\)?
Then Player 1's best response is \(UD\), since high students prefer
college conditional on it leading to a job. So there is also a
\textbf{separating} equilibrium, namely \(\langle UD, lr \rangle\). The
employer prefers that equilibrium, since her payoff is now \(p\) rather
than 0. And students prefer it, since their payoff is \(3p\) rather than
0. So if we assume that people will end up at Pareto-dominant outcomes,
we have reason to think that high students, and only high students, will
go to college, and employers will hire college students, and only
college students. And all this is true despite there being no advantage
whatsoever to going to college in terms of how good an employee one will
be.

Especially in popular treatments of the case, the existence of this kind
of model can be used to motivate the idea that college \emph{only} plays
a signaling role. That is, some people argue that in the real world
college does not make students more economically valuable, and the fact
that college graduates have better employment rates, and better pay, can
be explained by the signaling function of college. We'll look in a bit
at the detailed evidence for this, but first lets look at some variants
of the model.

So far we assumed that \(p < \nicefrac{1}{2}\). If we drop that
assumption, and assume instead that \(p > \nicefrac{1}{2}\), the case
becomes more complicated. Now if Player 1 plays \(DD\), i.e., goes to
the beach no matter what, Player 2's best response is still to hire
them. But note that now a very odd equilibrium~becomes available. The
pair \(\langle DD, rl \rangle\) is a Nash equilibrium, and, with the
right assumptions, a Perfect Bayesian equilibrium. This pair says that
Player 1 goes to the beach whatever her type, and Player 2 hires only
beach goers.

This is a very odd strategy for Player 2 to adopt, but it is a little
hard to say just why it is odd. It is clearly a Nash equilibrium. Given
that Player 2 is playing \(rl\), then clearly beach-going dominates
college-going. And given that Player 1 is playing \(DD\), playing \(rl\)
gets as good a return as is available to Player 2, i.e., \(2p-1\). Could
it also be a Perfect Bayesian equilibrium? It could, provided Player 2
has a rather odd update policy. Say that Player 2 thinks that if someone
goes to college, they are low type with probability greater than
\(\nicefrac{1}{2}\). That's consistent with what we've said; given that
Player 1 is playing \(DD\), the probability that a student goes to
college is 0. So the conditional probability that a college-goer is high
type is left open, and can be anything one likes in Perfect Bayesian
equilibrium. So if Player 2 sets it to be, say, 0, then the rational
reaction is to play \(rl\).

But now note what an odd update strategy this is for Player 2. She has
to assume that if someone deviates from the \(DD\) strategy, it is
someone for whom the deviation is strictly dominated. A few economists
and game theorists have thought that we can put more substantive
constraints on probabilities conditional on `off-the-equilibrium-path'
behaviour. One such constraint, is, roughly, that deviation shouldn't
lead to playing a dominated strategy. This is the ``\textbf{intuitive
criterion}'' of Cho and Kreps (1987). In this game, all the criterion
rules out is the odd pair \(\langle DD, rl \rangle\). It doesn't rule
out the very similiar pair \(\langle DD, ll \rangle\). But the intuitive
criterion makes more substantial constraints in other games. We'll close
the discussion of the math of signaling games with such an example, and
a more careful statement of the criterion.

\section{The Intuitive Criterion}\label{sec-intuitive-criterion}

The tree in Figure~\ref{fig-canadian-signal} represents \textbf{Canadian
Signal}, which is also a guessing game. The usual statement of it
involves all sorts of unfortunate stereotypes about the type of men who
have quiche and/or beer for breakfast, and I'm underwhelmed by it. So
I'll run with an example that relies on different stereotypes.

A North American tourist is in a bar. 60\% of the North Americans who
pass through that bar are from Canada, the other 40\% are from the US.
(This is a little implausible for most parts of the world, but maybe it
is very cold climate bar. In Cho and Kreps' version the split is 90/10
not 60/40, but all that matters is which side of 50/50 it is.) The
tourist, call her Player 1, knows her nationality, although the barman
doesn't. The tourist can ask for the bar TV to be turned to hockey or to
baseball, and she knows once she does that the barman will guess at her
nationality. (The barman might also try to rely on her accent, but she
has a fairly neutral upper-Midwest/central Canada accent.) Here are the
tourist's preferences.

\begin{itemize}
\tightlist
\item
  If she is American, she has a (weak) preference for watching baseball
  rather than hockey.
\item
  If she is Canadian, she has a (weak) preference for watching hockey
  rather than baseball.
\item
  Either way, she has a strong preference for being thought of as
  Canadian rather than American. This preference is considerably
  stronger than her preference over which sport to watch.
\end{itemize}

The barman's preferences are simpler; he prefers to make true guesses to
false guesses about the tourist's nationality. All of this is common
knowledge. So the decision tree is given in
Figure~\ref{fig-canadian-signal}. In this tree, \(B\) means asking for
baseball, \(H\) means asking for hockey, \(a\) means guessing the
tourist is from the USA, \(c\) means guessing she is from Canada.

\begin{figure}

\centering{

\pandocbounded{\includegraphics[keepaspectratio,alt={The Canadian signalling game}]{index_files/mediabag/figures/svg/canadian-signal.pdf}}

}

\caption{\label{fig-canadian-signal}Canadian Signal}

\end{figure}%

There are two Nash equilibria for this game. One is
\(\langle HH, ca \rangle\); everyone asks for hockey, and the barman
guesses Canadian if hockey, American if baseball. It isn't too hard to
check this is an equilibrium. The Canadian gets her best outcome, so it
is clearly an equilibrium~for her. The American gets the second-best
outcome, but asking for baseball would lead to a worse outcome. And
given that everyone asks for hockey, the best the barman can do is go
with the prior probabilities, and that means guessing Canadian. It is
easy to see how to extend this to a perfect Bayesian equilibrium; simply
posit that conditional on baseball being asked for, the probability that
the tourist is American is greater than \(\nicefrac{1}{2}\).

The other equilibrium~is rather odder. It is \(\langle BB, ac \rangle\);
everyone asks for baseball, and the barman guesses American if hockey,
Canadian if baseball. Again, it isn't too hard to see how it is Nash
equilibrium. The Canadian would rather have hockey, but not at the cost
of being thought American, so she has no incentive to defect. The
American gets her best outcome, so she has no incentive to defect. And
the barman does as well as he can given that he gets no information out
of the request, since everyone requests baseball.

Surprisingly, this could also be a perfect Bayesian equilibrium. This
requires that the barman maximise utility at every node. The tricky
thing is to ensure he maxmises utility at the node where hockey is
chosen. This can be done provided that conditional on hockey being
chosen, the probability of American rises to above \(\nicefrac{1}{2}\).
Well, nothing we have said rules that out, so there \emph{exists} a
perfect Bayesian equilibrium. But it doesn't seem very plausible. Why
would the barman adopt just \emph{this} updating disposition?

There were a bunch of suggestions for formal rules to rule out
intuitively crazy updating dispositions like this. I'm just going to
note one of the very early attempts in this area, Cho and Kreps'
\emph{Intuitive Criterion}. (Cho and Kreps 1987, 202)

Let an outcome \(o\) of the game be any combination of moves by nature
and the players. Call the players' moves collectively \(P\) and nature's
move \(N\). So an outcome is a pair \(\langle P, N \rangle\). Divide the
players into three categories.

\begin{itemize}
\tightlist
\item
  The \textbf{happy} players in \(o\) are those such that given \(N\),
  \(o\) maximises their possible returns. That is, for all possible
  \(P^\prime\), their payoff in \(\langle P, N \rangle\) is at least as
  large as their payoff in \(\langle P^\prime, N \rangle\).
\item
  The \textbf{content} players are those who are not happy, but who are
  such that for no strategy \(s^\prime\) which is an alternative to
  their current strategy \(s\), would they do better given what nature
  and the other players would do if they played \(s^\prime\).
\item
  The \textbf{unhappy} players are those who are neither happy nor
  content.
\end{itemize}

The standard Nash equilibrium~condition is that no player is unhappy. So
assume we have a Nash equilibrium  with only happy and content players.
And assume it is also a Perfect Bayesian equilibrium. That is, for any
`off-equilibrium' outcome, each player would have some credence function
such that their play maximises utility given that function.

Now add one constraint to those credence functions:

\begin{quote}
\textbf{Intuitive Criterion - Weak Version}

Assume a player has probability 1 that a strategy combination \(P\) will
be played. Consider their credence function conditional on another
player playing \(s^\prime\), wh\-ich is different to the strategy \(s\)
they play in \(P\). If it is consistent with everything else the player
believes that \(s^\prime\) could be played by a player who is content
with the actual outcome \(\langle P, N \rangle\), then give probability
1 to \(s^\prime\) being played by a content player.
\end{quote}

That is, if some deviation from equilibrium~happens, give probability 1
to it being one of the content players, not one of the happy players,
who makes the deviation.

That's enough to rule out the odd equilibrium~for the baseball-hockey
game. The only way that is a Perfect Bayesian equilibrium~is if the
barman responds to a hockey request by \emph{increasing} the probability
that the tourist is American. But in the odd equilibrium, the American
is happy, and the Canadian is merely content. So the Intuitive Criterion
says that the barman should give credence 1 to the hockey requester,
i.e., the deviator from equilibrium, being Canadian. And if so, the
barman won't respond to a hockey request by guessing the requestor is
American, so the Canadian would prefer to request hockey. And that means
the outcome is no longer an equilibrium, since the Canadian is no longer
even content with requesting baseball.

In games where everyone has only two options, the weak version I've
given is equivalent to what Cho and Kreps offers. The official version
is a little more complicated. First some terminology of mine, then their
terminology next.

\begin{itemize}
\tightlist
\item
  A player is \textbf{happy to have played} \(s\) rather than
  \(s^\prime\) if the payoff for \(s\) in \(o\) is greater than any
  possible outcome for \(s^\prime\) given \(N\).
\item
  A player is \textbf{content to have played} \(s\) rather than
  \(s^\prime\) if they are not happy to have played \(s\) rather than
  \(s^\prime\), but the payoff for \(s\) in \(o\) is as great as the
  payoff for \(s^\prime\) given \(N\) and the other players' strategies.
\item
  A player is \textbf{unhappy to have played} \(s\) rather than
  \(s^\prime\) if they are neither happy nor content to have played
  \(s\) rather than \(s^\prime\).
\end{itemize}

If a player is happy to have played \(s\) rather than \(s^\prime\), and
\(s\) is part of some equilibrium~outcome \(o\), then Cho and Kreps say
that \(s^\prime\) is an \textbf{equilibrium-dominated} strategy. The
full version of the Intuitive Criterion is then:

\begin{quote}
\textbf{Intuitive Criterion - Full Version}

Assume a player has probability 1 that a strategy combination \(P\) will
be played. Consider their credence function conditional on another
player playing \(s^\prime\), wh\-ich is different to the strategy \(s\)
they play in \(P\). If it is consistent with everything else the player
believes that \(s^\prime\) could be played by a player who is merely
content to have played \(s\) rather than \(s^\prime\), then give
probability 1 to \(s^\prime\) being played by someone who is content to
have played \(s\), rather than someone who is happy to have played
\(s\).
\end{quote}

This refinement matters in games where players have three or more
options. It might be that a player's options are \(s_1, s_2\) and
\(s_3\), and their type is \(t_1\) or \(t_2\). In the Nash
equilibrium~in question, they play \(s_1\). If they are type \(t_1\),
they are happy to have played \(s_1\) rather than \(s_2\), but merely
content to have played \(s_1\) rather than \(s_3\). If they are type
\(t_2\), they are happy to have played \(s_1\) rather than \(s_3\), but
merely content to have played \(s_1\) rather than \(s_2\). In my
terminology above, both players are content rather than happy with the
outcome. So the weak version of the Intuitive Criterion wouldn't put any
restrictions on what we can say about them conditional on them playing,
say \(s_2\). But the full version does say something; it says other
players should assign probability 1 to the player being type \(t_2\)
conditional on them playing \(s_2\), since the alternative is that
\(s_2\) is played by a player who is happy they played \(s_1\) rather
than \(s_2\). Similarly, it says that conditional on the player playing
\(s_3\), other players should assign probability 1 to their being of
type \(t_1\), since the alternative is to assign positive probability to
a player deviating to a strategy they are happy not to play.

The Intuitive Criterion has been the subject of an enormous literature.
Goo\-gle Scholar lists over 5000 citations for the Cho and Kreps paper
alone, and similar rules are discussed in other papers. So there are
arguments that it is too weak and arguments too strong, and refinements
in all sorts of directions. But we'll leave those. What I most wanted to
stress was the \emph{form} a refinement of Perfect Bayesian
equilibrium~would take. It is a constraint on the conditional credences
of agents conditional on some probability 0 event happening. It's
interesting that there are some plausible, even intuitive constraints;
in general it seems the economic literature has investigated rationality
under 0 probability evidence more than philosophers have.

\section{Spence and the Value of a Useless Signal}\label{sec-spence}

Let's turn away from these very theoretical questions, and back to
applications of these models to real world problems.

\textbf{College Signal} has an important history. It is a simplified
version of the model Spence (1973) published as an account of the labor
market, and which won him a share of a Nobel Prize. It triggered a huge
debate, going to this day, about the economic significance of college.
The debate is not about the mathematics of the model, which all check
out. It's about whether the model is an accurate representation of the
real world.

Let's look at those assumptions. In the separating equilibrium, high
type students go to college and low type ones do not, employers hire
college-goers and reject the rest, and everyone is behaving perfectly
rationally. In the model, \emph{no one learns anything} in college. A
high type student is more valuable to employers, but they would have
been just as valuable with no college. Indeed, them being hired without
going to college would be a weak Pareto-improvement, since college is
costly to the student. The student's productivity is settled by their
type, which Nature fixed before the game started. College does not
change it. All college does is cost less to a high type student than to
a low type one.

That asymmetry in cost is why the model works. Because the high type
student finds college `cheaper', and we'll want to come back to what
that means, there is a range of prices at which the high type student is
willing to pay and the low type one is not, and paying it therefore
reveals the type. The signal works because it is expensive, and because
it is differentially expensive. It does not need to be useful.

Notice also how different this is from \textbf{Simple Signal}. There the
players' interests were aligned, and a costless signal did the job. You
can get the separating equilibrium in a case where it doesn't cost
anything to leave one or the other signal.\footnote{The canonical
  example of this signalling is in David Lewis's book \emph{Convention},
  which grew out of his PhD Thesis (Lewis 1969). The patriots in Boston
  want to warn their soldiers out at Lexington which way the English are
  attacking. The sexton in the Old North Church can see what the English
  are doing, but he can't be seen talking to other patriots. So they
  have a plan; he'll put out one lantern if the English are marching up
  the highway, and two if they are sailing up the river. This works,
  even though it doesn't cost more to put up two lanterns.} In this
case, where the interests diverge, a costless signal couldn't work. If
college was costless, everyone would go. If it cost the \emph{same
amount} to everyone, it wouldn't work.

So we might worry what this says about what we do here. Is college just
a sorting mechanism? Does it produce a private gain but no social gain?
Is it a massively expensive system that just makes us collectively worse
off?

Now note that the model only gives us a consistency proof. We know, and
we'll go over this evidence more in Section~\ref{sec-wage-premium}, that
college grads earn more than non grads. What Spence's model shows is
that this fact alone is consistent with college not actually teaching
anyone anything. But lots of things are consistent without being all
true. To do more, we need to look more at the details.

\section{The College Wage Premium}\label{sec-wage-premium}

There is a key fact that needs explaining. In almost every rich country,
people who finish a degree earn substantially more over a lifetime than
people who do not. According to recent OECD (2025) figures (see
especially Table A4.1), on average across a range of countries, degree
holders earn 40\% more over their lifetime than people with high school
degrees who have not completed college. In the United States, this
premium is even more. Here, grads earn 64\% more than people whose
highest qualification is either a high school degree \emph{or a
diploma}. That's a gap of well over a million dollars, probably closer
to two million dollars. These are enormous numbers, and they need
explaining.

There are three natural explanations which are worth distinguishing.

\begin{itemize}
\tightlist
\item
  \textbf{Human capital.} College teaches things that are worth money.
  The premium measures the skills. This explanation doesn't say
  \emph{where} college teaches these skills; as far as this explanation
  goes it could be that you learn more from your fellow students in the
  dorms than in class.\footnote{Prima facie this is plausible, but note
    that most students around the world don't live in dorms.} What's
  central to this story is that people come out of college with more
  value to employers than they came in with.
\item
  \textbf{Selection.} Colleges admit people who were already going to do
  well. Just knowing that one group of students had higher SAT scores on
  average would suggest they'd earn more money. College grads do, on
  average, have much higher SAT scores than the general populace. Maybe
  that's the whole story.
\item
  \textbf{Signalling.} College is a hurdle that is cheaper to clear for
  the people employers want. The premium is a consequence of this fact.
\end{itemize}

These are not exclusive, and the full explanation is likely to involve
all three. But they are different, and a story where one of them
explains 80\% of the premium, and the other two explain 10\% each leaves
us with a very different picture than if a different one was the 80\%.

Still, let's see if we can make progress by thinking about what the
world would look like if one of them were the whole story. Two of the
explanations look fairly bad by that score.

If it were all human capital, then partial degrees should earn partial
returns, and three years of a four-year degree should be worth roughly
three quarters of the premium. If it were all selection, an admission
letter should be worth almost as much as a diploma, since the letter is
where the sorting happened.

Neither prediction holds. Two or three years of college buys very
little, and the premium mostly arrives when the degree does. This is the
\textbf{sheepskin effect}, named for the old practice of printing
diplomas on sheepskin. This is the single strongest piece of evidence
for the signalling story. It is hard to see why the fourth year of skill
acquisition should be worth several times the third. It is easy to see
why clearing a hurdle should be worth more than nearly clearing it.

The effect is easiest to see in the United States, which has an
unusually large population of people with some college and no degree. In
most countries that group is too small to say much about. In the
American data there are millions of them, and their advantage over high
school graduates is small. While it would definitely be nicer to have
more data, we shouldn't dismiss this American data: it does seem to tell
against the purest version of the human capital story.

\section{Does the Signalling Model Fit?}\label{sec-does-signalling}

So that's the first bit of evidence that the signalling theory is right.
It predicts the sheepskin effect, the other two, \emph{if taken as
comprehensive theories}, predict it won't exist. In fact, it does exist.

The model also explains why the different in wage premia between degrees
is relatively small. So the model fits a world where a degree in
medieval epistemology and a degree in accounting both open doors. We
have to be careful, because there is \emph{some} difference in premia
difference between degrees, but overall I'd say there's less than you'd
expect on a pure capital model.\footnote{Or maybe medieval epistemology
  has more labor market value than I realise. Disclosure: I've never
  taught a \emph{course} on medieval epistemology, but I have done
  lectures on it in epistemology classes.}

The signalling model also explains something else about public
universities like this one: the doors to the lecture theatres are
unlocked. Anyone can walk in to any lecture.\footnote{Some of you may
  recall this fact plays an important part in the plot of \emph{On the
  Calculation of Volume} (Balle {[}2021{]} 2025).} We are, on the whole,
not overrun with people taking advantage of all these effectively free
lectures. This is a little surprising on the human capital
story.\footnote{Though a human capital story that focussed on what you
  learn in discussion sections and in discussions with fellow students
  will also predict this, and is frankly a little more plausible.}

But there are also some difficulties.

\textbf{The premium grows with age.} Again, the Among workers in their
late twenties the advantage is real but moderate; among workers in their
late forties it is substantially larger. That feels like the wrong way
round for a signal. The whole apparatus rests on employers being unable
to see a worker's type directly, and after twenty years of performance
reviews, promotions and colleagues' opinions, the type is not hidden any
more. The signal should be losing value, and it is gaining. The key
thing is that the signalling model turns on the presence or absence of a
degree being a key part of how employers judge the quality of a
would-be-employee, and that's just not plausible after many years of
work.

There are things the defender of the signalling model could say. The
premium at forty-five may reflect twenty years of sorting inside firms
rather than the degree itself, in which case the degree got someone
through the first door and the rest is career. Or the wage structure may
have been fixed early, by contracts or internal pay scales, and simply
persist. Both are defensible, but we're back in a situation where one
theory is explaining things away, and the other theory is readily
predicting them.

\textbf{The cost assumption may not hold.} The model needs college to be
more costly for the low type than the high type, and it needs the
correlation to be tight; this is the single condition that makes the
separating equilibrium work at all. It is worth asking whether it is
true.

The `cost' here mostly isn't literal dollar cost. Some people get more
scholarships for academic success, and some people have to repeat
courses they can't complete. But none of these are making college
literally two million dollars more expensive for low type students.

Rather, the point is that college is more \emph{pleasant} for high types
than low types. OK, who actually finds college fun, who pays the least
psychic cost for being here? The students who flourish, who enjoy the
reading and the seminars and are untroubled by exams, are not obviously
the students who will be the most valuable employees. The traits that
make academic life comfortable are not the traits that make someone
reliable under a deadline, or good with clients, or sound in practical
judgement. When \emph{I} was a student, the students who were paying the
least psychic cost were the perpetual student types, and I would be very
hesitant of having hired them. More generally, one can be an excellent
at many jobs that require college degrees while having found college
genuinely painful, and one can be very happy indeed in a library and a
poor fit for most workplaces.

This is a more serious objection than the first, because it attacks an
assumption rather than a prediction. If there isn't a correlation
between the cost of degree and the value to the employer, the model
breaks. Plenty of low types, from the perspective of employers, will
sail through college because they love it here.

\textbf{College isn't expensive enough.} The wage premium is
\emph{huge}. In the model, everyone is acting rationally. Some people
ratioanlly turn down millions of dollars because the unpleasantness of
those calculus classes is just too much. This beggars belief. In the
modern university, with almost all courses being electives, not all
courses will be that painful.

There is a technical point here. There is a variant on Spence's model
where the difference between high and low is not how much college
\emph{costs}, but whether it is an option at all. In the biological
literature, signalling models often start with a costly signal that
other animals \emph{cannot} send. That's still a signalling model, and a
model where low types are simply incapable of finishing a degree, rather
than able but unwilling, does not require thinking that there are all
these people who will literally not take a calculus class even if
they'll be paid a million dollars for it.

\textbf{Changes in degree length provide natural experiments.} If the
hurdle is what matters, then changing its height should not change the
premium \emph{much}, because employers will recalibrate and a three-year
degree will come to mean what a four-year degree used to mean. If skills
are what matter, cutting the coursework should cut the premium.

Some of the evidence favours the signalling reading. In some countries
bachelor's degrees take three years, and in other countries, like the
US, they take four. On the whole, the premium doesn't obviously reflect
this length difference. Though the US, which has four year degrees, does
have a larger than average premium. In the EU, the Bologna reforms
shortened degrees across much of the continent without eliminating the
graduate premium, though the master's degree has become more expected
since, which looks like the hurdle moving rather than disappearing.

But the cleanest social data we have points towards the human capital
theory. Carolina Arteaga (2018) studied a reform at Universidad de los
Andes in Colombia, which in 2006 cut the required coursework in
economics by about 20\% and in business by about 14\%. The entering
classes, their test scores and the graduation rates stayed put, and so
did the university, the degree title and the diploma. The signal, in
other words, was held fixed by construction. Earnings were not:
graduates after the reform earned about 16\% less in economics and about
13\% less in business, and they did visibly worse in recruitment and
ended up at weaker firms. Something in the missing courses was doing
work that the diploma was not doing.

The reasonable reading is that both mechanisms operate. A degree is
partly a hurdle and partly an education, and the proportions vary by
institution and by labor market. What the model establishes is not that
college teaches nothing. It is that a world in which college taught
nothing could look, at least in broad outline, very much like ours.

\section{Akerlof and the Market for Lemons}\label{sec-akerlof-lemons}

Spence (1973) shows what the informed side can do about asymmetric
information. George Akerlof (1970) asks what happens when they cannot do
anything about it, and the answer is that the market can break in
unexpected ways.

For most of the twentieth century, the used car market had a striking
feature. Of course the older a car is, and the more miles it has on it,
the less value it has. And for most of the car's life, this is a
\emph{reasonably} linear decline. There was, however, a big
discontinuity. When you bought a car, and drove it off the lot, it
\emph{immediately} lost a big chunk of its value, maybe as much as 20\%.
Everyone knew this, but no one had a good theory for why it would be
so.\footnote{There were lots of bad theories, which I won't go into
  here. Suffice to say, they all made bad predictions.}

Akerlof suggested this worked because of asymmetric information. Let's
work with a really simple model. Most cars are good, but some are, in
the terminology of the mid-twentieth century, \emph{lemons}. They are
just bad cars. You can drive them around the block on a test drive, but
after a few days, weeks, or months at home, you'll realise you bought a
dud. This is private information; a buyer, even one who kicks the tires
and drives it around the block a few times, can't tell whether they are
buying a lemon.

Let's work with a \emph{really} simplified model. Assume that sellers
are motivated to sell; that's to say, restrict attention to those people
who have cars and would other things being equal like to sell them at
market price. Assume that half the cars are lemons. That's absurd, but
it keeps the numbers simple. The values of good and bad cars are in
Table~\ref{tbl-lemons}.

\begin{longtable}[]{@{}lcc@{}}
\caption{An extreme version of the market for
lemons}\label{tbl-lemons}\tabularnewline
\toprule\noalign{}
\endfirsthead
\endhead
\bottomrule\noalign{}
\endlastfoot
& \textbf{To seller} & \textbf{To buyer} \\
\textbf{Good car} & 100 & 120 \\
\textbf{Lemon} & 40 & 60 \\
\end{longtable}

Every car is worth more to a buyer than to its seller, so if everyone
could see quality, every car would change hands, the market would clear,
and everyone would be better off.

Now suppose buyers cannot tell the two apart. A buyer who thinks the
odds are even values a random car at \(\frac{120 + 60}{2} = 90\), and so
is willing to pay up to 90.

But a good car is worth 100 to its owner. At 90, no owner of a good car
sells. So the only cars on the market are lemons, and the buyer who
works this out revises: a car actually on offer is worth 60. Buyers now
offer at most 60. Lemon owners, who value their cars at 40, are happy to
sell.

The market clears, in the sense that the cars put up for sale sell, but
only lemons trade, at 60. Good cars never change hands, even though
every good-car sale would have produced a surplus of 20.

In more realistic models, where cars have continuous quality, there
aren't \emph{that many} lemons, and some holders of good cars are
\emph{really} motivated to sell, you'll get more trade. But still, the
fact that a car is up for sale will be evidence that it's bad, and the
used-car discount is what buyers rationally demand in order to overcome
this evidence.

The underlying structure here ends up being simmilar to the Spence
model. Indeed, Spence was motivated to come up with his model after his
advisors recommended reading Akerlof's paper.\footnote{Spence tells this
  story in his Nobel lecture (Spence 2002). Those advisors included
  Thomas Schelling, who we've already met, and Kenneth Arrow, who we
  will discuss \emph{at length} in the next part of the course.} The
decision to put a car up for sale is a signal; a signal that you're an
unfortunate owner of a lemon.

As well as explaining the used-car discount that was prevalent for the
twentieth century, Akerlof's model also predicts the conditions under
which the problem should ease. Those conditions have by now largely been
met, and the discount is really small now. The model requires two things
to hold. One is that there is enough variation in manufacturing quality
that the possibility of lemons is a serious problem. That's just not as
big a deal now; most car companies have excellent quality control and
the cars rolling off the line are pretty similar to one another. The
other condition is that there is an information asymmetry between buyers
and sellers. That isn't nearly as dramatic now. Vehicle history services
and cheap independent inspections, reduce the information gap, and for
what's left manufacturer certification removes the downside of buying a
lemon. Both of the model's assumptions were weakened, and the effect it
was built to explain faded. That's a really neat result; it's one thing
to explain a striking phenomenon, it's another altogether to predict
when the phenomena will fade.

One last very striking point for the philosophy students. Akerlof's
model was really unpopular with economists when it was first published.
They thought he was abandoning rational choice explanation. This is odd
because in his model, everyone is completely rational by the standards
that philosophers would use. They all have sensible priors, which they
update by conditionalisation, and then they utility maximise. What's
true is that in his model, not everyone is \emph{fully informed}.
Nowadays, I think economists have largely come around to philosophers'
way of talking; ignorance is unfortunate, but it's not irrational. But
there was something really sociologically puzzling about how economists
reacted to it, and a different course would explore this more.

\bookmarksetup{startatroot}

\chapter*{References}\label{references}
\addcontentsline{toc}{chapter}{References}

\markboth{References}{References}

\protect\phantomsection\label{refs}
\begin{CSLReferences}{1}{1}
\bibitem[\citeproctext]{ref-akerlof1970}
Akerlof, George A. 1970. {``The Market for {`Lemons'}: Quality
Uncertainty and the Market Mechanism.''} \emph{The Quarterly Journal of
Economics} 84 (3): 488--500. \url{https://doi.org/10.2307/1879431}.

\bibitem[\citeproctext]{ref-allais1953}
Allais, Maurice. 1953. {``Le Comportement de l'homme Rationnel Devant Le
Risque: Critique Des Postulats Et Axiomes de l'{É}cole Am{é}ricaine.''}
\emph{Econometrica} 21 (4): 503--46.
\url{https://doi.org/10.2307/1907921}.

\bibitem[\citeproctext]{ref-arteaga2018}
Arteaga, Carolina. 2018. {``The Effect of Human Capital on Earnings:
Evidence from a Reform at {C}olombia's Top University.''} \emph{Journal
of Public Economics} 157: 212--25.
\url{https://doi.org/10.1016/j.jpubeco.2017.10.007}.

\bibitem[\citeproctext]{ref-aumannbrandenburger1995}
Aumann, Robert, and Adam Brandenburger. 1995. {``Epistemic Conditions
for Nash Equilibrium.''} \emph{Econometrica} 63 (5): 1161--80.
\url{https://doi.org/10.2307/2171725}.

\bibitem[\citeproctext]{ref-balle2025volIII}
Balle, Solvej. (2021) 2025. \emph{On the Calculation of Volume (Book
{III})}. Translated by Sophia Hersi Smith and Jennifer Russell. New
Directions.

\bibitem[\citeproctext]{ref-benporathdekel1992}
Ben-Porath, Elchanan, and Eddie Dekel. 1992. {``Signaling Future Actions
and the Potential for Sacrifice.''} \emph{The Review of Economic
Studies} 59 (1): 127--51. \url{https://doi.org/10.2307/2297858}.

\bibitem[\citeproctext]{ref-blavatskyy2022}
Blavatskyy, Pavlo, Andreas Ortmann, and Valentyn Panchenko. 2022. {``On
the Experimental Robustness of the Allais Paradox.''} \emph{American
Economic Journal: Microeconomics} 14 (1): 143--63.
\url{https://doi.org/10.1257/mic.20190153}.

\bibitem[\citeproctext]{ref-bonanno2024}
Bonanno, Giacomo. 2024. \emph{Game Theory}. 3rd ed.
\url{https://faculty.econ.ucdavis.edu/faculty/bonanno/PDF/GT_book.pdf}.

\bibitem[\citeproctext]{ref-chokreps1987}
Cho, In-Koo, and David M. Kreps. 1987. {``Signaling Games and Stable
Equilibria.''} \emph{The Quarterly Journal of Economics} 102 (2):
179--221. \url{https://doi.org/10.2307/1885060}.

\bibitem[\citeproctext]{ref-ellsberg1961}
Ellsberg, Daniel. 1961. {``Risk, Ambiguity, and the Savage Axioms.''}
\emph{The Quarterly Journal of Economics} 75 (4): 643--69.
\url{https://doi.org/10.2307/1884324}.

\bibitem[\citeproctext]{ref-harsanyi1967}
Harsanyi, John C. 1967. {``Games with Incomplete Information Played by
{`{B}ayesian'} Players, {P}art {I}: The Basic Model.''} \emph{Management
Science} 14 (3): 159--82. \url{https://doi.org/10.1287/mnsc.14.3.159}.

\bibitem[\citeproctext]{ref-harsanyi1973}
Harsanyi, John C. 1973. {``Games with Randomly Disturbed Payoffs: A New
Rationale for Mixed-Strategy Equilibrium Points.''} \emph{International
Journal of Game Theory} 2: 1--23.
\url{https://doi.org/10.1007/BF01737554}.

\bibitem[\citeproctext]{ref-hotelling1929}
Hotelling, Harold. 1929. {``Stability in Competition.''} \emph{The
Economic Journal} 39 (153): 41--57.

\bibitem[\citeproctext]{ref-humberstone2015}
Humberstone, Lloyd. 2015. \emph{Philosophical Applications of Modal
Logic}. College {P}ublications.

\bibitem[\citeproctext]{ref-hume1739}
Hume, David. (1739) 2003. \emph{A Treatise of Human Nature}. Project
Gutenberg. \url{https://gutenberg.org/ebooks/4705}.

\bibitem[\citeproctext]{ref-jevons1871}
Jevons, W. Stanley. 1871. \emph{The Theory of Political Economy}.
Macmillan; Co.

\bibitem[\citeproctext]{ref-lewis1969}
Lewis, David. 1969. \emph{Convention: A Philosophical Study}. Harvard
University Press.

\bibitem[\citeproctext]{ref-mckelvey1992}
McKelvey, Richard D., and Thomas R. Palfrey. 1992. {``An Experimental
Study of the Centipede Game.''} \emph{Econometrica} 60 (4): 803--36.
\url{https://doi.org/10.2307/2951567}.

\bibitem[\citeproctext]{ref-nash1951}
Nash, John. 1951. {``Non-Cooperative Games.''} \emph{Annals of
Mathematics} 54 (2): 286--95. \url{https://doi.org/10.2307/1969529}.

\bibitem[\citeproctext]{ref-vonneumann1944}
Neumann, John von, and Oskar Morgenstern. 1944. \emph{Theory of Games
and Economic Behavior}. Princeton University Press.

\bibitem[\citeproctext]{ref-OECD2025}
OECD. 2025. \emph{Education at a Glance 2025: {OECD} Indicators}. OECD
Publishing. \url{https://doi.org/10.1787/1c0d9c79-en}.

\bibitem[\citeproctext]{ref-polak2007}
Polak, Ben. 2007. \emph{Nash Equilibrium: Bad Fashion and Bank Runs}.
Lecture 5, ECON 159: Game Theory, Open Yale Courses, Yale University.
\url{https://oyc.yale.edu/economics/econ-159/lecture-5}.

\bibitem[\citeproctext]{ref-potts2008}
Potts, Christopher. 2008. {``Interpretive Economy, {S}chelling Points,
and Evolutionary Stability.''}
\url{http://semanticsarchive.net/Archive/jExYWZlN/potts-interpretive-economy-mar08.pdf}.

\bibitem[\citeproctext]{ref-reepbenjamin1968}
Reep, Charles, and Bernard Benjamin. 1968. {``Skill and Chance in
Association Football.''} \emph{Journal of the Royal Statistical Society,
Series A} 131 (4): 581--85. \url{https://doi.org/10.2307/2343726}.

\bibitem[\citeproctext]{ref-reichenbach1949}
Reichenbach, Hans. 1949. \emph{The Theory of Probability: An Inquiry
into the Logical and Mathematical Foundations of the Calculus of
Probability}. 2nd ed. Translated by Ernest H. Hutten and Maria
Reichenbach. University of California Press.

\bibitem[\citeproctext]{ref-rosenthal1981}
Rosenthal, Robert W. 1981. {``Games of Perfect Information, Predatory
Pricing and the Chain-Store Paradox.''} \emph{Journal of Economic
Theory} 25 (1): 92--100.
\url{https://doi.org/10.1016/0022-0531(81)90018-1}.

\bibitem[\citeproctext]{ref-rousseau1755}
Rousseau, Jean-Jacques. (1755) 2004. \emph{A Discourse Upon the Origin
and the Foundation of the Inequality Among Mankind}. Project Gutenberg.
\url{https://gutenberg.org/ebooks/11136}.

\bibitem[\citeproctext]{ref-schaeffer2007}
Schaeffer, Jonathan, Neil Burch, Yngvi Björnsson, et al. 2007.
{``Checkers Is Solved.''} \emph{Science} 317 (5844): 1518--22.
\url{https://doi.org/10.1126/science.1144079}.

\bibitem[\citeproctext]{ref-schelling1960}
Schelling, Thomas C. 1960. \emph{The Strategy of Conflict}. Harvard
University Press.

\bibitem[\citeproctext]{ref-sen2017}
Sen, Amartya. 2017. \emph{Collective Choice and Social Welfare: An
Expanded Edition}. Harvard University Press.

\bibitem[\citeproctext]{ref-skyrms1996}
Skyrms, Brian. 1996. \emph{Evolution of the Social Contract}. Cambridge
University Press.

\bibitem[\citeproctext]{ref-spence2002}
Spence, A. Michael. 2002. {``Signaling in Retrospect and the
Informational Structure of Markets.''} \emph{American Economic Review}
92 (3): 434--59. \url{https://doi.org/10.1257/00028280260136200}.

\bibitem[\citeproctext]{ref-spence1973}
Spence, Michael. 1973. {``Job Market Signaling.''} \emph{The Quarterly
Journal of Economics} 87 (3): 355--74.
\url{https://doi.org/10.2307/1882010}.

\bibitem[\citeproctext]{ref-stalnaker1996}
Stalnaker, Robert. 1996. {``Knowledge, Belief and Counterfactual
Reasoning in Games.''} \emph{Economics and Philosophy} 12 (2): 133--63.
\url{https://doi.org/10.1017/S0266267100004132}.

\bibitem[\citeproctext]{ref-stalnaker1998}
Stalnaker, Robert. 1998. {``Belief Revision in Games: Forward and
Backward Induction.''} \emph{Mathematical Social Sciences} 36 (1):
31--56. \url{https://doi.org/10.1016/S0165-4896(98)00007-9}.

\bibitem[\citeproctext]{ref-williamson2000}
Williamson, Timothy. 2000. \emph{Knowledge and Its Limits}. {Oxford
University Press}.

\bibitem[\citeproctext]{ref-wilson2008}
Wilson, Jonathan. 2008. \emph{Inverting the Pyramid: The History of
Football Tactics}. Orion.

\end{CSLReferences}


\backmatter


\end{document}
